% L'émigration — Allemand 1ère A — Cours signé OnBuch+
% Programme officiel MINESEC. Compiler avec : tectonic ALA24-l-emigration.tex
\newcommand{\DOCMATIERE}{Allemand}
\newcommand{\DOCNIVEAU}{1ère A}
\newcommand{\DOCMODULE}{Chapitre 5 : Citoyenneté}
\newcommand{\DOCLECON}{ALA24}
\newcommand{\DOCTITRE}{L'émigration}
% OnBuch+ — préambule « cours signés » ENRICHI (compilé avec tectonic / XeLaTeX)
% Système visuel « Pop » d'OnBuch+ : fond crème (jamais blanc), bordures encre
% épaisses, ombres « dures » décalées, titres Archivo Black en capitales.
% Version MATHÉMATIQUES : graphes (pgfplots/tikz), boîtes pédagogiques
% (définition, propriété, méthode, attention, le savais-tu, exercice résolu,
% exercice, TP, bilan), page de garde et signature OnBuch+ ancrée en bas.
%
% Macros attendues dans main.tex AVANT \input{preamble} :
%   \DOCMATIERE  \DOCNIVEAU  \DOCMODULE  \DOCLECON (numéro)  \DOCTITRE
\documentclass[11pt]{article}

\usepackage{fontspec}
\usepackage[ngerman,french]{babel} % français = langue principale ; l'allemand via \all{..} / textebox
\usepackage[a4paper,margin=2cm,top=3.4cm,bottom=2.8cm,headheight=1.4cm,footskip=1.3cm]{geometry}
\usepackage{amsmath,amssymb,amsfonts}
\usepackage{mathtools} % \xmapsto, \coloneqq... souvent utilisés par les modèles
\usepackage{cancel} % simplifications barrées (bilans de forces)
% \abs{z} (module, valeur absolue) et \norm{u} (norme), utilisés par les modèles
\providecommand{\abs}{}\let\abs\relax\DeclarePairedDelimiter{\abs}{\lvert}{\rvert}
\providecommand{\norm}{}\let\norm\relax\DeclarePairedDelimiter{\norm}{\lVert}{\rVert}
\usepackage[table]{xcolor} % [table] : fournit aussi \rowcolors (alternance de lignes)
\usepackage{pagecolor}
\usepackage{tikz}
\usetikzlibrary{arrows.meta,positioning,calc,shapes.geometric,shapes.misc,shapes.arrows,decorations.pathmorphing,decorations.pathreplacing,decorations.markings,patterns,shadows,mindmap,trees,fit,backgrounds,matrix,intersections,angles,quotes}
\usepackage{pgfplots}
\usepackage[version=4]{mhchem} % équations nucléaires \ce{^{235}_{92}U}
\usepackage[european,siunitx]{circuitikz} % schémas électriques
\pgfplotsset{compat=1.18}
\usepgfplotslibrary{fillbetween,groupplots} % aires entre courbes (calcul intégral)
\usepackage{enumitem}
\usepackage{fancyhdr}
% [most] charge toutes les bibliothèques tcolorbox (listings, minted, raster,
% documentation...) et gonfle inutilement la mémoire TeX sur un document long
% (~300 boîtes) : on ne charge que ce qui est réellement utilisé.
\usepackage[skins,breakable]{tcolorbox}
\usepackage{graphicx}
\usepackage{colortbl}
\usepackage{booktabs}
\usepackage{tabularx}
\usepackage{diagbox} % case d'en-tête diagonale des tableaux à double entrée
% Colonnes extensibles centrée / gauche / droite (C, L, R), souvent utilisées par les modèles
\newcolumntype{C}{>{\centering\arraybackslash}X}
\newcolumntype{L}{>{\raggedright\arraybackslash}X}
\newcolumntype{R}{>{\raggedleft\arraybackslash}X}
\usepackage{array}
\usepackage{multicol}
\usepackage{multirow}
\usepackage{siunitx}
% parse-numbers=false : accepte aussi \SI{2,5 \times 10^{-3}}{...}
\sisetup{locale=FR,output-decimal-marker={,},group-minimum-digits=4,parse-numbers=false}
\DeclareSIUnit\ohm{\ensuremath{\symup{\Omega}}} % U+2126 absent de la police
\DeclareSIUnit\division{div} % réglages d'oscilloscope (ms/div, V/div)
\DeclareSIUnit\tr{tr} % tours (vitesse de rotation, ex. \SI{1500}{\tr\per\minute})
\DeclareSIUnit\molar{M} % concentration molaire (ex. \SI{3}{\molar}, \SI{1}{\micro\molar})
\DeclareSIUnit\an{an} % année (le modèle confond parfois avec \year, un compteur TeX) — ex. \SI{5}{\kilo\meter\per\an}
\DeclareSIUnit\FCFA{FCFA} % franc CFA (ex. \SI{500}{\FCFA\per\kilo\gram})
\DeclareSIUnit\degreeBrix{\textdegree Brix} % teneur en sucre d'un sirop/jus (réfractométrie)
\DeclareSIUnit\month{mois} % le modèle utilise parfois \month, un compteur TeX, comme unité de durée
\usepackage{qrcode}
% \degree : commande fréquemment utilisée par les modèles pour un angle ou une
% température en dehors de \SI{}{} (non fournie par les paquets chargés).
\providecommand{\degree}{^{\circ}}
% \qed (fin de démonstration), utilisé par les modèles sans amsthm
\providecommand{\qed}{\hfill$\square$}
% \barre : double barre des valeurs interdites dans un tableau de variations (tkz-tab)
\providecommand{\barre}{\ensuremath{\|}}
% environnement proof (amsthm non chargé) utilisé par les modèles
\ifdefined\proof\else\newenvironment{proof}[1][Démonstration]{\par\noindent\textit{#1.}\ }{\qed\par}\fi
\usepackage{lastpage}
\usepackage{hyperref}
\hypersetup{hidelinks}

% ============================================================
% Palette « Pop » OnBuch+ — mêmes hex que la classe P (plus_ui.dart)
% ============================================================
\definecolor{poporange}{HTML}{F59321}
\definecolor{popdark}{HTML}{E07A0C}
\definecolor{popcream}{HTML}{FFF6E8}
\definecolor{popink}{HTML}{1C1714}
\definecolor{poppurple}{HTML}{7A5AE0}
\definecolor{popgreen}{HTML}{2E9E6A}
\definecolor{poppink}{HTML}{E0457B}
\definecolor{popblue}{HTML}{2F7DE1}
\definecolor{popgold}{HTML}{C98A12}
\definecolor{popmuted}{HTML}{8A7F76}
\definecolor{popline}{HTML}{EADFD2}
\definecolor{popgray}{HTML}{8A7F76}
\colorlet{popgrey}{popgray}
% Alias tolérants (le modèle nomme parfois les couleurs différemment)
\colorlet{popwhite}{white}
\colorlet{popblack}{popink}
\colorlet{popred}{poppink}
\colorlet{popyellow}{popgold}
% Teintes claires (fonds de tableaux/encadrés) — le modèle nomme parfois
% les couleurs avec un suffixe L (ex. popblueL) sans les avoir définies.
\colorlet{poporangeL}{poporange!20}
\colorlet{popdarkL}{popdark!20}
\colorlet{poppurpleL}{poppurple!20}
\colorlet{popgreenL}{popgreen!20}
\colorlet{poppinkL}{poppink!20}
\colorlet{popblueL}{popblue!20}
\colorlet{popgoldL}{popgold!20}
\colorlet{popredL}{popred!20}
\colorlet{popgrayL}{popgray!20}
% Teintes claires pour fonds de tableaux / zones
\colorlet{popblueL}{popblue!10!white}
\colorlet{poporangeL}{poporange!14!white}
\colorlet{popgreenL}{popgreen!12!white}
\colorlet{poppurpleL}{poppurple!12!white}
\colorlet{poppinkL}{poppink!12!white}
\colorlet{popgoldL}{popgold!14!white}

\pagecolor{popcream}

% ============================================================
% Typographie
% ============================================================
\defaultfontfeatures{Ligatures=TeX}
\setmainfont{PlusJakartaSans-Medium.ttf}[Path=../../../fonts/,BoldFont=PlusJakartaSans-Bold.ttf]
% Symboles, API et emojis absents de la police principale : repli sur DejaVu Sans (caractères actifs)
\newfontface\symfont{DejaVuSans.ttf}[Path=../../../fonts/]
\makeatletter
\newcommand\symactive[1]{\catcode#1=13 \begingroup\lccode`\~=#1 \lowercase{\endgroup\def~}{{\symfont\char#1}}}
\newcommand\symblank[1]{\catcode#1=13 \begingroup\lccode`\~=#1 \lowercase{\endgroup\def~}{}}
\newcommand\symhyph[1]{\catcode#1=13 \begingroup\lccode`\~=#1 \lowercase{\endgroup\def~}{\mbox{-}}}
\newcount\symc
\newcommand\symrange[2]{\symc=#1 \@whilenum\symc<#2 \do{\expandafter\symactive\expandafter{\the\symc}\advance\symc by 1 }}
\newcommand\symblankrange[2]{\symc=#1 \@whilenum\symc<#2 \do{\expandafter\symblank\expandafter{\the\symc}\advance\symc by 1 }}
\makeatother
\symrange{"0250}{"02FF}\symactive{"2713}\symactive{"2714}\symactive{"2717}\symactive{"2718}\symactive{"2610}\symactive{"2611}\symactive{"2612}
\symhyph{"2011}\symhyph{"2E1A}\symblankrange{"1F300}{"1FAFF}\symblank{"FE0F}
% Maths en sans-serif (Fira Math) avec chiffres et lettres droites de Plus
% Jakarta, pour que les formules se fondent dans le texte.
\usepackage[math-style=ISO,bold-style=ISO]{unicode-math}
\setmathfont{FiraMath-Regular.otf}[Path=../../../fonts/]
\setmathfont{PlusJakartaSans-Medium.ttf}[Path=../../../fonts/,range={up/num,up/{Latin,latin},\mathup}]
\usepackage{icomma} % virgule décimale sans espace en mode math
% Caractères mathématiques Unicode absents des polices de texte (Plus Jakarta,
% Archivo Black) : rendus par leur équivalent mathématique, en texte comme en
% formule — sinon ils s'affichent en carré vide (ex. « Arithmétique dans ℤ »).
\usepackage{newunicodechar}
\newunicodechar{ℝ}{\ensuremath{\mathbb{R}}}
\newunicodechar{ℤ}{\ensuremath{\mathbb{Z}}}
\newunicodechar{ℂ}{\ensuremath{\mathbb{C}}}
\newunicodechar{ℕ}{\ensuremath{\mathbb{N}}}
\newunicodechar{ℚ}{\ensuremath{\mathbb{Q}}}
\newunicodechar{𝔻}{\ensuremath{\mathbb{D}}}
\newunicodechar{α}{\ensuremath{\alpha}}
\newunicodechar{β}{\ensuremath{\beta}}
\newunicodechar{γ}{\ensuremath{\gamma}}
\newunicodechar{δ}{\ensuremath{\delta}}
\newunicodechar{ε}{\ensuremath{\varepsilon}}
\newunicodechar{θ}{\ensuremath{\theta}}
\newunicodechar{λ}{\ensuremath{\lambda}}
\newunicodechar{μ}{\ensuremath{\mu}}
\newunicodechar{σ}{\ensuremath{\sigma}}
\newunicodechar{τ}{\ensuremath{\tau}}
\newunicodechar{φ}{\ensuremath{\varphi}}
\newunicodechar{ψ}{\ensuremath{\psi}}
\newunicodechar{ω}{\ensuremath{\omega}}
\newunicodechar{Σ}{\ensuremath{\Sigma}}
% majuscules produites par \MakeUppercase dans les titres (α-aminés -> Α)
\newunicodechar{Α}{\ensuremath{\alpha}}
\newunicodechar{Β}{\ensuremath{\beta}}
\newunicodechar{Γ}{\ensuremath{\Gamma}}
\newunicodechar{Δ}{\ensuremath{\Delta}}
\newunicodechar{Ε}{\ensuremath{\varepsilon}}
\newunicodechar{Θ}{\ensuremath{\Theta}}
\newunicodechar{Λ}{\ensuremath{\Lambda}}
\newunicodechar{Μ}{\ensuremath{\mu}}
\newunicodechar{Τ}{\ensuremath{\tau}}
\newunicodechar{Φ}{\ensuremath{\Phi}}
\newunicodechar{Ψ}{\ensuremath{\Psi}}
\newunicodechar{Ω}{\ensuremath{\Omega}}
\newunicodechar{↦}{\ensuremath{\mapsto}}
\newunicodechar{∈}{\ensuremath{\in}}
\newunicodechar{∉}{\ensuremath{\notin}}
\newunicodechar{∧}{\ensuremath{\wedge}}
\newunicodechar{⇔}{\ensuremath{\Leftrightarrow}}
\newunicodechar{⟺}{\ensuremath{\Longleftrightarrow}}
\newunicodechar{⇒}{\ensuremath{\Rightarrow}}
\newunicodechar{⟹}{\ensuremath{\Longrightarrow}}
\newunicodechar{∘}{\ensuremath{\circ}}
\newunicodechar{∪}{\ensuremath{\cup}}
\newunicodechar{∩}{\ensuremath{\cap}}
\newunicodechar{≡}{\ensuremath{\equiv}}
\newunicodechar{⊂}{\ensuremath{\subset}}
\newunicodechar{⊥}{\ensuremath{\perp}}
\newunicodechar{∃}{\ensuremath{\exists}}
\newunicodechar{∀}{\ensuremath{\forall}}
\newunicodechar{∛}{\ensuremath{\sqrt[3]{\,}}}
\newunicodechar{∜}{\ensuremath{\sqrt[4]{\,}}}
\newunicodechar{⃗}{\ensuremath{{}^{\to}}}
\newunicodechar{ⁿ}{\ifmmode^{n}\else\textsuperscript{\ensuremath{n}}\fi}
\newunicodechar{ˣ}{\ifmmode^{x}\else\textsuperscript{\ensuremath{x}}\fi}
\newunicodechar{ᵇ}{\ifmmode^{b}\else\textsuperscript{\ensuremath{b}}\fi}
\newunicodechar{ʳ}{\ifmmode^{r}\else\textsuperscript{\ensuremath{r}}\fi}
\newunicodechar{ᵘ}{\ifmmode^{u}\else\textsuperscript{\ensuremath{u}}\fi}
\newunicodechar{ʸ}{\ifmmode^{y}\else\textsuperscript{\ensuremath{y}}\fi}
\newunicodechar{ᵃ}{\ifmmode^{a}\else\textsuperscript{\ensuremath{a}}\fi}
\newunicodechar{ᵉ}{\ifmmode^{e}\else\textsuperscript{\ensuremath{e}}\fi}
\newunicodechar{ᵏ}{\ifmmode^{k}\else\textsuperscript{\ensuremath{k}}\fi}
\newunicodechar{ᵖ}{\ifmmode^{p}\else\textsuperscript{\ensuremath{p}}\fi}
\newunicodechar{ᴺ}{\ifmmode^{N}\else\textsuperscript{\ensuremath{N}}\fi}
\newunicodechar{ᶜ}{\ifmmode^{c}\else\textsuperscript{\ensuremath{c}}\fi}
\newunicodechar{ʰ}{\ifmmode^{h}\else\textsuperscript{\ensuremath{h}}\fi}
\newunicodechar{ᵠ}{\ifmmode^{q}\else\textsuperscript{\ensuremath{q}}\fi}
\newunicodechar{ₙ}{\ifmmode_{n}\else\textsubscript{\ensuremath{n}}\fi}
\newunicodechar{ₐ}{\ifmmode_{a}\else\textsubscript{\ensuremath{a}}\fi}
\newunicodechar{ᵢ}{\ifmmode_{i}\else\textsubscript{\ensuremath{i}}\fi}
\newunicodechar{ᵣ}{\ifmmode_{r}\else\textsubscript{\ensuremath{r}}\fi}
\newunicodechar{ₖ}{\ifmmode_{k}\else\textsubscript{\ensuremath{k}}\fi}
\newunicodechar{ᵦ}{\ifmmode_{\beta}\else\textsubscript{\ensuremath{\beta}}\fi}
\newunicodechar{ₓ}{\ifmmode_{x}\else\textsubscript{\ensuremath{x}}\fi}
\newfontfamily\popdisplay{ArchivoBlack-Regular.ttf}[Path=../../../fonts/]
\newfontfamily\poplogo{SpaceGrotesk-Bold.ttf}[Path=../../../fonts/]
\newfontfamily\popsemi{PlusJakartaSans-SemiBold.ttf}[Path=../../../fonts/]
\newfontfamily\popxbold{PlusJakartaSans-ExtraBold.ttf}[Path=../../../fonts/]
\newfontfamily\popxboldls{PlusJakartaSans-ExtraBold.ttf}[Path=../../../fonts/,LetterSpace=3.0]

\setlist[enumerate]{leftmargin=1.7em,itemsep=5pt,topsep=4pt}
\setlist[itemize]{leftmargin=1.7em,itemsep=4pt,topsep=4pt,label={\color{poporange}\textbullet}}
\setlength{\parindent}{0pt}
\setlength{\parskip}{5pt}
\renewcommand{\arraystretch}{1.25}

% pgfplots : style Pop par défaut
\pgfplotsset{
  every axis/.append style={axis line style={popink,line width=1pt},tick style={popink},
    grid style={popline,line width=0.5pt},label style={font=\small},tick label style={font=\footnotesize}},
  popcurve/.style={poporange,line width=1.8pt,no markers,smooth},
}
% Même style disponible dans un \draw TikZ simple (hors pgfplots)
\tikzset{popcurve/.style={poporange,line width=1.8pt}}
\tikzset{popgrid/.style={popline,very thin}}
\colorlet{popcurve}{poporange} % le nom du style est parfois utilisé comme couleur

% ============================================================
% Pill / squiggle
% ============================================================
\newcommand{\poppill}[3]{%
  \tikz[baseline=-0.55ex]\node[draw=popink,line width=1.2pt,rounded corners=2.6pt,%
    fill=#2,inner xsep=6.5pt,inner ysep=3pt]{%
    \popxboldls\fontsize{7}{7}\selectfont\color{#3}\MakeUppercase{#1}};%
}
\newcommand{\popsquiggle}[1]{%
  \tikz[baseline=-0.4ex]{
    \draw[#1,line width=2.6pt,line cap=round]
      plot [smooth] coordinates {(0,0.09) (0.34,0.01) (0.68,0.21) (1.02,0.01) (1.36,0.10)};
  }%
}
% Mot-clé surligné (marqueur orange sous le texte)
\newcommand{\cle}[1]{{\popxbold\color{popdark}#1}}

% ============================================================
% En-tête / pied
% ============================================================
\pagestyle{fancy}
\fancyhf{}
\renewcommand{\headrulewidth}{0pt}
\renewcommand{\footrulewidth}{0pt}
\lhead{\raisebox{-0.15ex}{\poplogo\fontsize{13}{13}\selectfont\color{popink}OnBuch\color{poporange}+}}
\rhead{\poppill{\DOCMATIERE\ \textbullet\ \DOCNIVEAU\ \textbullet\ Leçon \DOCLECON}{white}{popink}}
\lfoot{\popxbold\fontsize{7.6}{7.6}\selectfont\color{popmuted}\MakeUppercase{Cours signé OnBuch+}\ \textbullet\ {\popsemi\DOCTITRE}}
\rfoot{\tikz[baseline=-0.6ex]\node[rounded corners=8.5pt,fill=popink,minimum height=17pt,minimum width=17pt,inner xsep=5pt,inner sep=0]{\popxbold\fontsize{8}{8}\selectfont\color{popcream}\thepage\,/\,\pageref*{LastPage}};}

% ============================================================
% Titres
% ============================================================
\newcommand{\chaptitre}[2]{%
  \begin{tcolorbox}[enhanced,colback=poporange,colframe=popink,boxrule=2.6pt,arc=11pt,
      drop shadow=popink,left=15pt,right=15pt,top=12pt,bottom=11pt,before skip=3pt,after skip=18pt]
    \poppill{#2}{popink}{popcream}\par\vspace{9pt}
    {\raggedright\hyphenpenalty=10000\exhyphenpenalty=10000\popdisplay\fontsize{22}{24}\selectfont\color{popcream}\MakeUppercase{#1}\par}
  \end{tcolorbox}
}
% Section numérotée : pastille orange avec le numéro + titre Archivo
\newcounter{popsec}
\newcommand{\coursec}[1]{%
  \stepcounter{popsec}\setcounter{popsub}{0}%
  \par\vspace{16pt}\noindent
  \begin{minipage}{\linewidth}
  \tikz[baseline=-0.7ex]\node[draw=popink,line width=1.6pt,fill=poporange,rounded corners=4pt,minimum size=22pt,inner sep=0]{\popdisplay\fontsize{12}{12}\selectfont\color{popcream}\thepopsec};%
  \hspace{8pt}{\popdisplay\fontsize{15}{17}\selectfont\color{popink}\MakeUppercase{#1}}\par
  \vspace{4pt}\noindent{\color{poporange}\rule{36pt}{3.6pt}}
  \end{minipage}\par\nopagebreak\vspace{8pt}
}
\newcounter{popsub}
\newcommand{\courssub}[1]{%
  \stepcounter{popsub}%
  \par\vspace{9pt}\noindent{\popxbold\fontsize{11.5}{13}\selectfont\color{popink}{\color{poporange}\thepopsec.\thepopsub}\hspace{6pt}#1}\par\nopagebreak\vspace{4pt}
}

% ============================================================
% Boîtes pédagogiques — même recette : fond blanc, bordure épaisse colorée,
% ombre dure encre, badge plein. Titre optionnel [#1].
% ============================================================
\tcbset{popbase/.style={breakable,enhanced,colback=white,boxrule=2.3pt,arc=9pt,
  shadow={3pt}{-3pt}{0pt}{popink},left=14pt,right=14pt,top=11pt,bottom=11pt,before skip=11pt,after skip=15pt,
  fontupper=\popsemi\fontsize{10.6}{14.5}\selectfont\color{popink},
  fontlower=\popsemi\fontsize{10.6}{14.5}\selectfont\color{popink}}}
% Coche dessinée (le glyphe ✓ n'existe pas dans les polices du document)
\newcommand{\popcheck}[1]{\tikz[baseline=-0.5ex]\draw[#1,line width=1.6pt,line cap=round,line join=round](-0.13,0.0)--(-0.04,-0.09)--(0.13,0.1);}
\providecommand{\coche}{\popcheck{popgreen}}
\providecommand{\resultat}[1]{\par\smallskip\noindent\fcolorbox{popgreen}{popgreenL}{\parbox{\dimexpr\linewidth-2\fboxsep-2\fboxrule\relax}{\textbf{Résultat.} #1}}\par\smallskip}
\newcommand{\popboxhead}[3]{\poppill{#1}{#2}{white}\ifx\relax#3\relax\else\hspace{6pt}{\popxbold\fontsize{10.6}{12}\selectfont\color{popink}#3}\fi\par\vspace{7pt}}

\newtcolorbox{aretenir}[1][]{popbase,colframe=popgreen,before upper={\popboxhead{À retenir}{popgreen}{#1}}}
% Texte en allemand (document de lecture, dialogue, extrait) : cadre
% d'étude. Un titre optionnel (source, auteur) s'affiche dans l'en-tête.
\newtcolorbox{textebox}[1][]{popbase,colframe=popblue,before upper={\popboxhead{Text}{popblue}{#1}\begin{otherlanguage}{ngerman}},after upper={\end{otherlanguage}}}
% Marques ✓ / ✗ (forme correcte / fautive) souvent utilisées par les modèles
\providecommand{\cross}{\textcolor{poppink}{\ensuremath{\times}}}
\providecommand{\xmark}{\cross}\providecommand{\crossmark}{\cross}\providecommand{\cmark}{\ensuremath{\checkmark}}
% Ligne de réponse / justification à compléter (inventée par les modèles)
\providecommand{\justification}[1]{\par\noindent\textit{Justification~:} \dotfill\par}
% \textipa (package tipa non chargé) : la police ne couvre pas l'API -> simple passage
\providecommand{\textipa}[1]{#1}
% Soulignement (ulem non chargé) : \uline, \ul
\providecommand{\uline}[1]{\underline{#1}}\providecommand{\ul}[1]{\underline{#1}}
% \glossaire{\item ...} inventée par les modèles : liste de glossaire
\providecommand{\glossaire}[1]{\begin{itemize}#1\end{itemize}}
% \alert (beamer) inventée par les modèles : texte mis en évidence
\providecommand{\alert}[1]{\textbf{\textcolor{poppink}{#1}}}
% variantes ASCII d'ordinaux écrites en macro
\providecommand{\iere}{\textsuperscript{re}}\providecommand{\ieme}{\textsuperscript{e}}\providecommand{\ere}{\textsuperscript{re}}\providecommand{\eme}{\textsuperscript{e}}\providecommand{\er}{\textsuperscript{er}}\providecommand{\ie}{\textsuperscript{e}}
% Ordinaux français écrits en macro par les modèles : 1\ière, 3\ième
\newcommand{\ière}{\textsuperscript{re}}\newcommand{\ième}{\textsuperscript{e}}
% \exemple (inventée par les modèles) : étiquette « Exemple : »
\providecommand{\exemple}{\textbf{Exemple~:}\ }
% \square utilisable aussi hors mode math (cases à cocher)
\AtBeginDocument{\let\squareMath\square\renewcommand{\square}{\ensuremath{\squareMath}}}
% Mot ou phrase en allemand à l'intérieur d'un paragraphe français
\newcommand{\all}[1]{\ifmmode\text{\textit{#1}}\else\begingroup\selectlanguage{ngerman}\itshape #1\endgroup\fi}
\let\esp\all % alias toléré
\newtcolorbox{exemplebox}[1][]{popbase,colframe=poporange,before upper={\popboxhead{Exemple}{poporange}{#1}}}
\newtcolorbox{definition}[1][]{popbase,colframe=popblue,before upper={\popboxhead{Définition}{popblue}{#1}}}
\newtcolorbox{propriete}[1][]{popbase,colframe=poppurple,before upper={\popboxhead{Propriété}{poppurple}{#1}}}
\newtcolorbox{methode}[1][]{popbase,colframe=popgold,before upper={\popboxhead{Méthode}{popgold}{#1}}}
\newtcolorbox{attention}[1][]{popbase,colframe=poppink,before upper={\popboxhead{Attention}{poppink}{#1}}}
\newtcolorbox{experience}[1][]{popbase,colframe=popblue,colback=popblueL,before upper={\popboxhead{Expérience}{popblue}{#1}}}
% « Le savais-tu ? » : carte violette pleine (contexte, culture, Cameroun)
\newtcolorbox{savaistu}[1][]{popbase,colback=poppurple,colframe=popink,
  fontupper=\popsemi\fontsize{10.4}{14.2}\selectfont\color{white},
  before upper={\poppill{Le savais-tu\,?}{white}{poppurple}\ifx\relax#1\relax\else\hspace{6pt}{\popxbold\color{white}#1}\fi\par\vspace{7pt}}}
% Exercice résolu : énoncé en haut, \tcblower puis la solution sur fond crème
\newcounter{popexo}\newcounter{popexr}
\newtcolorbox{exoresolu}[1][]{popbase,colframe=poporange,colbacklower=popcream,
  segmentation style={popink,line width=1.2pt,dashed},
  before upper={\stepcounter{popexr}\popboxhead{Exercice résolu \thepopexr}{poporange}{#1}},
  before lower={\poppill{Solution}{popink}{popcream}\par\vspace{6pt}}}
% Exercice (non résolu en place) — la correction est dans la partie Corrigés
\newtcolorbox{exercice}[1][]{popbase,colframe=popink,
  before upper={\stepcounter{popexo}\popboxhead{Exercice \thepopexo}{popink}{#1}}}
% Corrigé
\newtcolorbox{corrige}[1][]{popbase,colframe=popgreen,colback=popgreenL,
  before upper={\popboxhead{Corrigé}{popgreen}{#1}}}
% Objectifs / prérequis / bilan
\newtcolorbox{objectifs}{popbase,colframe=popink,colback=poporangeL,
  before upper={\popboxhead{Objectifs de la leçon}{popink}{}}}
\newtcolorbox{prerequis}{popbase,colframe=popink,colback=popblueL,
  before upper={\popboxhead{Prérequis}{popblue}{}}}
\newtcolorbox{bilan}{popbase,colframe=popink,colback=white,boxrule=2.8pt,
  before upper={\popboxhead{Fiche bilan}{poporange}{L'essentiel à connaître pour l'examen}}}
% Figure encadrée avec légende
\newtcolorbox{popfigure}[1][]{enhanced,breakable=false,colback=white,colframe=popline,boxrule=1.4pt,arc=8pt,
  left=10pt,right=10pt,top=10pt,bottom=8pt,before skip=10pt,after skip=12pt,halign=center,
  lower separated=false,halign lower=center,
  fontlower=\popsemi\fontsize{9}{11.5}\selectfont\color{popmuted},
  #1}
\newcommand{\legende}[1]{\par\vspace{6pt}{\popsemi\fontsize{9}{11.5}\selectfont\color{popmuted}#1}}

% ============================================================
% Signature OnBuch+
% ============================================================
\newcommand{\poplogoinline}[1]{{\poplogo\fontsize{#1}{#1}\selectfont\color{popink}OnBuch\color{poporange}+}}
\newcommand{\signatureonbuch}{%
  \begin{tcolorbox}[enhanced,colback=popink,colframe=popink,boxrule=2.2pt,arc=10pt,
      drop shadow=poporange,left=14pt,right=14pt,top=10pt,bottom=10pt,before skip=0pt,after skip=0pt]
    \tikz[baseline=-0.7ex]\node[circle,fill=popgreen,draw=popcream,line width=1.3pt,minimum size=22pt,inner sep=0]{\popcheck{white}};%
    \hspace{9pt}\begin{minipage}[c]{0.62\linewidth}
      {\popxbold\fontsize{12}{13}\selectfont\color{popcream}Signé\ \textbullet\ {\poplogo OnBuch\color{poporange}+}}\par\vspace{2pt}
      {\raggedright\popsemi\fontsize{8}{10.5}\selectfont\color{popline}\MakeUppercase{Équipe pédagogique OnBuch+}\\\MakeUppercase{Conforme au programme officiel MINESEC}\par}
    \end{minipage}\hfill
    {\poplogo\fontsize{22}{22}\selectfont\color{popcream}OnBuch\color{poporange}+}
  \end{tcolorbox}%
}

% ============================================================
% Page de garde
% #1 titre · #2 matière · #3 niveau · #4 module · #5 n° leçon · #6 durée ·
% #7 description / objectifs (texte ou itemize)
% ============================================================
\newcommand{\pagegarde}[7]{%
  \thispagestyle{empty}
  \newgeometry{a4paper,margin=1.7cm,top=1.6cm,bottom=1.4cm}%
  \begin{tikzpicture}[remember picture,overlay]
    \fill[poporange!18] ([xshift=-3.2cm,yshift=-3.6cm]current page.north east) circle (4.6cm);
    \fill[poppurple!14] ([xshift=2.4cm,yshift=7.6cm]current page.south west) circle (3.8cm);
  \end{tikzpicture}%
  {\poplogo\fontsize{30}{30}\selectfont\color{popink}OnBuch\color{poporange}+}\hfill
  \raisebox{0.6ex}{\poppill{Cours signé \textbullet\ Premium}{popink}{popcream}}\par
  \vspace{1pt}\popsquiggle{poppurple}\par
  \vspace{8pt}
  \begin{tcolorbox}[enhanced,colback=poporange,colframe=popink,boxrule=2.8pt,arc=13pt,
      drop shadow=popink,left=18pt,right=18pt,top=12pt,bottom=13pt,after skip=10pt]
    \poppill{#2 \textbullet\ #3}{popink}{popcream}\hspace{5pt}\poppill{#4}{white}{popink}\par\vspace{8pt}
    {\popdisplay\fontsize{14}{14}\selectfont\color{popink}LEÇON #5}\par\vspace{4pt}
    {\raggedright\hyphenpenalty=10000\exhyphenpenalty=10000\popdisplay\fontsize{25}{27}\selectfont\color{popcream}\MakeUppercase{#1}\par}
    \vspace{8pt}
    {\popxbold\fontsize{9.5}{9.5}\selectfont\color{popink}\MakeUppercase{Durée conseillée : #6}}
  \end{tcolorbox}
  \begin{tcolorbox}[enhanced,colback=white,colframe=popink,boxrule=2.2pt,arc=10pt,
      drop shadow=popink,left=16pt,right=16pt,top=10pt,bottom=10pt,after skip=8pt]
    \poppill{Dans cette leçon}{poppurple}{white}\par\vspace{5pt}
    \popsemi\fontsize{9.3}{12.6}\selectfont\color{popink}#7
  \end{tcolorbox}
  \noindent\begin{minipage}[t]{0.487\linewidth}
    \begin{tcolorbox}[enhanced,colback=popgreen,colframe=popink,boxrule=2.2pt,arc=10pt,
        drop shadow=popink,left=11pt,right=11pt,top=10pt,bottom=10pt,height=3.1cm,valign=center]
      \tcbox[colback=white,colframe=popink,boxrule=1.3pt,arc=5pt,boxsep=0pt,nobeforeafter,
        left=3pt,right=3pt,top=3pt,bottom=3pt,valign=center]{\qrcode[height=1.9cm]{https://play.google.com/store/apps/details?id=com.luvvix.onbuch&hl=fr}}%
      \hspace{8pt}\begin{minipage}[c]{0.5\linewidth}
        {\popdisplay\fontsize{11}{12.5}\selectfont\color{white}\MakeUppercase{Télécharge OnBuch}}\par\vspace{4pt}
        {\raggedright\popsemi\fontsize{8.4}{11}\selectfont\color{white}Gratuit \textbullet\ Google Play\\Tuteur IA, annales, concours, résultats\par}
      \end{minipage}
    \end{tcolorbox}
  \end{minipage}\hfill
  \begin{minipage}[t]{0.487\linewidth}
    \begin{tcolorbox}[enhanced,colback=poppurple,colframe=popink,boxrule=2.2pt,arc=10pt,
        drop shadow=popink,left=11pt,right=11pt,top=10pt,bottom=10pt,height=3.1cm,valign=center]
      \tikz[baseline=-0.7ex]\node[circle,fill=white,draw=popink,line width=1.3pt,minimum size=1.6cm,inner sep=0]{\popdisplay\fontsize{24}{24}\selectfont\color{poppurple}+};%
      \hspace{8pt}\begin{minipage}[c]{0.6\linewidth}
        {\popdisplay\fontsize{11}{12.5}\selectfont\color{white}\MakeUppercase{Passe à OnBuch+}}\par\vspace{4pt}
        {\raggedright\popsemi\fontsize{8.4}{11}\selectfont\color{white}Tous les cours signés, Tuteur IA illimité, fiches PDF — onglet OnBuch+ de l'app\par}
      \end{minipage}
    \end{tcolorbox}
  \end{minipage}
  \vspace{8pt}
  \popcontenu
  \vfill
  \signatureonbuch
  \restoregeometry
  \newpage
}

% Rangée « Ce cours contient » (page de garde)
\newcommand{\poptile}[3]{% #1 couleur #2 gros texte #3 légende
  \begin{tcolorbox}[enhanced,colback=#1,colframe=popink,boxrule=2pt,arc=9pt,shadow={2.5pt}{-2.5pt}{0pt}{popink},
      left=6pt,right=6pt,top=9pt,bottom=9pt,halign=center,valign=center,height=1.75cm,width=\linewidth,nobeforeafter]
    {\popdisplay\fontsize{12.5}{13}\selectfont\color{white}\MakeUppercase{#2}}\par\vspace{3pt}
    {\centering\popsemi\fontsize{7.8}{9.5}\selectfont\color{white}#3\par}
  \end{tcolorbox}}
\newcommand{\popcontenu}{%
  \par\noindent{\popxboldls\fontsize{7.5}{7.5}\selectfont\color{popmuted}CE COURS SIGNÉ CONTIENT}\par\vspace{7pt}
  \noindent\begin{minipage}[t]{0.235\linewidth}\poptile{popblue}{Cours}{complet, illustré, pas à pas}\end{minipage}\hfill
  \begin{minipage}[t]{0.235\linewidth}\poptile{poporange}{Méthodes}{et exercices résolus}\end{minipage}\hfill
  \begin{minipage}[t]{0.235\linewidth}\poptile{poppink}{Exercices}{type examen + corrigés}\end{minipage}\hfill
  \begin{minipage}[t]{0.235\linewidth}\poptile{popgreen}{Fiche bilan}{l'essentiel à retenir}\end{minipage}\par}

% Sommaire
\newtcolorbox{sommaire}{enhanced,colback=white,colframe=popink,boxrule=2.2pt,arc=10pt,
  drop shadow=popink,left=18pt,right=18pt,top=13pt,bottom=13pt,before skip=6pt,after skip=20pt,
  before upper={\poppill{Sommaire}{poppurple}{white}\par\vspace{9pt}\popsemi\fontsize{10.6}{14.5}\selectfont\color{popink}}}

% Signature de fin de cours — poussée tout en bas de la dernière page
\newcommand{\findecours}{%
  \par\vfill\nopagebreak
  \begin{center}\popsquiggle{poporange}\end{center}\vspace{4pt}
  \signatureonbuch
}
\AtBeginDocument{\def\llbracket{\mathopen{[\mkern-3.5mu[}}\def\rrbracket{\mathclose{]\mkern-3.5mu]}}} % intervalles d'entiers (glyphe ⟦ absent de la police)
% Glyphes absents de Fira Math : redéfinis à partir de glyphes disponibles
\AtBeginDocument{%
  \def\setminus{\mathbin{\backslash}}\let\smallsetminus\setminus
  \def\vdots{\mathord{\vbox{\baselineskip3.2pt\lineskiplimit0pt\kern4pt\hbox{.}\hbox{.}\hbox{.}}}}%
  \def\ddots{\mathinner{\mkern1mu\raise6.5pt\hbox{.}\mkern2mu\raise3.6pt\hbox{.}\mkern2mu\raise0.7pt\hbox{.}\mkern1mu}}%
  \def\triangle{\mathord{\scalebox{0.9}{$\symup{\Delta}$}}}%
  \def\nabla{\mathord{\raisebox{\depth}{\scalebox{1}[-1]{$\symup{\Delta}$}}}}%
  \def\checkmark{\popcheck{popdark}}%
  \def\star{\mathbin{\ast}}\def\bigstar{\mathord{\ast}}%
  \def\diamond{\mathbin{\scalebox{0.6}{$\Diamond$}}}%
  \def\bigcup{\mathop{\mathchoice{\raisebox{-0.3em}{\scalebox{1.7}{$\cup$}}}{\scalebox{1.25}{$\cup$}}{\cup}{\cup}}\displaylimits}%
  \def\bigcap{\mathop{\mathchoice{\raisebox{-0.3em}{\scalebox{1.7}{$\cap$}}}{\scalebox{1.25}{$\cap$}}{\cap}{\cap}}\displaylimits}%
  \def\lesseqqgtr{\mathrel{\lessgtr}}\def\bigoplus{\mathop{\oplus}\displaylimits}%
}
\providecommand{\ssi}{si et seulement si} % abréviation fréquente
\AtBeginDocument{\providecommand{\wideparen}{\overparen}} % arc de cercle

%% FIGFIX-BEGIN v5 — correctifs globaux des figures (docs/onbuch-generation/tools/figfix)
\usepackage{adjustbox}\usepackage{etoolbox}\tcbuselibrary{raster}
% toute figure TikZ/pgfplots plus large que la ligne est réduite à la largeur disponible
\BeforeBeginEnvironment{tikzpicture}{\begin{adjustbox}{max width=\linewidth}}
\AfterEndEnvironment{tikzpicture}{\end{adjustbox}}
% légendes pgfplots lisibles par-dessus les courbes
\pgfplotsset{every axis/.append style={legend style={fill=white,fill opacity=0.92,text opacity=1,draw=black!15,inner sep=3pt}}}
% étiquettes d'arêtes (syntaxe edge["..."]) avec fond blanc
\tikzset{every edge quotes/.append style={fill=white,inner sep=1.2pt}}
%% FIGFIX-END
\begin{document}

\pagegarde{L'émigration}{Allemand}{1ère A}{Chapitre 5 : Citoyenneté}{ALA24}{2 h}{Cette leçon propose la lecture et l'exploitation d'un texte sur l'émigration d'un jeune Camerounais en Allemagne. Elle permet d'acquérir le vocabulaire thématique du départ et de l'exil, de maîtriser le parfait (Perfekt) et les subordonnées en dass et wenn, et de produire un court texte personnel.
\vspace{4pt}
\begin{multicols}{2}\small
\begin{itemize}
\item À la fin de cette leçon, je sais comprendre globalement puis en détail un texte allemand sur l'émigration.
\item Je sais employer correctement le vocabulaire du thème (auswandern, das Ausland, das Heimweh, die Hoffnung, der Pass, die Grenze) avec articles et pluriels.
\item Je sais former et utiliser le parfait (Perfekt) avec haben ou sein pour raconter des événements passés.
\item Je sais construire des subordonnées introduites par dass et wenn en plaçant correctement le verbe conjugué.
\item Je sais distinguer les emplois de wenn (condition, temps) et éviter les confusions avec als et wann.
\item Je sais répondre par écrit et oralement à des questions de compréhension sur un texte.
\end{itemize}\end{multicols}}

\begin{sommaire}
\begin{multicols}{2}
\begin{enumerate}
\item Texte support : Kevins neues Leben in Deutschland
\item Compréhension du texte
\item Lexique thématique : die Emigration
\item Grammaire 1 : das Perfekt (le parfait)
\item Grammaire 2 : les subordonnées en dass et wenn
\item Méthode du commentaire et commentaire modèle
\item Production guidée : mon histoire d'émigration
\item Activité d'intégration
\item Méthodes et exercices résolus
\item Exercices
\item Corrigés des exercices
\item Fiche bilan
\end{enumerate}
\end{multicols}
\end{sommaire}

\begin{objectifs}
\begin{itemize}
\item À la fin de cette leçon, je sais comprendre globalement puis en détail un texte allemand sur l'émigration.
\item Je sais employer correctement le vocabulaire du thème (auswandern, das Ausland, das Heimweh, die Hoffnung, der Pass, die Grenze) avec articles et pluriels.
\item Je sais former et utiliser le parfait (Perfekt) avec haben ou sein pour raconter des événements passés.
\item Je sais construire des subordonnées introduites par dass et wenn en plaçant correctement le verbe conjugué.
\item Je sais distinguer les emplois de wenn (condition, temps) et éviter les confusions avec als et wann.
\item Je sais répondre par écrit et oralement à des questions de compréhension sur un texte.
\item Je sais résumer un texte en quelques phrases au parfait.
\item Je sais rédiger un court paragraphe guidé racontant un départ ou une arrivée à l'étranger.
\end{itemize}
\end{objectifs}

\begin{prerequis}
\begin{itemize}
\item Le présent de l'indicatif des verbes réguliers et irréguliers (Seconde)
\item Les auxiliaires haben et sein au présent (Seconde)
\item La place du verbe dans la phrase déclarative (verbe en 2e position) (Seconde)
\item Les articles définis et indéfinis au nominatif et à l'accusatif (Seconde)
\item Le vocabulaire du voyage et des transports (Seconde)
\end{itemize}
\end{prerequis}

\newpage

\coursec{Texte support : Kevins neues Leben in Deutschland}

\courssub{Situation d'accroche et problématique}

Chaque année, des milliers de jeunes Camerounais rêvent de partir étudier ou travailler à l'étranger : certains choisissent l'Allemagne, attirés par ses universités gratuites et par la qualité de ses formations. Mais que signifie vraiment quitter son pays, sa famille, sa langue ? Derrière le mot \all{Emigration} se cachent des préparatifs administratifs (\all{der Pass}, \all{das Visum}), des adieux émouvants à l'aéroport, des premières difficultés dans le pays d'accueil, le mal du pays (\all{das Heimweh}), mais aussi de grands espoirs (\all{die Hoffnung}).

En Allemagne aussi, l'émigration est un grand sujet de société : des millions de personnes y ont immigré, et beaucoup d'Allemands ont eux-mêmes émigré autrefois vers l'Amérique. Aujourd'hui, nous lisons l'histoire de Kevin, un jeune Camerounais de 19 ans arrivé à Munich, et nous apprenons à raconter des événements passés avec \cle{le parfait} (\all{das Perfekt}) et à exprimer causes et conditions avec les subordonnées en \all{dass} et \all{wenn}.

\textbf{Problématique :} comment raconter en allemand le parcours d'un émigrant, de son départ de Douala jusqu'à ses espoirs pour l'avenir ?

\courssub{Avant la lecture : anticiper grâce au titre}

\begin{methode}[Bien lire un texte allemand en quatre étapes]
\begin{enumerate}
\item \textbf{Lire le titre et anticiper.} Le titre \all{Kevins neues Leben in Deutschland} (« La nouvelle vie de Kevin en Allemagne ») annonce un changement de vie. Avant même de lire, formulez des hypothèses : Qui est Kevin ? Pourquoi va-t-il en Allemagne ? Que va-t-il y faire ?
\item \textbf{Première lecture globale, sans dictionnaire.} Écoutez d'abord la lecture du professeur, puis lisez silencieusement. Cherchez seulement le sens général : de quoi parle le texte ?
\item \textbf{Deuxième lecture détaillée, avec le glossaire.} Soulignez les mots inconnus, vérifiez-les dans le glossaire, repérez les verbes conjugués.
\item \textbf{Repérer les mots-repères.} Répondez aux questions : \all{Wer?} (qui ?), \all{Wo?} (où ?), \all{Wann?} (quand ?), \all{Warum?} (pourquoi ?).
\end{enumerate}
\end{methode}

\begin{attention}[Ne traduisez pas mot à mot !]
L'erreur classique du francophone est de traduire chaque mot séparément. Cherchez d'abord le \cle{sens global}. Les \cle{mots transparents} vous aident énormément : \all{der Pass} (le passeport), \all{die Grenze} (la frontière), \all{die Familie} (la famille), \all{die Hoffnung} (l'espérance), \all{der Flughafen} (l'aéroport) ressemblent au français ou à l'anglais. Fiez-vous aussi au contexte : si Kevin pleure à l'aéroport, vous devinez le sens de \all{weinen} (pleurer) sans dictionnaire.
\end{attention}

\courssub{Le texte : Kevins neues Leben in Deutschland}

\begin{textebox}[Kevins neues Leben in Deutschland]
Kevin ist letztes Jahr von Douala nach München ausgewandert. Er ist neunzehn Jahre alt und hat in Yaoundé sein Abitur gemacht. Weil er sehr gute Noten hatte, hat er ein Stipendium für ein Studium in Deutschland bekommen. Sein Traum: Er will Ingenieur werden.

Vor der Abreise hatte Kevin viel zu tun. Er hat einen Pass beantragt und ein Visum bei der deutschen Botschaft bekommen. Seine Eltern haben ihm einen großen Koffer gekauft, und seine kleine Schwester hat ihm ein Foto der Familie geschenkt. Am Tag der Abreise ist die ganze Familie zum Flughafen von Douala gekommen. Seine Mutter hat geweint, als er „Tschüss“ gesagt hat. Kevin hat versprochen: „Ich rufe jede Woche an!“

An der Grenze hat der Beamte seinen Pass kontrolliert. Nach zwölf Stunden Flug ist Kevin in München gelandet. Am Anfang war alles schwierig: Das Wetter war kalt, das Essen war anders, und Kevin hat die Sprache nicht gut verstanden. In Deutschland hat Kevin zuerst Heimweh gehabt, weil er seine Familie und die kamerunische Küche vermisst hat. Er sagt, dass er oft an Kamerun denkt, besonders an den Fußball mit seinen Freunden und an den Markt von Douala.

Aber jetzt geht es ihm besser. Er hat neue Freunde aus vielen Ländern gefunden und hat einen kleinen Job in der Universitätsbibliothek. Wenn er Heimweh hat, kocht er Eru oder Ndolé, und dann fühlt er sich fast wie zu Hause. Kevin arbeitet viel, denn er weiß, dass das Studium nicht leicht ist. Wenn er sein Studium beendet hat, will er nach Kamerun zurückkehren und seinem Land helfen. Er hat große Hoffnung für die Zukunft: „Kamerun braucht Ingenieure“, sagt er stolz.
\end{textebox}

\textbf{Glossaire (mots nouveaux) :}

\begin{center}
\begin{tabularx}{\linewidth}{|X|X|}
\hline
\rowcolor{popblueL} \textbf{Deutsch} & \textbf{Français} \\
\hline
auswandern (ist ausgewandert) & émigrer, quitter son pays \\
\hline
das Stipendium, -ien & la bourse (d'études) \\
\hline
die Abreise, -n & le départ \\
\hline
beantragen (hat beantragt) & demander officiellement (un document) \\
\hline
die Botschaft, -en & l'ambassade \\
\hline
der Beamte, -n & le fonctionnaire, l'agent \\
\hline
das Heimweh & le mal du pays \\
\hline
vermissen (hat vermisst) & regretter l'absence de, manquer (qqn) \\
\hline
versprechen (hat versprochen) & promettre \\
\hline
zurückkehren (ist zurückgekehrt) & retourner, revenir \\
\hline
die Zukunft & l'avenir \\
\hline
stolz & fier, fière \\
\hline
\end{tabularx}
\end{center}

\courssub{Compréhension globale : qui ? où ? quand ? pourquoi ?}

\begin{exemplebox}[Les mots-repères appliqués au texte]
\begin{itemize}
\item \all{Wer?} — Kevin, 19 ans, un jeune Camerounais ; sa famille (Mutter, Eltern, Schwester).
\item \all{Wo?} — De Douala (\all{von Douala}) à Munich (\all{nach München}) ; l'aéroport de Douala, l'université de Munich.
\item \all{Wann?} — \all{letztes Jahr} (l'année dernière) ; après son Abitur (le baccalauréat).
\item \all{Warum?} — Pour étudier : \all{Er will Ingenieur werden} ; il a obtenu une bourse (\all{ein Stipendium}) grâce à ses bonnes notes.
\end{itemize}
\end{exemplebox}

\begin{exoresolu}[Questions de compréhension globale]
Énoncé : Répondez en allemand par des phrases complètes. 1) \all{Woher kommt Kevin?} 2) \all{Warum ist er nach Deutschland gegangen?} 3) \all{Was hat er an der Grenze gezeigt?}
\tcblower
Solution détaillée : 1) \all{Kevin kommt aus Kamerun, aus Douala.} (Kevin vient du Cameroun, de Douala.) 2) \all{Er ist nach Deutschland gegangen, weil er dort studieren will. Er hat ein Stipendium bekommen.} (Il est allé en Allemagne parce qu'il veut y étudier. Il a obtenu une bourse.) 3) \all{Er hat seinen Pass gezeigt; der Beamte hat ihn kontrolliert.} (Il a montré son passeport ; le fonctionnaire l'a contrôlé.) Remarquez le parfait : auxiliaire \all{haben}/\all{sein} + participe II en fin de phrase.
\end{exoresolu}

\courssub{Deux notions clés : Emigration et Immigration}

\begin{definition}[Emigration et Immigration]
\cle{Die Emigration} (l'émigration) : le fait de \textbf{quitter} son pays pour vivre à l'étranger. Le verbe : \all{auswandern}. Kevin \emph{émigre} du Cameroun : \all{Kevin ist aus Kamerun ausgewandert.}

\cle{Die Immigration} (l'immigration) : le fait d'\textbf{entrer} dans un pays étranger pour s'y installer. Le verbe : \all{einwandern}. Vu d'Allemagne, Kevin \emph{immigre} : \all{Kevin ist in Deutschland eingewandert.}

La même personne est donc \all{Auswanderer} (émigrant) pour son pays de départ et \all{Einwanderer} (immigrant) pour son pays d'arrivée.
\end{definition}

\begin{center}
\begin{tabularx}{\linewidth}{|X|X|X|}
\hline
\rowcolor{popblueL} \textbf{Français} & \textbf{Deutsch} & \textbf{Remarque} \\
\hline
émigrer (quitter) & auswandern & verbe à particule séparable : \all{Er wandert aus.} \\
\hline
immigrer (entrer) & einwandern & \all{Er wandert in Deutschland ein.} \\
\hline
l'émigrant & der Auswanderer, - & aussi : \all{der Emigrant, -en} \\
\hline
l'immigrant & der Einwanderer, - & aussi : \all{der Immigrant, -en} \\
\hline
l'étranger (le pays) & das Ausland & toujours singulier : \all{im Ausland leben} \\
\hline
\end{tabularx}
\end{center}

\courssub{Carte mentale du thème}

\begin{popfigure}
\begin{center}
\begin{tikzpicture}[mindmap, grow cyclic, every node/.style={concept, align=center, font=\small}, concept color=popblue!40, level 1/.append style={level distance=3.2cm, sibling angle=90}]
\node[concept color=popblue!60, font=\small\bfseries] {die Emigration}
  child[concept color=poporange!50] { node[align=center]{Gründe\\ Studium, Arbeit,\\ Stipendium} }
  child[concept color=popgreen!50] { node[align=center]{Vorbereitung\\ Pass, Visum,\\ Koffer, Abschied} }
  child[concept color=poppink!50] { node[align=center]{Probleme\\ Heimweh, Sprache,\\ Wetter, Grenze} }
  child[concept color=popgold!60] { node[align=center]{Hoffnungen\\ Job, Freunde,\\ zurückkehren} };
\end{tikzpicture}
\end{center}
\legende{Carte mentale : les quatre axes du thème « die Emigration » dans le texte de Kevin.}
\end{popfigure}

\courssub{Observation grammaticale : préparer les sections suivantes}

Relisons le texte avec des yeux de grammairien. Deux phénomènes sautent aux yeux ; ils feront l'objet des sections 4 et 5.

\textbf{1. Le parfait (\all{das Perfekt}).} Pour raconter les événements passés, le texte utilise massivement le parfait : auxiliaire \all{haben} ou \all{sein} conjugué en position 2 + participe II renvoyé à la fin de la phrase.

\begin{center}
\begin{tabular}{|l|l|}
\hline
\rowcolor{popgreenL} \textbf{Phrase du texte} & \textbf{Traduction} \\
\hline
\all{Kevin ist ... ausgewandert.} & Kevin a émigré. \\
\hline
\all{Er hat sein Abitur gemacht.} & Il a eu son baccalauréat. \\
\hline
\all{Seine Mutter hat geweint.} & Sa mère a pleuré. \\
\hline
\all{Der Beamte hat seinen Pass kontrolliert.} & Le fonctionnaire a contrôlé son passeport. \\
\hline
\all{Kevin hat Heimweh gehabt.} & Kevin a eu le mal du pays. \\
\hline
\all{Er hat neue Freunde gefunden.} & Il a trouvé de nouveaux amis. \\
\hline
\end{tabular}
\end{center}

\textbf{2. Les subordonnées en \all{dass} et \all{wenn}.} Pour exprimer ce que Kevin dit, pense ou sait (\all{dass} = que) et pour exprimer une condition ou un moment futur (\all{wenn} = quand, si), le texte emploie des subordonnées où le verbe conjugué se place \textbf{à la fin} :

\all{Er sagt, dass er oft an Kamerun denkt.} « Il dit qu'il pense souvent au Cameroun. »

\all{Wenn er sein Studium beendet hat, will er zurückkehren.} « Quand il aura terminé ses études, il veut rentrer. »

\begin{exercice}[À vous de jouer : repérage dans le texte]
\begin{enumerate}
\item Relevez dans le texte cinq autres verbes conjugués au parfait et soulignez l'auxiliaire et le participe II.
\item Trouvez dans le texte une subordonnée introduite par \all{weil} (parce que) et une autre par \all{als} (quand, au passé). Où se trouve le verbe conjugué ?
\item Classez les mots suivants dans la carte mentale : \all{das Visum, die Hoffnung, das Heimweh, das Stipendium, der Pass, die Grenze}.
\end{enumerate}
\end{exercice}

\begin{corrige}[Exercice de repérage]
1) Exemples : \all{hat ... bekommen} (a obtenu), \all{haben ... gekauft} (ont acheté), \all{ist ... gelandet} (a atterri), \all{hat ... versprochen} (a promis), \all{hat ... vermisst} (a manqué/regretté). 2) \all{..., weil er seine Familie vermisst hat.} et \all{..., als er „Tschüss“ gesagt hat.} : dans les deux cas, le verbe conjugué est rejeté \textbf{à la fin} de la subordonnée. 3) Vorbereitung : \all{das Visum, der Pass, die Grenze} ; Gründe : \all{das Stipendium} ; Probleme : \all{das Heimweh} ; Hoffnungen : \all{die Hoffnung}.
\end{corrige}

\begin{savaistu}[L'Allemagne, pays d'immigration... et d'émigration]
L'Allemagne compte aujourd'hui plus de 20 millions de personnes issues de l'immigration (\all{Menschen mit Migrationshintergrund}), soit environ un quart de la population. Mais l'histoire a aussi connu le mouvement inverse : entre 1850 et 1930, environ 6 millions d'Allemands ont eux-mêmes émigré vers les États-Unis, fuyant la pauvreté ou cherchant une vie meilleure. L'allemand s'est d'ailleurs mêlé à l'anglais américain : le mot \emph{hamburger} vient de la ville de Hambourg ! Le mot \all{Heimweh} lui-même est un mot allemand que plusieurs langues ont emprunté, tant il exprime bien ce sentiment universel de l'émigrant.
\end{savaistu}

\begin{aretenir}[L'essentiel de cette section]
\begin{itemize}
\item Le thème : \all{die Emigration} (quitter son pays) et \all{die Immigration} (entrer dans un pays) ; mots-clés : \all{auswandern, das Ausland, der Pass, die Grenze, das Heimweh, die Hoffnung}.
\item Méthode de lecture : anticiper avec le titre, lire globalement sans dictionnaire, puis en détail avec le glossaire, et repérer \all{Wer? Wo? Wann? Warum?}
\item Le texte raconte le passé avec \cle{le parfait} : auxiliaire \all{haben}/\all{sein} + participe II en fin de phrase.
\item Les subordonnées en \all{dass} et \all{wenn} rejettent le verbe conjugué à la fin : \all{Er sagt, dass er oft an Kamerun denkt.}
\end{itemize}
\end{aretenir}

\coursec{Compréhension du texte}

Dans la section précédente, nous avons lu et découvert le texte support \emph{Kevins neues Leben in Deutschland} : Kevin, un jeune Camerounais de Yaoundé, a quitté son pays pour étudier en Allemagne. Nous avons observé le vocabulaire nouveau et la structure du récit. Il s'agit maintenant de vérifier et d'approfondir notre \cle{compréhension} du texte : d'abord globalement (de quoi parle le texte ?), puis dans le détail (pourquoi, comment, quand ?), enfin personnellement (et toi, qu'en penses-tu ?). C'est la démarche classique de tout travail de lecture en allemand : \all{globales Verstehen} (compréhension globale), \all{detailliertes Verstehen} (compréhension détaillée), \all{persönliche Stellungnahme} (prise de position personnelle).

\courssub{La méthode de compréhension : cinq étapes}

\begin{popfigure}
\begin{tikzpicture}[
  node distance=0.55cm,
  etape/.style={rectangle, rounded corners=3pt, draw=popblue, fill=popblueL, text width=4.6cm, align=center, minimum height=0.9cm, font=\small},
  fleche/.style={-{Stealth[length=2.5mm]}, thick, poporange}
]
\node[etape, align=center] (e1){\textbf{1. Lecture globale}\\ \all{Worum geht es ?}};
\node[etape, below=of e1, align=center] (e2){\textbf{2. Lecture détaillée}\\ \all{Wer ? Warum ? Wann ? Wo ?}};
\node[etape, below=of e2, align=center] (e3){\textbf{3. Questions}\\ répondre par des phrases complètes};
\node[etape, below=of e3, align=center] (e4){\textbf{4. Résumé}\\ \all{Zusammenfassung} en 3 phrases};
\node[etape, below=of e4, align=center] (e5){\textbf{5. Production personnelle}\\ \all{Meine Meinung}};
\draw[fleche] (e1) -- (e2);
\draw[fleche] (e2) -- (e3);
\draw[fleche] (e3) -- (e4);
\draw[fleche] (e4) -- (e5);
\end{tikzpicture}
\legende{Organigramme de la méthode de compréhension d'un texte}
\end{popfigure}

\begin{methode}[Comment répondre à une question de compréhension]
\begin{enumerate}
\item \textbf{Repérer les mots-clés de la question} : le mot interrogatif (\all{warum} = pourquoi, \all{was} = que/quoi, \all{wer} = qui, \all{wann} = quand, \all{wo} = où) et le verbe.
\item \textbf{Localiser le passage du texte} qui contient la réponse : souligner la phrase ou le groupe de mots utile.
\item \textbf{Reprendre les mots de la question} dans la réponse : si la question est \all{Warum ist Kevin ausgewandert ?}, la réponse commence par \all{Kevin ist ausgewandert, weil ...}.
\item \textbf{Rédiger une phrase complète} : sujet + verbe conjugué en 2e position + compléments. Jamais un seul mot !
\item \textbf{Citer le texte} si la consigne le demande, entre guillemets allemands : \all{Im Text steht: „ ... “}.
\item \textbf{Relire} : vérifier la place du verbe, la majuscule aux noms, l'accord sujet-verbe.
\end{enumerate}
\end{methode}

\courssub{Compréhension globale : de quoi parle le texte ?}

Avant de répondre aux questions, relisons un extrait du texte support et un court document complémentaire : un journaliste allemand interviewe Kevin quelques mois après son arrivée.

\begin{textebox}[Ein Interview mit Kevin]
\textbf{Journalist:} Kevin, du bist vor sechs Monaten nach Deutschland gekommen. Warum bist du ausgewandert?

\textbf{Kevin:} Ich bin ausgewandert, weil ich in Deutschland studieren will. In Kamerun habe ich mein Abitur gemacht, und dann habe ich ein Visum beantragt.

\textbf{Journalist:} War die Reise schwierig?

\textbf{Kevin:} Ja, sehr! An der Grenze hat ein Beamter meinen Pass und mein Visum kontrolliert. Ich war sehr nervös, aber alles war in Ordnung.

\textbf{Journalist:} Und wie ist dein Leben hier?

\textbf{Kevin:} Am Anfang hatte ich Heimweh. Ich habe meine Familie und meine Freunde in Yaoundé vermisst. Das Essen und das Wetter sind hier auch anders. Aber jetzt habe ich neue Freunde gefunden, und ich habe Deutsch gelernt.

\textbf{Journalist:} Hast du Hoffnung für die Zukunft?

\textbf{Kevin:} Ja, natürlich! Ich habe große Hoffnung: Wenn ich mein Studium beendet habe, will ich Ingenieur werden. Vielleicht kehre ich dann nach Kamerun zurück, um meinem Land zu helfen.

\textbf{Journalist:} Danke für das Gespräch, Kevin!

\textbf{Kevin:} Gern geschehen!
\end{textebox}

\begin{definition}[Glossaire de l'interview]
\begin{itemize}
\item \all{beantragen} (Partizip II: \all{beantragt}) : demander officiellement (un visa) — \all{Ich habe ein Visum beantragt.} « J'ai demandé un visa. » (Verbe à préfixe inséparable \all{be-} : pas de \all{ge-} au Partizip II).
\item \all{der Beamte, die Beamten} : le fonctionnaire, l'agent — à la frontière : le douanier, le policier (nom à déclinaison \all{-n}).
\item \all{kontrollieren} : contrôler, vérifier.
\item \all{nervös} : nerveux, stressé.
\item \all{vermissen} (accusatif) : manquer (quelqu'un) — \all{Ich vermisse meine Familie.} « Ma famille me manque. »
\item \all{das Studium, die Studien} : les études universitaires.
\item \all{zurückkehren} (verbe à particule séparable, auxiliaire \all{sein}) : retourner, rentrer.
\item \all{das Gespräch, die Gespräche} : la conversation, l'entretien.
\end{itemize}
\end{definition}

\begin{exercice}[Compréhension globale]
Réponds par des phrases complètes, en allemand.
\begin{enumerate}
\item \all{Worum geht es in dem Text ?} (De quoi parle le texte ?)
\item \all{Wer ist die Hauptperson ?} (Qui est le personnage principal ?)
\item \all{Woher kommt Kevin, und wo wohnt er jetzt ?} (D'où vient Kevin, et où habite-t-il maintenant ?)
\end{enumerate}
\end{exercice}

\begin{corrige}[Compréhension globale]
\begin{enumerate}
\item \all{Der Text handelt von der Emigration eines jungen Kameruners nach Deutschland.} « Le texte parle de l'émigration d'un jeune Camerounais vers l'Allemagne. » — Formule utile : \all{Der Text handelt von ...} = « Le texte traite de ... » (génitif après \all{handeln von}).
\item \all{Die Hauptperson ist Kevin, ein junger Mann aus Yaoundé.} « Le personnage principal est Kevin, un jeune homme de Yaoundé. » (Apposition au nominatif).
\item \all{Kevin kommt aus Kamerun, und jetzt wohnt er in Deutschland.} « Kevin vient du Cameroun, et maintenant il habite en Allemagne. » (\all{aus} + datif sans article pour pays ; \all{in} + datif pour lieu).
\end{enumerate}
\end{corrige}

\courssub{Compréhension détaillée : les questions en W-}

Les questions détaillées portent sur un point précis du texte. Elles commencent souvent par un \cle{mot interrogatif} (\all{W-Frage}) : le verbe conjugué est alors en \textbf{2e position}, juste après le mot interrogatif, comme en français après « pourquoi », « quand », « que ».

\begin{exemplebox}[Observer avant de comprendre la règle]
\all{Warum ist Kevin ausgewandert ? — Er ist ausgewandert, weil er in Deutschland studieren will.}

« Pourquoi Kevin a-t-il émigré ? — Il a émigré parce qu'il veut étudier en Allemagne. »

\medskip
\all{Was hat der Beamte an der Grenze kontrolliert ? — Der Beamte hat Kevins Pass und Visum kontrolliert.}

« Qu'a contrôlé le fonctionnaire à la frontière ? — Le fonctionnaire a contrôlé le passeport et le visa de Kevin. » (Génitif \all{Kevins} pour les deux noms).

\medskip
\all{Wann hatte Kevin Heimweh ? — Am Anfang hatte er Heimweh.}

« Quand Kevin a-t-il eu le mal du pays ? — Au début, il avait le mal du pays. » (Prétérit de \all{haben}).
\end{exemplebox}

\begin{center}
\begin{tabularx}{\linewidth}{|X|X|X|}
\hline
\rowcolor{popblueL} Français & Deutsch & Remarque \\
\hline
Pourquoi ... ? & \all{Warum ... ?} & réponse avec \all{weil} (verbe à la fin) \\
\hline
Que / Quoi ... ? & \all{Was ... ?} & \all{Was hat er kontrolliert ?} \\
\hline
Qui ... ? & \all{Wer ... ?} & \all{Wer ist ausgewandert ?} \\
\hline
Quand ... ? & \all{Wann ... ?} & \all{Wann ist er gekommen ?} \\
\hline
Où ... ? & \all{Wo ... ?} / \all{Wohin ... ?} & \all{wo} = lieu (datif), \all{wohin} = direction (accusatif) \\
\hline
Comment ... ? & \all{Wie ... ?} & \all{Wie war die Reise ?} \\
\hline
\end{tabularx}
\end{center}

\begin{exoresolu}[Question détaillée sur le texte]
\textbf{Énoncé.} Réponds par une phrase complète : \all{Warum hat Kevin Heimweh gehabt ?}

\tcblower
\textbf{Solution détaillée.}
\begin{enumerate}
\item Mot interrogatif : \all{warum} → la réponse contiendra \all{weil} (parce que).
\item Passage du texte : \all{„Ich habe meine Familie und meine Freunde in Yaoundé vermisst.“}
\item Rédaction : \all{Kevin hat Heimweh gehabt, weil er seine Familie und seine Freunde in Yaoundé vermisst hat.}
\end{enumerate}
Traduction : « Kevin a eu le mal du pays parce que sa famille et ses amis de Yaoundé lui manquaient. » Attention : avec \all{weil}, le verbe conjugué (\all{hat}) va à la \textbf{fin} de la subordonnée.
\end{exoresolu}

\courssub{Les questions fermées : oui ou non, mais en phrase complète !}

\begin{propriete}[La question sans mot interrogatif : le verbe en 1re position]
Quand la question n'a \textbf{pas} de mot interrogatif (question fermée, réponse \all{ja} / \all{nein}), le \textbf{verbe conjugué occupe la 1re position} et le sujet vient juste après :

\medskip
\all{Hat Kevin Heimweh gehabt ? — Ja, er hat Heimweh gehabt.}

« Kevin a-t-il eu le mal du pays ? — Oui, il a eu le mal du pays. »

\medskip
Au parfait, c'est l'auxiliaire (\all{haben} / \all{sein}) qui passe en tête ; le participe passé reste à la fin :

\medskip
\all{Ist Kevin nach Deutschland gegangen ? — Ja, er ist nach Deutschland gegangen.}

« Kevin est-il allé en Allemagne ? — Oui, il est allé en Allemagne. »
\end{propriete}

\begin{attention}[Ne réponds jamais seulement par « ja » ou « nein » !]
À l'oral comme à l'écrit, un examinateur attend une \textbf{phrase complète}. Erreur typique du francophone :

\medskip
\all{Hat Kevin Heimweh gehabt ? — Ja.} \quad (trop court, note perdue !)

\medskip
Correct : \all{Ja, Kevin hat Heimweh gehabt.} — Le verbe conjugué est en \textbf{2e position} dans la réponse affirmative : \all{Ja,} ne compte pas comme un élément de la phrase. Autre piège : ne pas traduire « oui » par \all{doch} ; \all{doch} ne sert qu'à contredire une question négative : \all{Hat Kevin kein Heimweh gehabt ? — Doch, er hat Heimweh gehabt !} (« Si ! »).
\end{attention}

\begin{exercice}[Vrai ou faux — \all{Richtig oder falsch ?}]
Indique si la phrase est \all{richtig} (vrai) ou \all{falsch} (faux) et justifie par une citation du texte.
\begin{enumerate}
\item \all{Kevin ist vor einem Jahr nach Deutschland gekommen.}
\item \all{Der Beamte hat Kevins Pass kontrolliert.}
\item \all{Kevin hat am Anfang kein Heimweh gehabt.}
\item \all{Kevin will Ingenieur werden.}
\end{enumerate}
\end{exercice}

\begin{corrige}[Vrai ou faux]
\begin{enumerate}
\item \textbf{Falsch.} \all{Im Text steht: „Du bist vor sechs Monaten nach Deutschland gekommen.“} — Il y a six mois, pas un an.
\item \textbf{Richtig.} \all{Im Text steht: „An der Grenze hat ein Beamter meinen Pass und mein Visum kontrolliert.“} (Discours direct de Kevin : \all{meinen} = mon ; rapporté à la 3e personne : \all{Kevins}).
\item \textbf{Falsch.} \all{Im Text steht: „Am Anfang hatte ich Heimweh.“} — Kevin a eu le mal du pays au début.
\item \textbf{Richtig.} \all{Im Text steht: „Wenn ich mein Studium beendet habe, will ich Ingenieur werden.“}
\end{enumerate}
\end{corrige}

\courssub{Retrouver les équivalents dans le texte}

\begin{exoresolu}[De la version au texte]
\textbf{Énoncé.} Retrouve dans l'interview l'équivalent allemand des phrases françaises suivantes :
\begin{enumerate}
\item « J'ai demandé un visa. »
\item « J'étais très nerveux, mais tout était en règle. »
\item « Maintenant j'ai trouvé de nouveaux amis. »
\end{enumerate}

\tcblower
\textbf{Solution détaillée.}
\begin{enumerate}
\item \all{„... dann habe ich ein Visum beantragt.“} — Verbe \all{beantragen}, Partizip II \all{beantragt} (verbe à préfixe inséparable \all{be-} : pas de \all{ge-} !).
\item \all{„Ich war sehr nervös, aber alles war in Ordnung.“} — Prétérit de \all{sein} : \all{war} ; Prétérit de \all{sein} (3e pers. sg.) : \all{war}.
\item \all{„Aber jetzt habe ich neue Freunde gefunden.“} — Parfait de \all{finden} (auxiliaire \all{haben}) : \all{hat gefunden} → \all{habe ... gefunden} (1e pers. sg.).
\end{enumerate}
\end{exoresolu}

\courssub{Questions à réponse personnelle et résumé oral}

Les questions personnelles n'ont pas de réponse dans le texte : tu donnes \textbf{ton} avis, avec \all{ich}. Modèles de début de réponse : \all{Ich möchte auch ...}, \all{Ich denke, dass ...}, \all{Meiner Meinung nach ...} (« à mon avis »).

\begin{exemplebox}[Réponse personnelle modèle]
\all{Möchtest du auch im Ausland studieren ? Warum (nicht) ?}

\medskip
\all{Ja, ich möchte auch im Ausland studieren, weil ich eine neue Kultur kennenlernen will.}

« Oui, je voudrais aussi étudier à l'étranger parce que je veux découvrir une nouvelle culture. » (\all{weil} + verbe à la fin ; \all{kennenlernen} verbe à particule inséparable).

\medskip
\all{Nein, ich möchte nicht im Ausland studieren, weil ich meine Familie vermissen würde.}

« Non, je ne voudrais pas étudier à l'étranger parce que ma famille me manquerait. » (Konjunktiv II \all{würde vermissen} pour hypothèse).
\end{exemplebox}

\begin{experience}[Résumé oral en trois phrases]
Par groupes de deux, résumez le texte à l'oral en \textbf{trois phrases au parfait} (autant que possible), en utilisant au moins quatre mots du lexique : \all{auswandern}, \all{das Ausland}, \all{das Heimweh}, \all{die Hoffnung}, \all{der Pass}, \all{die Grenze}.

\medskip
Modèle : \all{Kevin ist nach Deutschland ausgewandert, weil er dort studieren will. An der Grenze hat ein Beamter Kevins Pass kontrolliert. Am Anfang hat er Heimweh gehabt, aber jetzt hat er neue Freunde gefunden und große Hoffnung.}
\end{experience}

\begin{aretenir}[L'essentiel de la section]
\begin{itemize}
\item Compréhension en trois temps : globale → détaillée → personnelle.
\item Question avec mot interrogatif : verbe en 2e position ; question fermée : verbe en 1re position.
\item Toujours répondre par une \textbf{phrase complète}, en reprenant les mots de la question.
\item Pour justifier : citer le texte entre guillemets allemands „ ... “ avec \all{Im Text steht: ...}.
\item Réponse avec \all{weil} : le verbe conjugué va à la fin de la subordonnée.
\end{itemize}
\end{aretenir}

\coursec{Lexique thématique : die Emigration}

Après avoir lu et compris l'histoire de Kevin, qui a quitté le Cameroun pour commencer une nouvelle vie en Allemagne, nous avons rencontré beaucoup de mots nouveaux : \all{auswandern}, \all{das Ausland}, \all{der Pass}, \all{das Heimweh}... Cette section a pour but de les organiser, de les expliquer et de les mémoriser durablement. Un bon vocabulaire ne s'apprend pas mot par mot, mais par \cle{champs lexicaux} : on regroupe les mots qui appartiennent à une même situation. Ici, la situation est claire : \textbf{quitter son pays, franchir une frontière, vivre à l'étranger}.

\courssub{Observation : le vocabulaire dans le texte}

\begin{textebox}[Mini-texte d'observation : An der Grenze]
Kevin steht am Flughafen von Douala. In der Hand hält er seinen Pass und sein Visum. Er verabschiedet sich von seiner Familie. Seine Mutter weint, aber Kevin hat Hoffnung: Er will in Deutschland studieren. „Ich komme zurück!“, verspricht er. Dann geht er an Bord des Flugzeugs. Nach zehn Stunden landet er in Frankfurt. An der Grenze kontrolliert ein Beamter seine Dokumente. „Willkommen in Deutschland!“, sagt der Beamte freundlich. Kevin ist jetzt im Ausland. Am Anfang ist alles schwierig: die Fremdsprache, das Essen, das Wetter. Manchmal hat er Heimweh und denkt an Yaoundé, an seine Freunde und an den Fußballplatz im Quartier. Aber er lernt jeden Tag Deutsch und findet neue Freunde. Er will sich integrieren und später vielleicht die deutsche Staatsangehörigkeit bekommen. Eines Tages, das weiß er, wird er in seine Heimat zurückkehren.
\end{textebox}

\textbf{Traduction :} Kevin est à l'aéroport de Douala. À la main, il tient son passeport et son visa. Il dit au revoir à sa famille. Sa mère pleure, mais Kevin a de l'espoir : il veut étudier en Allemagne. « Je reviendrai ! », promet-il. Puis il monte à bord de l'avion. Après dix heures, il atterrit à Francfort. À la frontière, un fonctionnaire contrôle ses documents. « Bienvenue en Allemagne ! », dit le fonctionnaire avec gentillesse. Kevin est maintenant à l'étranger. Au début, tout est difficile : la langue étrangère, la nourriture, le temps. Parfois il a le mal du pays et pense à Yaoundé, à ses amis et au terrain de football du quartier. Mais il apprend l'allemand chaque jour et trouve de nouveaux amis. Il veut s'intégrer et peut-être obtenir plus tard la nationalité allemande. Un jour, il le sait, il retournera dans son pays.

\textbf{Glossaire :} \all{der Flughafen} (die Flughäfen) : l'aéroport ; \all{sich verabschieden von} : dire au revoir à ; \all{versprechen} (hat versprochen) : promettre ; \all{an Bord gehen} : monter à bord ; \all{landen} (ist gelandet) : atterrir ; \all{das Wetter} : le temps (météo) ; \all{das Quartier} (die Quartiere) : le quartier.

\textbf{Questions d'observation :}
\begin{enumerate}
\item Quels mots du texte désignent des \textbf{documents} ?
\item Quels mots expriment des \textbf{sentiments} ?
\item Quels verbes expriment le \textbf{départ} et le \textbf{retour} ?
\end{enumerate}

Réponses attendues : 1. \all{der Pass}, \all{das Visum}, \all{die Dokumente} ; 2. \all{die Hoffnung}, \all{das Heimweh}, \all{weinen} ; 3. \all{sich verabschieden}, \all{an Bord gehen}, \all{zurückkommen}, \all{zurückkehren}.

\courssub{Champ lexical 1 : le départ et le voyage}

Voici les verbes et les noms essentiels pour raconter un départ et un retour. Observe bien la colonne « Remarque » : elle indique les pièges de conjugaison.

\begin{center}
\begin{tabularx}{\linewidth}{|X|X|X|}
\hline
\rowcolor{popblueL} Deutsch & Français & Remarque \\
\hline
auswandern (ist ausgewandert) & émigrer (quitter son pays) & verbe de mouvement : auxiliaire \all{sein} au Perfekt \\
\hline
einwandern (ist eingewandert) & immigrer (entrer dans un pays) & auxiliaire \all{sein} \\
\hline
abreisen (ist abgereist) & partir (en voyage) & particule séparable : \all{er reist ab} \\
\hline
zurückkehren (ist zurückgekehrt) & retourner, revenir & particule séparable : \all{sie kehrt zurück} \\
\hline
die Ausreise (-n) & le départ (du pays) & nom féminin \\
\hline
die Einreise (-n) & l'entrée (dans un pays) & nom féminin \\
\hline
\end{tabularx}
\end{center}

\begin{exemplebox}[Exemples traduits]
\all{Kevin ist nach Deutschland ausgewandert.} « Kevin a émigré en Allemagne. »\par
\all{Viele Menschen wandern jedes Jahr in Europa ein.} « Beaucoup de personnes immigrent chaque année en Europe. »\par
\all{Mein Bruder ist gestern abgereist.} « Mon frère est parti hier. »\par
\all{Nach fünf Jahren ist sie in ihre Heimat zurückgekehrt.} « Après cinq ans, elle est retournée dans son pays. »
\end{exemplebox}

\begin{attention}[Mouvement = auxiliaire \all{sein}]
Les verbes \all{auswandern}, \all{einwandern}, \all{abreisen}, \all{zurückkehren} expriment un \textbf{changement de lieu} : au Perfekt, ils se conjuguent avec \all{sein}, pas avec \all{haben}. On dit \all{er \textbf{ist} ausgewandert} et jamais « er hat ausgewandert ». C'est la même logique qu'en français avec « il \textbf{est} parti ».
\end{attention}

\courssub{Champ lexical 2 : les documents et l'administration}

Pour franchir une frontière, il faut des documents. Ce vocabulaire est indispensable pour comprendre les dialogues à l'aéroport ou à l'ambassade.

\begin{center}
\begin{tabularx}{\linewidth}{|X|X|X|}
\hline
\rowcolor{popgreenL} Deutsch & Français & Pluriel \\
\hline
der Pass & le passeport & die Pässe \\
\hline
das Visum & le visa & die Visa (ou : die Visen) \\
\hline
die Grenze & la frontière & die Grenzen \\
\hline
der Beamte & le fonctionnaire, l'agent & die Beamten \\
\hline
die Botschaft & l'ambassade & die Botschaften \\
\hline
kontrollieren (hat kontrolliert) & contrôler & verbe régulier \\
\hline
\end{tabularx}
\end{center}

\begin{propriete}[Majuscule, article et pluriel : la règle d'or]
En allemand, \textbf{tous les noms prennent une majuscule}, même au milieu de la phrase : \all{der Pass}, \all{die Grenze}, \all{das Ausland}. C'est une différence majeure avec le français. De plus, on n'apprend jamais un nom seul : on l'apprend toujours avec \textbf{son article} (der / die / das) et \textbf{son pluriel}. Dis-toi : « \all{der Pass, die Pässe} », « \all{die Grenze, die Grenzen} ». Le genre ne se devine pas : il se mémorise !
\end{propriete}

\begin{attention}[Le nom décliné : der Beamte]
\all{Der Beamte} appartient à la déclinaison faible (n-Deklination) : à l'accusatif et au datif singulier, il prend un \textbf{-n} : \all{Ich frage den Beamte\textbf{n}.} « Je demande au fonctionnaire. » Ne sois pas surpris de voir cette forme dans les textes.
\end{attention}

\begin{exemplebox}[Dialogue à la frontière]
\all{— Ihren Pass und Ihr Visum, bitte!} « Votre passeport et votre visa, s'il vous plaît ! »\par
\all{— Hier, bitte.} « Voilà. »\par
\all{— Der Beamte kontrolliert die Dokumente.} « Le fonctionnaire contrôle les documents. »\par
\all{— Alles in Ordnung. Willkommen!} « Tout est en règle. Bienvenue ! »
\end{exemplebox}

\courssub{Champ lexical 3 : les sentiments}

Émigrer, c'est aussi vivre des émotions fortes : la tristesse du départ, l'espoir d'une vie meilleure.

\begin{definition}[das Heimweh]
\all{Das Heimweh} (n., \textbf{sans pluriel}) signifie « le mal du pays ». Il se construit avec le verbe \all{haben} : \all{Kevin hat Heimweh.} « Kevin a le mal du pays. » On peut préciser : \all{Er hat Heimweh nach Kamerun.} « Il a le mal du Cameroun. »
\end{definition}

\begin{center}
\begin{tabularx}{\linewidth}{|X|X|X|}
\hline
\rowcolor{poppinkL} Deutsch & Français & Construction \\
\hline
das Heimweh (Sg.) & le mal du pays & Heimweh haben \\
\hline
die Hoffnung (-en) & l'espoir, l'espérance & Hoffnung haben auf + Akk. \\
\hline
die Sehnsucht (-üe) & le désir profond, la nostalgie & Sehnsucht haben nach + Dat. \\
\hline
vermissen (hat vermisst) & manquer (quelqu'un) & jemanden vermissen \\
\hline
weinen (hat geweint) & pleurer & verbe régulier \\
\hline
sich freuen auf + Akk. & se réjouir de (à venir) & ich freue mich auf... \\
\hline
\end{tabularx}
\end{center}

\begin{exemplebox}[Exemples traduits]
\all{Sie ist ins Ausland gegangen.} « Elle est partie à l'étranger. »\par
\all{An der Grenze zeigt man seinen Pass.} « À la frontière, on montre son passeport. »\par
\all{Er hat Hoffnung auf ein Stipendium.} « Il a l'espoir d'obtenir une bourse. »\par
\all{Ich vermisse meine Familie sehr.} « Ma famille me manque beaucoup. »\par
\all{Wir freuen uns auf die Ferien in der Heimat.} « Nous nous réjouissons des vacances au pays. »
\end{exemplebox}

\begin{attention}[« Manquer » : le piège de la construction]
En français, « tu me manques » place la personne qui ressent le manque au \textbf{complément} (« me »). En allemand, \all{vermissen} fonctionne de façon directe : le \textbf{sujet} est celui qui ressent le manque, et la personne qui manque est à l'accusatif : \all{Ich vermisse dich.} « Tu me manques. » Ne dis donc jamais « du vermisst mich » pour « tu me manques » : cela signifierait « je te manque » !
\end{attention}

\courssub{Champ lexical 4 : la vie à l'étranger}

\begin{center}
\begin{tabularx}{\linewidth}{|X|X|X|}
\hline
\rowcolor{popgoldL} Deutsch & Français & Pluriel / Remarque \\
\hline
das Ausland & l'étranger & \textbf{toujours au singulier} \\
\hline
die Fremdsprache & la langue étrangère & die Fremdsprachen \\
\hline
die Kultur & la culture & die Kulturen \\
\hline
sich integrieren & s'intégrer & verbe réfléchi : \all{ich integriere mich} \\
\hline
die Staatsangehörigkeit & la nationalité & die Staatsangehörigkeiten \\
\hline
\end{tabularx}
\end{center}

\begin{attention}[das Ausland : singulier et cas]
\all{Das Ausland} est \textbf{toujours au singulier} en allemand, même quand le français dit « les pays étrangers ». Attention aussi à la différence de cas : on dit \all{ins Ausland} (accusatif = mouvement : \all{Sie geht ins Ausland.} « Elle part à l'étranger. ») mais \all{im Ausland} (datif = lieu : \all{Er lebt im Ausland.} « Il vit à l'étranger. »). Question « wohin ? » (où, vers quel lieu ?) → accusatif ; question « wo ? » (où, dans quel lieu ?) → datif.
\end{attention}

\courssub{Expressions utiles à mémoriser}

\begin{aretenir}[Les expressions clés de la leçon]
\begin{itemize}
\item \all{ins Ausland gehen} : partir à l'étranger
\item \all{im Ausland leben} : vivre à l'étranger
\item \all{an der Grenze} : à la frontière
\item \all{Heimweh haben} : avoir le mal du pays
\item \all{Hoffnung haben auf + Akk.} : avoir l'espoir de
\item \all{in die Heimat zurückkehren} : retourner au pays
\item \all{den Pass kontrollieren} : contrôler le passeport
\end{itemize}
\end{aretenir}

\courssub{Familles de mots : un mot en cache trois}

En allemand, les mots d'une même famille partagent une racine. Apprendre la famille, c'est multiplier ton vocabulaire sans effort.

\begin{center}
\begin{tabularx}{\linewidth}{|X|X|X|}
\hline
\rowcolor{popblueL} Verbe & Nom (personne / action) & Adjectif ou nom lié \\
\hline
auswandern (émigrer) & der Auswanderer (l'émigrant) / die Auswanderung (l'émigration) & --- \\
\hline
einwandern (immigrer) & der Einwanderer (l'immigrant) / die Einwanderung (l'immigration) & --- \\
\hline
hoffen (espérer) & die Hoffnung (l'espoir) & hoffnungsvoll (plein d'espoir) \\
\hline
heim (à la maison) & die Heimat (la patrie, le pays natal) & das Heimweh (le mal du pays) \\
\hline
reisen (voyager) & die Reise (le voyage) / der Reisende (le voyageur) & --- \\
\hline
\end{tabularx}
\end{center}

Observe : \all{aus-} (vers l'extérieur) et \all{ein-} (vers l'intérieur) sont des préfixes de sens opposé. Le même mécanisme donne \all{die Ausreise} / \all{die Einreise}, \all{der Ausgang} / \all{der Eingang} (la sortie / l'entrée).

\begin{savaistu}[Le « Heimweh », un mot intraduisible ?]
Le mot \all{Heimweh} est un mot typiquement allemand, souvent cité parmi les mots les plus difficiles à traduire dans le monde. Il désigne une douleur physique et morale causée par l'éloignement du pays natal. Autrefois, les Suisses parlaient même du \all{Schweizerheimweh} (« le mal du pays suisse ») : leurs soldats, engagés comme mercenaires à l'étranger, tombaient littéralement malades de nostalgie loin de leurs montagnes. Aujourd'hui, le mot est passé dans d'autres langues, comme le mot \all{Wanderlust} (l'envie de voyager), son contraire !
\end{savaistu}

\courssub{Carte mentale du vocabulaire}

\begin{popfigure}
\begin{tikzpicture}[
  mindmap,
  grow cyclic,
  every node/.style={concept, align=center, font=\small},
  root concept/.append style={concept color=popblue, font=\bfseries},
  level 1 concept/.append style={level distance=3.6cm, sibling angle=90}
]
\node[root concept]{die Emigration}
  child[concept color=poporange]{ node[align=center]{der Abschied\\ \scriptsize auswandern\\ \scriptsize abreisen\\ \scriptsize zurückkehren} }
  child[concept color=popgreen]{ node[align=center]{die Dokumente\\ \scriptsize der Pass\\ \scriptsize das Visum\\ \scriptsize die Grenze} }
  child[concept color=poppink]{ node[align=center]{die Gefühle\\ \scriptsize das Heimweh\\ \scriptsize die Hoffnung\\ \scriptsize vermissen} }
  child[concept color=popgold]{ node[align=center]{das neue Leben\\ \scriptsize das Ausland\\ \scriptsize die Fremdsprache\\ \scriptsize sich integrieren} };
\end{tikzpicture}
\legende{Carte mentale : le vocabulaire de l'émigration en quatre champs lexicaux.}
\end{popfigure}

\courssub{Exercice résolu et entraînement}

\begin{exoresolu}[Complète avec le mot qui convient]
Énoncé — Complète les phrases avec : \all{Heimweh}, \all{Pass}, \all{Ausland}, \all{Hoffnung}, \all{Grenze}, \all{zurückgekehrt}.
\begin{enumerate}
\item Kevin lebt jetzt im \dots\ und lernt Deutsch.
\item An der \dots\ kontrolliert der Beamte die Dokumente.
\item Zeigen Sie bitte Ihren \dots !
\item Maria hat \dots\ : Sie denkt oft an ihre Familie in Douala.
\item Er hat \dots\ auf ein Stipendium in Deutschland.
\item Nach zehn Jahren ist sie in ihre Heimat \dots .
\end{enumerate}
\tcblower
Solution détaillée :
\begin{enumerate}
\item \all{Kevin lebt jetzt im \textbf{Ausland} und lernt Deutsch.} (lieu → datif : \all{im Ausland})
\item \all{An der \textbf{Grenze} kontrolliert der Beamte die Dokumente.} (le contrôle se fait à la frontière)
\item \all{Zeigen Sie bitte Ihren \textbf{Pass}!} (accusatif masculin : \all{Ihren Pass})
\item \all{Maria hat \textbf{Heimweh}: Sie denkt oft an ihre Familie in Douala.} (construction : \all{Heimweh haben})
\item \all{Er hat \textbf{Hoffnung} auf ein Stipendium in Deutschland.} (construction : \all{Hoffnung haben auf} + Akk.)
\item \all{Nach zehn Jahren ist sie in ihre Heimat \textbf{zurückgekehrt}.} (Perfekt avec \all{sein}, verbe de mouvement)
\end{enumerate}
\end{exoresolu}

\begin{exercice}[À toi de jouer !]
\begin{enumerate}
\item Traduis en allemand : « Mon frère a émigré en Allemagne. Il vit à l'étranger depuis deux ans. Parfois il a le mal du pays, mais il a l'espoir de trouver un bon travail. »
\item Trouve le nom qui correspond : \all{einwandern} → ? ; \all{hoffen} → ? ; \all{auswandern} (la personne) → ?
\item Vrai ou faux ? Corrige si c'est faux : a) \all{Sie hat Heimweh.} b) \all{Er ist ins Ausland gegangen.} c) \all{Kevin hat nach Deutschland ausgewandert.}
\end{enumerate}
\end{exercice}

\begin{corrige}[Exercice « À toi de jouer ! »]
\begin{enumerate}
\item \all{Mein Bruder ist nach Deutschland ausgewandert. Er lebt seit zwei Jahren im Ausland. Manchmal hat er Heimweh, aber er hat Hoffnung auf eine gute Arbeit.}
\item \all{einwandern} → \all{die Einwanderung} (ou \all{der Einwanderer}) ; \all{hoffen} → \all{die Hoffnung} ; \all{auswandern} → \all{der Auswanderer}.
\item a) Vrai. b) Vrai. c) Faux : \all{auswandern} se conjugue avec \all{sein} au Perfekt → \all{Kevin \textbf{ist} nach Deutschland ausgewandert.}
\end{enumerate}
\end{corrige}

\begin{aretenir}[L'essentiel du lexique]
\begin{itemize}
\item Chaque nom s'apprend avec son \textbf{article} et son \textbf{pluriel} ; tous les noms prennent une \textbf{majuscule}.
\item Verbes de mouvement (\all{auswandern}, \all{abreisen}, \all{zurückkehren}) → Perfekt avec \all{sein}.
\item \all{ins Ausland} (mouvement, accusatif) ≠ \all{im Ausland} (lieu, datif) ; \all{das Ausland} reste au singulier.
\item Constructions fixes : \all{Heimweh haben}, \all{Hoffnung haben auf} + Akk., \all{sich freuen auf} + Akk.
\item Les familles de mots (\all{auswandern → der Auswanderer → die Auswanderung}) permettent d'élargir vite son vocabulaire.
\end{itemize}
\end{aretenir}

\coursec{Grammaire 1 : das Perfekt (le parfait)}

Dans la section précédente, nous avons enrichi notre vocabulaire sur le thème de l'émigration : \all{auswandern}, \all{das Heimweh}, \all{die Grenze}... Mais en relisant le texte \emph{Kevins neues Leben in Deutschland}, vous avez remarqué que Kevin raconte son parcours \textbf{au passé} : il est parti, il a obtenu une bourse, il a pleuré à l'aéroport. Quel temps l'allemand utilise-t-il pour raconter ces événements ? C'est précisément le sujet de cette section : le \cle{Perfekt}, le temps du passé de la conversation.

\courssub{Observation : que remarque-t-on dans le texte ?}

Reprenons deux phrases du texte support et observons-les attentivement.

\begin{exemplebox}[Phrases tirées du texte]
\all{Kevin ist ausgewandert.} « Kevin a émigré. »\par
\all{Er hat ein Stipendium bekommen.} « Il a obtenu une bourse. »
\end{exemplebox}

Observons ensemble :

\begin{itemize}
\item Dans chaque phrase, il y a \textbf{deux verbes} : un petit verbe conjugué (\all{ist}, \all{hat}) en \textbf{deuxième position}, et une forme verbale (\all{ausgewandert}, \all{bekommen}) rejetée \textbf{à la toute fin de la phrase}.
\item Le petit verbe conjugué est \all{sein} (être) dans la première phrase, \all{haben} (avoir) dans la seconde.
\item La forme finale ressemble au participe passé français : \all{ausgewandert} = « émigré », \all{bekommen} = « obtenu ».
\end{itemize}

C'est exactement la structure du passé composé français (« il \textbf{a} \textbf{obtenu} », « il \textbf{est} \textbf{parti} »), avec une différence capitale : en allemand, le participe passé se place \textbf{en toute fin de phrase}, quelle que soit la longueur de la phrase.

\begin{definition}[Das Perfekt]
Le \cle{Perfekt} (le parfait) est le temps composé du passé le plus utilisé en allemand à l'oral et dans les récits personnels (lettres, conversations, journaux intimes). Il correspond au \textbf{passé composé français} : \all{Ich habe gegessen} = « j'ai mangé ».
\end{definition}

\courssub{Formation du Perfekt : la structure en cadre}

\begin{propriete}[Formation du Perfekt]
Le Perfekt se forme avec :\par
\textbf{l'auxiliaire \all{haben} ou \all{sein}, conjugué au présent, en 2\textsuperscript{e} position} \textbf{+} \textbf{le participe passé (Partizip II) placé en toute fin de phrase}.\par
\all{Er hat sein Abitur gemacht.} « Il a passé son baccalauréat. »\par
\all{Sie ist nach Deutschland geflogen.} « Elle est allée en Allemagne en avion. »
\end{propriete}

L'auxiliaire et le participe forment un \textbf{cadre} (en allemand \all{die Satzklammer}) qui entoure tous les autres éléments de la phrase. C'est la grande loi de la phrase allemande : le verbe conjugué ouvre le cadre en deuxième position, le participe le ferme à la fin.

\begin{popfigure}
\begin{tikzpicture}[node distance=0.4cm]
\node[draw=popblue, fill=popblueL, rounded corners, minimum width=2.2cm, minimum height=0.9cm, align=center] (suj) {\textbf{Kevin}\\ Sujet};
\node[draw=poporange, fill=poporangeL, rounded corners, minimum width=2.4cm, minimum height=0.9cm, align=center, right=of suj] (aux) {\textbf{hat}\\ auxiliaire\\ (2\textsuperscript{e} position)};
\node[draw=popmuted, fill=popcream, rounded corners, minimum width=3.4cm, minimum height=0.9cm, align=center, right=of aux] (mil) {ein Stipendium\\ in Deutschland\\ (le milieu de la phrase)};
\node[draw=popgreen, fill=popgreenL, rounded corners, minimum width=2.6cm, minimum height=0.9cm, align=center, right=of mil] (part) {\textbf{bekommen}\\ participe passé\\ (fin de phrase)};
\draw[-{Stealth}, very thick, poporange] (aux.north) to[bend left=35] node[above, align=center, text=popdark]{\small die Satzklammer : le cadre} (part.north);
\end{tikzpicture}
\legende{La structure en cadre (Satzklammer) : l'auxiliaire en 2\textsuperscript{e} position, le participe passé à la fin.}
\end{popfigure}

\begin{attention}[Le participe toujours à la fin !]
Le francophone a tendance à coller le participe à l'auxiliaire comme en français : *\all{Kevin hat bekommen ein Stipendium}. C'est une erreur grave. Tout ce qui complète le verbe se place \textbf{entre} l'auxiliaire et le participe : \all{Kevin hat \textbf{letztes Jahr in Deutschland} ein Stipendium bekommen.} « Kevin a obtenu une bourse en Allemagne l'année dernière. »
\end{attention}

Pour conjuguer correctement, il faut d'abord bien connaître le présent des deux auxiliaires :

\begin{center}
\begin{tabular}{|l|l|l|}\hline
\rowcolor{popblueL} \textbf{Personne} & \textbf{haben} & \textbf{sein} \\ \hline
ich & habe & bin \\ \hline
du & hast & bist \\ \hline
er/sie/es & hat & ist \\ \hline
wir & haben & sind \\ \hline
ihr & habt & seid \\ \hline
sie/Sie & haben & sind \\ \hline
\end{tabular}
\end{center}

\courssub{Le participe passé des verbes faibles : ge + radical + t}

Les \textbf{verbes faibles} (\all{schwache Verben}) sont les verbes réguliers. Leur participe se construit mécaniquement : \textbf{ge- + radical + -t} (ou \textbf{-et} si le radical se termine par -t, -d ou certaines consonnes difficiles à prononcer).

\begin{center}
\begin{tabularx}{\linewidth}{|X|X|X|}\hline
\rowcolor{popblueL} \textbf{Infinitif} & \textbf{Participe passé} & \textbf{Français} \\ \hline
machen & ge\textbf{mach}t & fait \\ \hline
sagen & ge\textbf{sag}t & dit \\ \hline
lernen & ge\textbf{lern}t & appris \\ \hline
wohnen & ge\textbf{wohn}t & habité \\ \hline
weinen & ge\textbf{wein}t & pleuré \\ \hline
arbeiten & ge\textbf{arbeit}et & travaillé \\ \hline
öffnen & ge\textbf{öffn}et & ouvert \\ \hline
\end{tabularx}
\end{center}

\all{Kevin hat viel geweint.} « Kevin a beaucoup pleuré. »\par
\all{Er hat Deutsch gelernt.} « Il a appris l'allemand. »

\textbf{Exception importante} : les verbes en \textbf{-ieren} (souvent d'origine française ou latine) ne prennent \textbf{pas de ge-} : \all{studieren} $\rightarrow$ \all{studiert}, \all{kontrollieren} $\rightarrow$ \all{kontrolliert}, \all{telefonieren} $\rightarrow$ \all{telefoniert}.

\all{Der Beamte hat seinen Pass kontrolliert.} « L'agent a contrôlé son passeport. »

\courssub{Le participe passé des verbes forts : ge + (radical modifié) + en}

Les \textbf{verbes forts} (\all{starke Verben}) sont les verbes irréguliers : leur radical change souvent de voyelle, et leur participe se termine en \textbf{-en}. Il faut les apprendre par cœur, comme les verbes irréguliers anglais.

\begin{center}
\begin{tabularx}{\linewidth}{|X|X|X|}\hline
\rowcolor{popblueL} \textbf{Infinitif} & \textbf{Participe passé} & \textbf{Français} \\ \hline
gehen & \textbf{gegangen} & allé \\ \hline
kommen & \textbf{gekommen} & venu \\ \hline
fliegen & \textbf{geflogen} & volé, pris l'avion \\ \hline
fahren & \textbf{gefahren} & allé (en véhicule) \\ \hline
essen & \textbf{gegessen} & mangé \\ \hline
trinken & \textbf{getrunken} & bu \\ \hline
schreiben & \textbf{geschrieben} & écrit \\ \hline
sehen & \textbf{gesehen} & vu \\ \hline
bleiben & \textbf{geblieben} & resté \\ \hline
sein & \textbf{gewesen} & été \\ \hline
\end{tabularx}
\end{center}

\all{Kevin ist nach München geflogen.} « Kevin est allé à Munich en avion. »\par
\all{Er hat einen Brief geschrieben.} « Il a écrit une lettre. »

\courssub{Verbes séparables et inséparables : où va le ge- ?}

\begin{propriete}[Place du ge- selon le type de verbe]
\begin{itemize}
\item \textbf{Verbes à particule séparable} (auswandern, zurückkehren, mitkommen...) : le \textbf{ge-} se glisse \textbf{entre la particule et le verbe} : \all{auswandern} $\rightarrow$ \all{aus\textbf{ge}wandert} ; \all{zurückkehren} $\rightarrow$ \all{zurück\textbf{ge}kehrt}.
\item \textbf{Verbes inséparables} (préfixes be-, ver-, er-, ent-, emp-, ge-, miss-) : \textbf{pas de ge-} du tout : \all{bekommen} $\rightarrow$ \all{bekommen} ; \all{vermissen} $\rightarrow$ \all{vermisst} ; \all{verstehen} $\rightarrow$ \all{verstanden}.
\item \textbf{Verbes en -ieren} : \textbf{pas de ge-} non plus : \all{studieren} $\rightarrow$ \all{studiert}.
\end{itemize}
\end{propriete}

\all{Kevin ist ausgewandert.} « Kevin a émigré. » (séparable : aus-ge-wandert)\par
\all{Er hat ein Stipendium bekommen.} « Il a obtenu une bourse. » (inséparable : pas de ge-)\par
\all{Er hat seine Familie vermisst.} « Sa famille lui a manqué. » (inséparable : vermisst)

\begin{attention}[Pas de ge- pour les verbes en -ieren et les verbes inséparables]
Erreurs classiques du francophone : *\all{Er hat gestudiert} au lieu de \all{Er hat studiert} ; *\all{Sie hat gebekommen} au lieu de \all{Sie hat bekommen}. Retenez : si le verbe commence par \textbf{be-, ver-, er-, ent-} ou finit par \textbf{-ieren}, le participe ne prend \textbf{jamais} de ge-.
\end{attention}

\courssub{Le choix de l'auxiliaire : haben ou sein ?}

C'est la grande question du Perfekt. La règle est claire :

\begin{propriete}[Choix de l'auxiliaire]
\textbf{On utilise \all{sein}} (conjugué au présent) avec :\par
1. les verbes de \textbf{mouvement} (qui impliquent un déplacement d'un point A vers un point B) : \all{gehen, kommen, fliegen, fahren, reisen, auswandern, zurückkehren, einreisen} ;\par
2. les verbes de \textbf{changement d'état} : \all{sterben} (mourir), \all{aufstehen} (se lever), \all{einschlafen} (s'endormir), \all{werden} (devenir) ;\par
3. les verbes \all{sein} (être) et \all{bleiben} (rester).\par\smallskip
\textbf{On utilise \all{haben}} avec \textbf{tous les autres verbes} : verbes d'action, d'opinion, de parole, verbes transitifs, verbes réfléchis...
\end{propriete}

\begin{exemplebox}[Comparer pour comprendre]
\all{Kevin ist nach München geflogen.} « Kevin est allé à Munich en avion. » — mouvement $\rightarrow$ \textbf{sein}.\par
\all{Er hat seine Familie vermisst.} « Sa famille lui a manqué. » — pas de mouvement $\rightarrow$ \textbf{haben}.\par
\all{Seine Mutter ist in Yaoundé geblieben.} « Sa mère est restée à Yaoundé. » — \all{bleiben} $\rightarrow$ \textbf{sein}.\par
\all{Sie hat jeden Tag an ihn gedacht.} « Elle a pensé à lui chaque jour. » — pensée $\rightarrow$ \textbf{haben}.
\end{exemplebox}

\begin{popfigure}
\begin{tikzpicture}[node distance=0.5cm]
\node[draw=popblue, fill=popblueL, rounded corners, minimum width=5.2cm, align=center] (h) {\textbf{mit haben}\\ (tous les autres verbes)};
\node[draw=popgreen, fill=popgreenL, rounded corners, minimum width=5.2cm, align=center, right=1.2cm of h] (s) {\textbf{mit sein}\\ (mouvement, changement d'état, sein/bleiben)};
\node[draw=popblue!60, fill=popcream, rounded corners, below=0.35cm of h, minimum width=5.2cm, align=left] (hl) {gemacht (fait)\\ gesagt (dit)\\ kontrolliert (contrôlé)\\ geweint (pleuré)\\ vermisst (manqué)};
\node[draw=popgreen!70, fill=popcream, rounded corners, below=0.35cm of s, minimum width=5.2cm, align=left] (sl) {ausgewandert (émigré)\\ gegangen (allé)\\ gekommen (venu)\\ zurückgekehrt (rentré)\\ geflogen (pris l'avion)};
\end{tikzpicture}
\legende{Deux colonnes à mémoriser : les participes avec \all{haben} et ceux avec \all{sein}.}
\end{popfigure}

\begin{attention}[Les deux erreurs préférées des francophones]
1. \textbf{Oublier le ge-} du participe : *\all{Er hat macht} au lieu de \all{Er hat gemacht}. Le participe allemand ressemble au participe français, mais il porte presque toujours sa petite coiffe \textbf{ge-}.\par
2. \textbf{Mal choisir l'auxiliaire} : *\all{Er hat ausgewandert} au lieu de \all{Er ist ausgewandert}. En français, « il a émigré » prend « avoir » ; en allemand, tout verbe de mouvement exige \textbf{sein}. Pensez au français soutenu « il \textbf{est} parti » : c'est le même réflexe !
\end{attention}

\courssub{Emploi du Perfekt : le passé de la conversation}

\begin{propriete}[Quand utilise-t-on le Perfekt ?]
Le Perfekt est le temps du passé de la \textbf{conversation}, des \textbf{dialogues}, des \textbf{lettres} et des \textbf{récits personnels}. Il correspond au passé composé français : \all{Was hast du gestern gemacht?} « Qu'as-tu fait hier ? »\par
À l'écrit littéraire et journalistique, l'allemand préfère souvent le Präteritum (que nous étudierons plus tard), sauf pour \all{sein} et \all{haben} dont le Präteritum (\all{war}, \all{hatte}) s'emploie partout.
\end{propriete}

Comparons les deux langues :

\begin{center}
\begin{tabularx}{\linewidth}{|X|X|X|}\hline
\rowcolor{popblueL} \textbf{Français} & \textbf{Deutsch} & \textbf{Remarque} \\ \hline
J'ai mangé. & Ich habe gegessen. & haben + participe en -en \\ \hline
Il est arrivé. & Er ist gekommen. & mouvement : sein \\ \hline
Nous avons voyagé. & Wir sind gereist. & mouvement : sein (pas « avoir » !) \\ \hline
Elle a étudié. & Sie hat studiert. & -ieren : pas de ge- \\ \hline
Tu as obtenu. & Du hast bekommen. & inséparable : pas de ge- \\ \hline
\end{tabularx}
\end{center}

\courssub{Un mini-texte pour tout réviser}

\begin{textebox}[Kevins erster Brief aus München]
Liebe Mama,\par
ich bin gestern gut in München angekommen. Der Flug hat vier Stunden gedauert. Am Flughafen hat ein Beamter meinen Pass kontrolliert, und dann bin ich mit dem Zug in die Stadt gefahren. Mein Onkel hat mich abgeholt und hat mir alles gezeigt. Ich habe heute schon Deutsch gesprochen und ein deutsches Brot gegessen — es hat gut geschmeckt ! Am Abend habe ich lange mit Onkel Thomas geredet. Er ist vor zwanzig Jahren aus Kamerun ausgewandert und ist hier geblieben. Er hat mir gesagt: „Das Heimweh geht vorbei.“ Ich habe geweint, aber nur ein bisschen. Morgen besuche ich die Universität. Ich habe große Hoffnungen für meine Zukunft.\par
Viele Grüße, dein Kevin
\end{textebox}

\textbf{Glossaire} : der Flug, die Flüge (le vol) — der Beamte, die Beamten (l'agent, le fonctionnaire) — abholen (aller chercher) — schmecken (avoir bon goût) — vorbeigehen (passer, cesser) — die Zukunft (l'avenir).

\textbf{Questions de compréhension} : 1. \all{Wie lange hat der Flug gedauert?} 2. \all{Wer hat Kevin abgeholt?} 3. \all{Wann ist Onkel Thomas ausgewandert?} Soulignez dans le texte tous les verbes au Perfekt et justifiez le choix de l'auxiliaire.

\courssub{Tableau récapitulatif des participes du texte}

\begin{center}
\begin{tabularx}{\linewidth}{|X|X|X|X|}\hline
\rowcolor{popblueL} \textbf{Infinitif} & \textbf{Participe} & \textbf{Auxiliaire} & \textbf{Français} \\ \hline
auswandern & ausgewandert & sein & émigré \\ \hline
bekommen & bekommen & haben & obtenu \\ \hline
fliegen & geflogen & sein & pris l'avion \\ \hline
vermissen & vermisst & haben & manqué \\ \hline
weinen & geweint & haben & pleuré \\ \hline
kontrollieren & kontrolliert & haben & contrôlé \\ \hline
ankommen & angekommen & sein & arrivé \\ \hline
bleiben & geblieben & sein & resté \\ \hline
essen & gegessen & haben & mangé \\ \hline
sagen & gesagt & haben & dit \\ \hline
\end{tabularx}
\end{center}

\begin{savaistu}[Le Perfekt, roi du sud !]
Dans le sud de l'Allemagne, en Autriche et en Suisse, on utilise le Perfekt presque partout à l'oral, même pour des verbes comme \all{sein} : \all{Ich bin müde gewesen} (« j'ai été fatigué ») au lieu du Präteritum \all{Ich war müde}. Un Nord-Allemand dira plutôt \all{Ich war müde}. Si tu écoutes un jour un Suisse ou un Bavarois parler de son passé, tu entendras le Perfekt à chaque phrase — tu es maintenant prêt à le comprendre !
\end{savaistu}

\courssub{Exercice résolu}

\begin{exoresolu}[Mettez les phrases au Perfekt]
1. Kevin (auswandern) nach Deutschland. \quad 2. Er (bekommen) ein Stipendium. \quad 3. Seine Mutter (weinen) am Flughafen. \quad 4. Der Beamte (kontrollieren) seinen Pass. \quad 5. Kevin (fliegen) nach München. \quad 6. Er (studieren) Informatik.
\tcblower
\textbf{Solution détaillée.}\par
1. \all{Kevin \textbf{ist} nach Deutschland \textbf{ausgewandert}.} — Verbe de mouvement $\rightarrow$ auxiliaire \all{sein} ; verbe séparable $\rightarrow$ ge- entre la particule et le radical : aus-ge-wandert.\par
2. \all{Er \textbf{hat} ein Stipendium \textbf{bekommen}.} — Pas de mouvement $\rightarrow$ \all{haben} ; verbe inséparable (be-) $\rightarrow$ pas de ge-.\par
3. \all{Seine Mutter \textbf{hat} am Flughafen \textbf{geweint}.} — Verbe faible régulier : ge + wein + t ; auxiliaire \all{haben}.\par
4. \all{Der Beamte \textbf{hat} seinen Pass \textbf{kontrolliert}.} — Verbe en -ieren $\rightarrow$ pas de ge- ; auxiliaire \all{haben}.\par
5. \all{Kevin \textbf{ist} nach München \textbf{geflogen}.} — Mouvement $\rightarrow$ \all{sein} ; verbe fort : participe en -en avec changement de voyelle (fliegen $\rightarrow$ geflogen).\par
6. \all{Er \textbf{hat} Informatik \textbf{studiert}.} — Verbe en -ieren $\rightarrow$ pas de ge- ; auxiliaire \all{haben}.
\end{exoresolu}

\begin{exercice}[À vous de jouer]
Mettez les phrases suivantes au Perfekt, puis traduisez-les en français :\par
1. Die Familie (bleiben) in Yaoundé. \quad 2. Kevin (schreiben) einen Brief. \quad 3. Wir (essen) ein deutsches Brot. \quad 4. Sie (zurückkehren) nach Kamerun. \quad 5. Ihr (vermissen) eure Freunde. \quad 6. Der Zug (kommen) um acht Uhr an.
\end{exercice}

\begin{corrige}[Exercice « À vous de jouer »]
1. \all{Die Familie ist in Yaoundé geblieben.} « La famille est restée à Yaoundé. » (\all{bleiben} $\rightarrow$ sein)\par
2. \all{Kevin hat einen Brief geschrieben.} « Kevin a écrit une lettre. » (verbe fort, haben)\par
3. \all{Wir haben ein deutsches Brot gegessen.} « Nous avons mangé un pain allemand. » (verbe fort, haben)\par
4. \all{Sie ist nach Kamerun zurückgekehrt.} « Elle est rentrée au Cameroun. » (mouvement $\rightarrow$ sein ; séparable : zurück-ge-kehrt)\par
5. \all{Ihr habt eure Freunde vermisst.} « Vos amis vous ont manqué. » (inséparable : pas de ge-)\par
6. \all{Der Zug ist um acht Uhr angekommen.} « Le train est arrivé à huit heures. » (mouvement $\rightarrow$ sein ; séparable : an-ge-kommen)
\end{corrige}

\begin{aretenir}[L'essentiel du Perfekt]
\begin{itemize}
\item \textbf{Formation} : \all{haben} ou \all{sein} au présent (2\textsuperscript{e} position) + participe passé \textbf{en toute fin de phrase} (Satzklammer).
\item \textbf{Participe des verbes faibles} : ge + radical + t (\all{gemacht}) ; des \textbf{verbes forts} : ge + radical souvent modifié + en (\all{gegangen}).
\item \textbf{Pas de ge-} pour les verbes en \textbf{-ieren} (\all{studiert}) et les verbes \textbf{inséparables} (\all{bekommen}, \all{vermisst}) ; ge- \textbf{à l'intérieur} des verbes séparables (\all{ausgewandert}).
\item \textbf{Auxiliaire sein} : mouvement, changement d'état, \all{sein} et \all{bleiben} ; \textbf{haben} pour tout le reste.
\item \textbf{Emploi} : le passé de la conversation et des récits personnels, comme le passé composé français.
\end{itemize}
\end{aretenir}

\coursec{Grammaire 2 : les subordonnées en dass et wenn}

Dans la section précédente, nous avons étudié le \cle{Perfekt}, le temps qui permet de raconter le passé : \all{Kevin ist nach Deutschland ausgewandert.} « Kevin a émigré en Allemagne. » Mais Kevin ne fait pas que raconter son passé : il exprime aussi ce qu'il \emph{pense}, ce qu'il \emph{espère} et ce qu'il fera \emph{quand} ses études seront terminées. Pour cela, l'allemand utilise des \cle{propositions subordonnées} introduites par \all{dass} « que » et \all{wenn} « quand / si ». Ces subordonnées obéissent à une règle capitale et sans exception : \textbf{le verbe conjugué se place en dernière position}. C'est l'une des plus grandes difficultés pour les francophones, car le français garde l'ordre sujet–verbe. Observons, réglons, entraînons-nous.

\courssub{Observation : deux phrases du texte}

\begin{textebox}[Extraits du texte « Kevins neues Leben in Deutschland »]
\all{Er sagt, dass er oft an Kamerun denkt.}

\all{Wenn er sein Studium beendet hat, will er zurückkehren.}
\end{textebox}

Traduisons mot à mot :

\begin{center}
\begin{tabularx}{\linewidth}{|X|X|X|}\hline
\rowcolor{popblueL} Allemand & Mot à mot & Français correct \\ \hline
\all{Er sagt, dass er oft an Kamerun denkt.} & Il dit, que lui souvent au Cameroun \textbf{pense}. & Il dit qu'il pense souvent au Cameroun. \\ \hline
\all{Wenn er sein Studium beendet hat, will er zurückkehren.} & Quand lui ses études terminé \textbf{a}, veut-il rentrer. & Quand il aura terminé ses études, il veut rentrer. \\ \hline
\end{tabularx}
\end{center}

\textbf{Que remarque-t-on ?}
\begin{itemize}
\item Après \all{dass} et après \all{wenn}, le verbe conjugué (\all{denkt}, \all{hat}) est rejeté \textbf{tout à la fin} de la subordonnée.
\item Une \textbf{virgule} sépare toujours la proposition principale de la subordonnée.
\item Quand la subordonnée est placée \emph{avant} la principale, le verbe de la principale vient \emph{immédiatement après la virgule} : \all{..., will er zurückkehren.}
\end{itemize}

\begin{popfigure}
\begin{tikzpicture}[font=\small]
\node[draw=popblue, fill=popblueL, rounded corners, align=center, minimum width=2.4cm] (main) at (0,0) {\textbf{Principale}\\ Er sagt,};
\node[draw=poporange, fill=poporangeL, rounded corners, align=center, minimum width=6.8cm, right=1.2cm of main] (sub) {\textbf{Subordonnée}\\ dass er oft an Kamerun \textbf{denkt.}};
\node[draw=popgreen, fill=popgreenL, rounded corners, align=center, below=1.1cm of sub] (verb) {verbe conjugué\\ en DERNIÈRE position};
\draw[-{Stealth[length=3mm]}, very thick, popgreen] (sub.south east) ++(-0.6,0) -- (verb.north east);
\draw[-{Stealth[length=3mm]}, very thick, popink] (main.east) -- node[above, align=center]{virgule\\ obligatoire} (sub.west);
\end{tikzpicture}
\legende{Structure de la subordonnée : la conjonction \all{dass} renvoie le verbe conjugué en fin de proposition.}
\end{popfigure}

\courssub{La règle d'or : le verbe conjugué en dernière position}

\begin{propriete}[Verbe en fin de subordonnée]
Dans une subordonnée introduite par \all{dass} ou \all{wenn}, le \cle{verbe conjugué} se place en \textbf{dernière position} de la subordonnée. L'ordre est donc : conjonction + sujet + compléments + \textbf{verbe}.

\all{Ich hoffe, dass ich ein Stipendium bekomme.} « J'espère que j'obtiendrai une bourse. »

\emph{Cas des temps composés :} avec un temps composé (Perfekt, Futur), le participe passé ou l'infinitif précède, et l'auxiliaire conjugué vient en tout dernier : \all{..., dass er sein Studium beendet hat.} / \all{..., wenn er zurückkehren wird.}

\emph{Cas des verbes à particule :} dans la subordonnée, la particule se \textbf{recolle} au verbe conjugué placé à la fin : \all{..., wenn er nach Kamerun zurückkehrt.} (et non *\all{kehrt zurück}).
\end{propriete}

\begin{attention}[Le piège de l'ordre français]
En français, on dit : « il dit qu'il \textbf{pense} souvent au Cameroun » (sujet–verbe). En allemand, le verbe part à la fin : \all{Er sagt, dass er oft an Kamerun \textbf{denkt}.} Ne jamais écrire *\all{dass er denkt oft an Kamerun} : c'est la faute la plus fréquente des francophones. Autre piège : oublier la virgule avant \all{dass} — elle est \textbf{obligatoire} en allemand, alors que le français n'en met pas devant « que ».
\end{attention}

\courssub{dass : la complétive après les verbes d'opinion et de parole}

\all{dass} correspond au « que » français. Il introduit une \cle{proposition complétive} après les verbes qui expriment la parole, la pensée, le savoir ou le sentiment :

\begin{center}
\begin{tabularx}{\linewidth}{|l|X|X|}\hline
\rowcolor{popblueL} Verbe & Sens & Exemple \\ \hline
\all{sagen} & dire & \all{Er sagt, dass er glücklich ist.} \\ \hline
\all{denken} & penser & \all{Sie denkt, dass er Recht hat.} \\ \hline
\all{hoffen} & espérer & \all{Wir hoffen, dass er zurückkommt.} \\ \hline
\all{wissen} & savoir & \all{Ich weiß, dass er im Ausland lebt.} \\ \hline
\all{glauben} & croire & \all{Er glaubt, dass alles gut wird.} \\ \hline
\all{meinen} & estimer & \all{Sie meint, dass Deutschland schön ist.} \\ \hline
\end{tabularx}
\end{center}

\begin{exemplebox}[Exemples traduits]
\all{Sie weiß, dass ihr Sohn im Ausland lebt.} « Elle sait que son fils vit à l'étranger. »

\all{Kevin hofft, dass seine Familie ihn besucht.} « Kevin espère que sa famille lui rendra visite. »

\all{Die Eltern denken, dass das Studium wichtig ist.} « Les parents pensent que les études sont importantes. »

\all{Er sagt, dass er oft an Kamerun denkt.} « Il dit qu'il pense souvent au Cameroun. »

\all{Ich glaube, dass er die Grenze gestern überquert hat.} « Je crois qu'il a franchi la frontière hier. » (Perfekt : participe \all{überquert} puis auxiliaire \all{hat}.)
\end{exemplebox}

\begin{attention}[dass $\neq$ das]
Ne pas confondre \all{dass} (conjonction, « que », avec \textbf{deux s}) et \all{das} (article ou pronom, « le / ce », avec \textbf{un seul s}).

\all{Das ist der Pass, den er braucht.} « Voici le passeport dont il a besoin. » (article)

\all{Er sagt, dass er den Pass braucht.} « Il dit qu'il a besoin du passeport. » (conjonction)

\emph{Test :} si l'on peut remplacer le mot par \all{dieses} « ce », on écrit \all{das} ; sinon, on écrit \all{dass}.
\end{attention}

\courssub{wenn : subordonnées temporelles et conditionnelles}

\all{wenn} a deux valeurs principales :

\begin{itemize}
\item \textbf{valeur temporelle} : « quand, lorsque, chaque fois que » — pour le présent, le futur et les habitudes. \all{Wenn ich Heimweh habe, rufe ich meine Mutter an.} « Quand j'ai le mal du pays, j'appelle ma mère. »
\item \textbf{valeur conditionnelle} : « si ». \all{Wenn er zurückkehrt, hilft er seinem Land.} « S'il rentre, il aidera son pays. »
\end{itemize}

\begin{exemplebox}[Exemples traduits]
\all{Wenn ich Heimweh habe, rufe ich meine Mutter an.} « Quand j'ai le mal du pays, j'appelle ma mère. »

\all{Wenn er zurückkehrt, hilft er seinem Land.} « S'il rentre, il aidera son pays. »

\all{Wenn er sein Studium beendet hat, will er zurückkehren.} « Quand il aura terminé ses études, il veut rentrer. »

\all{Wenn es regnet, bleibe ich zu Hause.} « Quand il pleut, je reste à la maison. »
\end{exemplebox}

\begin{propriete}[wenn / als / wann : trois mots pour « quand »]
\begin{itemize}
\item \all{wenn} = « quand / si » : présent, futur, habitudes (même au passé : « chaque fois que ») et conditions.
\item \all{als} = « quand, lorsque » : \textbf{passé, événement unique}. \all{Als er nach Deutschland kam, war er 18 Jahre alt.} « Quand il est arrivé en Allemagne, il avait 18 ans. »
\item \all{wann} = « quand » : uniquement dans les \textbf{questions} (directes ou indirectes). \all{Wann kommst du zurück?} « Quand rentres-tu ? »
\end{itemize}
\end{propriete}

\begin{popfigure}
\begin{tikzpicture}[font=\small]
\node[draw=popink, fill=popblueL, rounded corners, align=center, minimum width=2.6cm] (q) at (0,0) {\textbf{« quand »}\\ en allemand ?};
\node[draw=popgreen, fill=popgreenL, rounded corners, align=center, right=2.6cm of q, yshift=1.6cm] (wenn) {\textbf{wenn}\\ présent, futur,\\ habitude, si\\ \all{Wenn er kommt, ...}};
\node[draw=poporange, fill=poporangeL, rounded corners, align=center, right=2.6cm of q] (als) {\textbf{als}\\ passé,\\ événement unique\\ \all{Als er kam, ...}};
\node[draw=poppurple, fill=poppurpleL, rounded corners, align=center, right=2.6cm of q, yshift=-1.6cm] (wann) {\textbf{wann}\\ question\\ \all{Wann kommt er?}};
\draw[-{Stealth[length=3mm]}, thick, popgreen] (q.east) -- (wenn.west);
\draw[-{Stealth[length=3mm]}, thick, poporange] (q.east) -- (als.west);
\draw[-{Stealth[length=3mm]}, thick, poppurple] (q.east) -- (wann.west);
\end{tikzpicture}
\legende{Choisir entre \all{wenn}, \all{als} et \all{wann} selon le sens.}
\end{popfigure}

\courssub{Place de la subordonnée : après ou avant la principale}

La subordonnée peut occuper deux positions :

\begin{enumerate}
\item \textbf{Après la principale} (le cas le plus fréquent) : la principale garde son verbe en 2\textsuperscript{e} position. \all{Er sagt, dass er oft an Kamerun denkt.}
\item \textbf{Avant la principale} : la subordonnée occupe alors la 1\textsuperscript{re} position de la phrase ; le verbe de la principale vient \textbf{immédiatement après la virgule}, puis le sujet. \all{Wenn er zurückkehrt, will er helfen.}
\end{enumerate}

\begin{attention}[Subordonnée en tête = verbe de la principale juste après la virgule]
Quand la subordonnée est placée en tête, le verbe de la principale suit immédiatement la virgule : \all{Wenn er zurückkehrt, \textbf{will} er helfen.} et non *\all{Wenn er zurückkehrt, er will helfen.} La subordonnée entière compte comme la 1\textsuperscript{re} position ; le verbe de la principale reste en 2\textsuperscript{e} position.
\end{attention}

\begin{center}
\begin{tabularx}{\linewidth}{|X|X|}\hline
\rowcolor{popblueL} Subordonnée après & Subordonnée avant \\ \hline
\all{Ich rufe meine Mutter an, wenn ich Heimweh habe.} & \all{Wenn ich Heimweh habe, rufe ich meine Mutter an.} \\ \hline
\all{Er will helfen, wenn er zurückkehrt.} & \all{Wenn er zurückkehrt, will er helfen.} \\ \hline
\all{Sie ist froh, dass ihr Sohn studiert.} & \all{Dass ihr Sohn studiert, macht sie froh.} \\ \hline
\end{tabularx}
\end{center}

\courssub{Comparaison français / allemand}

\begin{center}
\begin{tabularx}{\linewidth}{|X|X|X|}\hline
\rowcolor{popblueL} Français & Deutsch & Remarque \\ \hline
Il dit qu'il pense souvent au Cameroun. & \all{Er sagt, dass er oft an Kamerun denkt.} & Verbe rejeté à la fin en allemand. \\ \hline
J'espère que j'obtiendrai une bourse. & \all{Ich hoffe, dass ich ein Stipendium bekomme.} & Futur français = présent allemand possible. \\ \hline
Quand j'ai le mal du pays, j'appelle ma mère. & \all{Wenn ich Heimweh habe, rufe ich meine Mutter an.} & Verbe de la principale après la virgule. \\ \hline
Quand il est arrivé, il avait 18 ans. & \all{Als er ankam, war er 18 Jahre alt.} & Passé unique : \all{als}, pas \all{wenn}. \\ \hline
\end{tabularx}
\end{center}

\courssub{Méthode : construire une subordonnée en dass ou wenn}

\begin{methode}[Quatre étapes pour ne jamais se tromper]
\begin{enumerate}
\item \textbf{Identifier la conjonction} : \all{dass} (« que ») ou \all{wenn} (« quand / si »).
\item \textbf{Écrire sujet + compléments} dans l'ordre habituel : \all{dass ich ein Stipendium ...}
\item \textbf{Placer le verbe conjugué tout à la fin} : \all{dass ich ein Stipendium bekomme.} (avec un temps composé : participe puis auxiliaire : \all{..., dass er angekommen ist.} ; avec un verbe à particule : la particule se recolle : \all{..., wenn er zurückkehrt.})
\item \textbf{Vérifier la virgule} obligatoire entre les deux propositions, et si la subordonnée est en tête, placer le verbe de la principale juste après : \all{Wenn ..., kommt er.}
\end{enumerate}
\end{methode}

\begin{textebox}[Kevins Pläne für die Zukunft]
\all{Kevin sitzt in seinem kleinen Zimmer in München und denkt an seine Zukunft. Er sagt, dass er oft Heimweh hat, besonders am Wochenende. Wenn er an Yaoundé denkt, sieht er den Markt, die Freunde und die Familie vor sich. Seine Mutter hofft, dass er gesund bleibt und gut lernt. Kevin weiß, dass das Studium in Deutschland schwer ist, aber er glaubt, dass er es schaffen kann. Wenn er sein Studium beendet hat, will er nach Kamerun zurückkehren. Er meint, dass sein Land junge Ingenieure braucht. Wenn er zurückkehrt, will er eine Firma gründen und viele Jugendliche ausbilden. Seine Freunde in Deutschland sagen, dass er ein Vorbild ist. Kevin lacht und antwortet: „Wenn ich Erfolg habe, komme ich oft zurück und besuche euch!" Er weiß, dass die Emigration nicht einfach ist, aber er hofft, dass seine Geschichte anderen Mut macht.}

\emph{Glossaire :} das Heimweh (nur Sg.) : le mal du pays ; der Markt, die Märkte : le marché ; schaffen : réussir ; die Firma, die Firmen : l'entreprise ; gründen : fonder ; ausbilden : former ; das Vorbild, die Vorbilder : le modèle ; der Erfolg, die Erfolge : le succès ; der Mut (nur Sg.) : le courage.

\emph{Questions :} 1) \all{Wo wohnt Kevin?} 2) \all{Was hofft seine Mutter?} 3) \all{Was will Kevin machen, wenn er zurückkehrt?} 4) \all{Warum ist Kevin ein Vorbild?}
\end{textebox}

\courssub{Exercice résolu et entraînement}

\begin{exoresolu}[Remets les mots dans l'ordre]
Énoncé. Construis une phrase correcte avec : \all{hofft / Kevin / dass / ein Stipendium / er / bekommt}. Puis transforme : \all{Er kehrt zurück. Dann hilft er seinem Land.} en une phrase avec \all{wenn} en tête.

\tcblower

Solution détaillée.
1) Conjonction : \all{dass}. Sujet + compléments : \all{er ein Stipendium}. Verbe à la fin : \all{bekommt}. Phrase : \all{Kevin hofft, dass er ein Stipendium bekommt.} « Kevin espère qu'il obtiendra une bourse. »

2) \all{wenn} en tête, verbe de la subordonnée à la fin (la particule \all{zurück} se recolle : \all{zurückkehrt}), verbe de la principale après la virgule : \all{Wenn er zurückkehrt, hilft er seinem Land.} « Quand il rentrera, il aidera son pays. »
\end{exoresolu}

\begin{exercice}[À toi de jouer]
1. Traduis en allemand : a) Elle sait que son fils vit à l'étranger. b) Quand j'ai le mal du pays, j'appelle ma mère. c) S'il obtient une bourse, il restera en Allemagne.

2. Choisis entre \all{wenn}, \all{als} et \all{wann} : a) \all{... ich ein Kind war, lebte ich in Douala.} b) \all{... kommt er zurück?} c) \all{... er Zeit hat, besucht er uns.}

3. Mets le verbe à la bonne place : a) \all{Er sagt, dass er oft an Kamerun (denken).} b) \all{Ich weiß, dass sie gestern (ankommen).} (Perfekt)

4. Relie les deux phrases avec \all{dass} : \all{Kevin hofft. Seine Familie besucht ihn bald.}
\end{exercice}

\begin{corrige}[Exercice]
1. a) \all{Sie weiß, dass ihr Sohn im Ausland lebt.} b) \all{Wenn ich Heimweh habe, rufe ich meine Mutter an.} c) \all{Wenn er ein Stipendium bekommt, bleibt er in Deutschland.}

2. a) \all{Als} (passé, événement unique) ; b) \all{Wann} (question) ; c) \all{Wenn} (habitude / condition).

3. a) \all{Er sagt, dass er oft an Kamerun denkt.} (verbe conjugué en dernière position) ; b) \all{Ich weiß, dass sie gestern angekommen ist.} (Perfekt : participe puis auxiliaire \all{ist}, car \all{ankommen} est un verbe de mouvement).

4. \all{Kevin hofft, dass seine Familie ihn bald besucht.}
\end{corrige}

\begin{aretenir}[L'essentiel]
\begin{itemize}
\item Après \all{dass} et \all{wenn}, le verbe conjugué va en \textbf{dernière position} : \all{..., dass er oft an Kamerun denkt.}
\item Avec un temps composé : participe ou infinitif, puis auxiliaire en tout dernier : \all{..., dass er angekommen ist.}
\item Avec un verbe à particule, la particule se recolle au verbe en fin de subordonnée : \all{..., wenn er zurückkehrt.}
\item \all{dass} = « que » après \all{sagen, denken, hoffen, wissen, glauben} ; il s'écrit avec \textbf{deux s}.
\item \all{wenn} = « quand / si » (présent, futur, habitude) ; \all{als} = « quand » au passé (événement unique) ; \all{wann} = « quand » dans les questions.
\item Subordonnée en tête $\Rightarrow$ verbe de la principale juste après la virgule : \all{Wenn er zurückkehrt, will er helfen.}
\item La virgule entre les deux propositions est \textbf{obligatoire}.
\end{itemize}
\end{aretenir}

\coursec{Méthode du commentaire et commentaire modèle}

Dans la section précédente, nous avons appris à construire des subordonnées avec \all{dass} et \all{wenn} : \all{Ich finde, dass Kevin mutig ist.} / \all{Wenn ich die Möglichkeit hätte, würde ich auch auswandern.} Ces structures ne sont pas un exercice gratuit : ce sont précisément les phrases dont vous aurez besoin pour \cle{commenter un texte}. En effet, au baccalauréat comme dans les évaluations de Première, on vous demandera souvent : \all{Fassen Sie den Text zusammen und geben Sie Ihre Meinung.} (Résumez le texte et donnez votre avis.) Cette section vous donne la méthode complète, étape par étape, puis un commentaire modèle rédigé sur notre texte support \emph{Kevins neues Leben in Deutschland}.

\courssub{Qu'est-ce qu'un commentaire de texte ?}

\begin{definition}[Der Kommentar — le commentaire de texte]
Un \cle{commentaire de texte} (\all{der Kommentar, die Kommentare}) est une production écrite ou orale qui comporte quatre parties :
\begin{enumerate}
\item la \cle{présentation} du texte (\all{die Einleitung}) : le thème, les personnages, la situation ;
\item le \cle{résumé} du contenu (\all{der Inhalt}) : l'idée principale et les étapes du récit, racontées \textbf{avec ses propres mots}, au parfait ;
\item l'\cle{analyse} (\all{die Analyse}) : les sentiments des personnages, les problèmes, les solutions ;
\item l'\cle{avis personnel} (\all{die Meinung}) : ce que vous pensez du texte, avec une justification.
\end{enumerate}
\end{definition}

\begin{attention}[Ce qu'un commentaire n'est PAS]
Un commentaire n'est \textbf{jamais} une recopie du texte. Recopier des phrases entières du texte support est sanctionné : l'examinateur veut voir \emph{votre} allemand, pas celui du texte. La règle d'or : \textbf{fermez le texte, puis racontez l'histoire avec vos mots}, en utilisant le vocabulaire du thème (\all{auswandern, das Ausland, das Heimweh, die Hoffnung, der Pass, die Grenze}) et le parfait.
\end{attention}

\courssub{Les quatre étapes du commentaire}

\begin{methode}[Commenter un texte en quatre étapes]
\begin{enumerate}
\item \textbf{Lire deux fois le texte.} À la première lecture, cherchez le sens global : de qui parle-t-on ? Où ? Que se passe-t-il ? À la deuxième lecture, soulignez les mots du thème et les moments importants du récit.
\item \textbf{Rédiger l'introduction (1 à 2 phrases).} Présentez le texte : \all{Der Text handelt von...} (Le texte parle de...) + le thème : \all{Das Thema ist die Emigration.} (Le thème est l'émigration.)
\item \textbf{Rédiger le résumé (3 à 4 phrases) au parfait.} Racontez les étapes du récit avec les connecteurs chronologiques : \all{zuerst} (d'abord), \all{dann} (puis), \all{danach} (ensuite), \all{schließlich} (finalement). N'oubliez pas : le participe II va \textbf{à la fin} de la phrase.
\item \textbf{Donner votre avis (2 à 3 phrases) avec justification.} Utilisez \all{Ich finde, dass...} ou \all{Meiner Meinung nach...} suivi d'une raison introduite par \all{weil} ou d'un souhait introduit par \all{wenn}.
\end{enumerate}
\end{methode}

\begin{popfigure}
\begin{center}
\begin{tikzpicture}[
node distance=0.55cm,
etape/.style={rectangle, rounded corners=3pt, draw=popblue, fill=popblueL, text width=9.5cm, align=center, inner sep=6pt, font=\small},
fleche/.style={-{Stealth[length=2.5mm]}, thick, poporange}
]
\node[etape] (e1) {\textbf{1. Einleitung} — Der Text handelt von... / Das Thema ist...};
\node[etape, below=of e1] (e2) {\textbf{2. Inhalt} — Résumé au Perfekt : zuerst..., dann..., danach..., schließlich...};
\node[etape, below=of e2] (e3) {\textbf{3. Analyse} — Gefühle und Probleme : Heimweh, Hoffnung, Schwierigkeiten};
\node[etape, below=of e3] (e4) {\textbf{4. Meinung} — Ich finde, dass... / Meiner Meinung nach... / Wenn ich...};
\draw[fleche] (e1) -- (e2);
\draw[fleche] (e2) -- (e3);
\draw[fleche] (e3) -- (e4);
\end{tikzpicture}
\end{center}
\legende{Organigramme du commentaire de texte : quatre étapes obligatoires, de la présentation à l'avis personnel.}
\end{popfigure}

\courssub{Les expressions utiles du commentaire}

Pour chaque étape, il existe des \cle{phrases toutes faites} (\all{die Redemittel}) qu'il faut mémoriser. Elles garantissent un allemand correct et donnent une structure claire à votre production.

\begin{center}
\begin{tabularx}{\linewidth}{|X|X|X|}
\hline
\rowcolor{popblueL} Étape & Deutsch & Français \\
\hline
Présentation & \all{Der Text handelt von einem jungen Kameruner.} & Le texte parle d'un jeune Camerounais. \\
\hline
Présentation & \all{Das Thema des Textes ist die Emigration.} & Le thème du texte est l'émigration. \\
\hline
Résumé & \all{Der Autor erzählt, dass Kevin nach Deutschland ausgewandert ist.} & L'auteur raconte que Kevin a émigré en Allemagne. \\
\hline
Résumé & \all{Zuerst hat er Heimweh gehabt, dann hat er Freunde gefunden.} & D'abord il a eu le mal du pays, puis il a trouvé des amis. \\
\hline
Analyse & \all{Kevin hat viele Probleme, aber er gibt nicht auf.} & Kevin a beaucoup de problèmes, mais il n'abandonne pas. \\
\hline
Avis & \all{Meiner Meinung nach ist Kevin sehr mutig.} & À mon avis, Kevin est très courageux. \\
\hline
Avis & \all{Ich finde es wichtig, dass junge Leute studieren können.} & Je trouve important que les jeunes puissent étudier. \\
\hline
Avis + souhait & \all{Wenn ich die Möglichkeit hätte, würde ich auch gern im Ausland studieren.} & Si j'avais la possibilité, j'aimerais aussi étudier à l'étranger. \\
\hline
\end{tabularx}
\end{center}

\begin{propriete}[Grammaire en action : ce que le commentaire mobilise]
Le commentaire est une \textbf{application directe} des deux points de grammaire de cette leçon :
\begin{itemize}
\item \textbf{Le parfait} pour le résumé : auxiliaire \all{haben} ou \all{sein} à la deuxième place, participe II à la fin : \all{Kevin \textbf{ist} nach München \textbf{ausgewandert}.} / \all{Er \textbf{hat} ein Stipendium \textbf{bekommen}.}
\item \textbf{Les subordonnées en \all{dass} et \all{wenn}} pour l'avis : verbe conjugué \textbf{à la fin} de la subordonnée : \all{Ich finde, dass Kevin sehr mutig \textbf{ist}.} / \all{Wenn ich die Möglichkeit \textbf{hätte}, ...}
\end{itemize}
Attention aux verbes à particule séparable au parfait : \all{auswandern} donne \all{ist ausgewandert} (le \all{-ge-} se place entre la particule et le radical).
\end{propriete}

\begin{exemplebox}[Deux phrases modèles à retenir par cœur]
\all{Der Text handelt von einem jungen Kameruner, der nach Deutschland ausgewandert ist.}

« Le texte parle d'un jeune Camerounais qui a émigré en Allemagne. »

\medskip
\all{Ich finde es mutig, dass Kevin allein ins Ausland gegangen ist.}

« Je trouve courageux que Kevin soit parti seul à l'étranger. »

\medskip
Observez la deuxième phrase : la construction \all{Ich finde es} + adjectif + \all{dass}... correspond au français « je trouve + adjectif + que... ». Le petit mot \all{es} est obligatoire en allemand : on ne peut pas dire \all{Ich finde mutig, dass...}.
\end{exemplebox}

\courssub{Commentaire modèle rédigé}

Voici maintenant un commentaire complet (8 à 10 lignes) sur notre texte support. Lisez-le attentivement : il respecte exactement les quatre étapes de l'organigramme.

\begin{textebox}[Kommentar — Commentaire modèle sur « Kevins neues Leben in Deutschland »]
Der Text handelt von Kevin, einem jungen Kameruner, der nach München ausgewandert ist. Das Thema ist die Emigration.

Kevin hat ein Stipendium bekommen und hat in Deutschland ein Studium begonnen. Zuerst hat er Heimweh gehabt, weil er seine Familie und seine Freunde vermisst hat. Dann hat er eine Wohnung gefunden und neue Freunde kennengelernt. Schließlich hat er sich an das neue Leben gewöhnt, aber er hat sein Land nicht vergessen.

Ich finde, dass Kevin sehr mutig ist, weil er allein ins Ausland gegangen ist. Meiner Meinung nach ist Emigration eine große Chance, aber auch ein Risiko. Wenn ich die Möglichkeit hätte, würde ich auch gern in Deutschland studieren, weil die Universitäten dort sehr gut sind. Aber ich würde nach dem Studium gern nach Kamerun zurückkehren, um meinem Land zu helfen.
\end{textebox}

\textbf{Glossaire :} \emph{das Stipendium, die Stipendien} (la bourse d'études) ; \emph{vermissen} (manquer à quelqu'un, regretter) ; \emph{kennenlernen} (faire la connaissance de) ; \emph{sich gewöhnen an + Akkusativ} (s'habituer à) ; \emph{mutig} (courageux) ; \emph{das Risiko, die Risiken} (le risque) ; \emph{zurückkehren} (retourner, revenir).

\medskip
\textbf{Analyse du modèle.} Relevons ensemble les outils utilisés :

\begin{center}
\begin{tabularx}{\linewidth}{|X|X|}
\hline
\rowcolor{popblueL} Marqueur relevé dans le modèle & Fonction \\
\hline
\all{Der Text handelt von...} & présenter le texte (Einleitung) \\
\hline
\all{zuerst, dann, schließlich} & ordonner le résumé chronologiquement \\
\hline
\all{weil er seine Familie vermisst hat} & expliquer un sentiment (cause) \\
\hline
\all{Ich finde, dass...} / \all{Meiner Meinung nach...} & exprimer son avis \\
\hline
\all{Wenn ich die Möglichkeit hätte, würde ich...} & exprimer un souhait (Konjunktiv II) \\
\hline
\all{aber ich würde gern zurückkehren, um... zu helfen} & nuancer et conclure \\
\hline
\end{tabularx}
\end{center}

Comptez les verbes au parfait dans le résumé : \all{hat bekommen, hat begonnen, hat gehabt, hat vermisst, hat gefunden, hat kennengelernt, hat sich gewöhnt, hat vergessen}. C'est exactement ce que l'examinateur attend : un résumé au passé, avec des participes II correctement placés en fin de phrase. Remarquez aussi que tous ces verbes se construisent avec \all{haben}, car ils expriment des actions sans déplacement ; en revanche, dans l'introduction, \all{auswandern} exige l'auxiliaire \all{sein} : \all{der nach München ausgewandert ist}.

\begin{attention}[Les critères de notation au baccalauréat]
Au Bac, le résumé doit rester \textbf{bref} : 3 à 4 phrases maximum. L'essentiel de la note porte sur :
\begin{itemize}
\item l'\textbf{exactitude grammaticale} : parfait correct (bon auxiliaire \all{haben}/\all{sein}, participe II en fin de phrase), place du verbe en position 2 dans les principales et en fin dans les subordonnées ;
\item l'\textbf{usage du vocabulaire du thème} : employez les mots de la leçon (\all{das Ausland, das Heimweh, die Hoffnung, auswandern}) plutôt que des paraphrases vagues ;
\item la \textbf{personnalisation} : un avis \emph{à vous}, justifié par \all{weil} ou \all{denn}, vaut toujours plus de points qu'une phrase générale copiée.
\end{itemize}
Erreur fréquente des francophones : oublier le \all{es} dans \all{Ich finde es wichtig, dass...} ou mettre le verbe de la subordonnée en position 2 comme en français : écrivez \all{dass Kevin mutig \textbf{ist}}, et non \all{dass Kevin \textbf{ist} mutig}.
\end{attention}

\begin{exoresolu}[Transformer des notes en commentaire]
Énoncé : voici des notes sur le texte. Rédigez un mini-commentaire de 4 phrases en utilisant : \all{handeln von} — \all{zuerst} — \all{Ich finde, dass...} — \all{Wenn ich...}.

Notes : Kevin / Kamerun $\rightarrow$ München / Stipendium / Heimweh / nicht aufgeben.

\tcblower
Solution détaillée :

\all{Der Text handelt von Kevin, einem jungen Kameruner, der nach München ausgewandert ist. Zuerst hat er ein Stipendium bekommen, aber dann hat er Heimweh gehabt. Trotzdem hat er nicht aufgegeben. Ich finde, dass Kevin ein gutes Beispiel für junge Leute ist. Wenn ich die Chance hätte, würde ich auch gern im Ausland lernen.}

« Le texte parle de Kevin, un jeune Camerounais qui a émigré à Munich. D'abord il a obtenu une bourse, mais ensuite il a eu le mal du pays. Pourtant il n'a pas abandonné. Je trouve que Kevin est un bon exemple pour les jeunes. Si j'avais la chance, j'aimerais aussi apprendre à l'étranger. »

Points de méthode : (1) la phrase 1 présente le texte ; (2) les phrases 2 et 3 résument au parfait avec \all{zuerst} et \all{dann}, participes II en fin (\all{bekommen, gehabt, aufgegeben}) ; (3) la phrase 4 donne un avis avec \all{dass} + verbe final ; (4) la phrase 5 exprime un souhait avec \all{wenn} + Konjunktiv II.
\end{exoresolu}

\begin{exercice}[À vous de commenter]
Rédigez un commentaire de 6 à 8 lignes sur le texte \emph{Kevins neues Leben in Deutschland}, en suivant l'organigramme : introduction (1 phrase), résumé au parfait (3 phrases avec \all{zuerst, dann, schließlich}), avis personnel justifié (2 phrases avec \all{dass} et \all{wenn}). Soulignez dans votre production tous les participes II et les verbes finaux des subordonnées.
\end{exercice}

\begin{corrige}[Exercice — éléments de correction]
Il n'existe pas une seule bonne réponse, mais tout bon commentaire contient :
\begin{itemize}
\item une introduction du type : \all{Der Text handelt von einem jungen Kameruner, der nach Deutschland ausgewandert ist. Das Thema ist die Emigration.}
\item un résumé au parfait, par exemple : \all{Zuerst hat Kevin ein Stipendium bekommen. Dann ist er nach München geflogen und hat ein Studium begonnen. Schließlich hat er sich an das neue Leben gewöhnt.}
\item un avis justifié : \all{Ich finde, dass Kevin sehr mutig ist, weil er seine Familie verlassen hat. Wenn ich die Möglichkeit hätte, würde ich auch gern in Deutschland studieren, weil man dort gute Universitäten hat.}
\end{itemize}
Vérifiez : auxiliaire correct (\all{sein} pour \all{fliegen, auswandern} ; \all{haben} pour \all{bekommen, beginnen}), participe II en fin de phrase, verbe conjugué en fin de subordonnée après \all{dass}, \all{weil} et \all{wenn}.
\end{corrige}

\courssub{Entraînement oral : présenter son commentaire en une minute}

\begin{experience}[Der Ein-Minuten-Kommentar — le commentaire en une minute]
Travail par binômes, en trois temps :
\begin{enumerate}
\item \textbf{Préparation (5 minutes).} Rédigez non pas un texte complet, mais une \textbf{fiche de mots-clés} : 4 lignes maximum, une par étape (Einleitung — Inhalt — Analyse — Meinung), avec seulement 3 ou 4 mots par ligne. Exemple : \emph{Einleitung : Kevin — Kameruner — München — Emigration}.
\item \textbf{Présentation (1 minute par élève).} Présentez votre commentaire à votre camarade \textbf{sans lire de phrases complètes} : regardez vos mots-clés et reconstruisez les phrases à voix haute. Le camarade chronomètre et vérifie la présence des quatre étapes.
\item \textbf{Retour (2 minutes).} Le camarade donne un point fort et un point à améliorer, en allemand si possible : \all{Dein Perfekt war gut.} (Ton parfait était bon.) / \all{Die Verbstellung im dass-Nebensatz war falsch.} (La place du verbe dans la subordonnée en \all{dass} était fausse.)
\end{enumerate}
Conseil du professeur : commencez toujours par la phrase toute faite \all{Der Text handelt von...} — elle vous lance et vous met en confiance.
\end{experience}

\begin{aretenir}[L'essentiel de la section]
\begin{itemize}
\item Un commentaire = 4 étapes : \textbf{Einleitung} (\all{Der Text handelt von...}), \textbf{Inhalt} (résumé au parfait avec \all{zuerst, dann, danach, schließlich}), \textbf{Analyse} (sentiments et problèmes), \textbf{Meinung} (\all{Ich finde, dass...} / \all{Meiner Meinung nach...}).
\item On ne recopie jamais le texte : on résume avec ses propres mots, en 3 à 4 phrases.
\item Le parfait : auxiliaire en position 2, participe II à la fin ; les subordonnées en \all{dass}, \all{weil}, \all{wenn} : verbe conjugué à la fin.
\item À l'oral : une fiche de mots-clés suffit ; parler une minute sans lire est un entraînement idéal pour l'examen.
\end{itemize}
\end{aretenir}

\coursec{Production guidée : mon histoire d'émigration}

Après avoir analysé le texte de Kevin, travaillé le vocabulaire de l'émigration et révisé les structures grammaticales du parfait et des subordonnées en \all{dass} et \all{wenn}, vous êtes maintenant prêt à produire votre propre récit. Cette section vous guide pas à pas, de l'étude d'un modèle à la rédaction de votre texte personnel, en passant par l'auto-évaluation et l'expression orale.

\courssub{Modèle de production écrite : « Mein Weg ins Ausland »}

Avant d'écrire, observons un texte modèle qui respecte toutes les contraintes imposées (plan en quatre parties, 6 verbes au parfait dont 5 en phrase principale, une subordonnée en \all{dass}, une en \all{wenn}, six mots du lexique thématique). Lisez-le attentivement, puis consultez le glossaire et répondez aux questions de compréhension.

\begin{textebox}{Modèle : Mein Weg ins Ausland}
Letztes Jahr bin ich von Douala nach München geflogen. Ich wollte in Deutschland Informatik studieren. Zuerst habe ich meinen Pass und mein Visum bei der Botschaft beantragt. Meine Eltern haben mich zum Flughafen begleitet. Der Abflug war ein schwieriger Moment, denn meine Mutter hat geweint. An der Grenze habe ich die Passkontrolle ohne Probleme passiert. Nach der Ankunft in München hatte ich zuerst großes Heimweh nach meiner Familie und dem kamerunischen Essen. Aber ich habe große Hoffnung gefasst, weil die Universität sehr gut ist. Ich weiß, \textbf{dass} das Leben im Ausland am Anfang schwer ist. \textbf{Wenn} ich mein Studium beendet habe, werde ich nach Kamerun zurückkehren und dort arbeiten.
\end{textebox}

\textbf{Glossaire du modèle}
\begin{center}
\begin{tabularx}{\linewidth}{|X|X|X|}
\hline
\rowcolor{popblueL} \textbf{Allemand} & \textbf{Français} & \textbf{Remarques} \\
\hline
\all{letztes Jahr} & l'année dernière & expression de temps (cas accusatif) \\
\all{von Douala nach München fliegen} & voler de Douala à Munich & \all{fliegen} (fort) : \all{flog}, \all{geflogen} ; auxiliaire \all{sein} \\
\all{Informatik studieren} & étudier l'informatique & \all{studieren} (faible, en \all{-ieren}) $\rightarrow$ Part. \all{studiert} (sans \all{ge-}) \\
\all{beantragen} & demander, solliciter & verbe à particule inséparable \all{be-} $\rightarrow$ Part. \all{beantragt} (sans \all{ge-}) \\
\all{die Botschaft, -en} & l'ambassade & \\
\all{begleiten} & accompagner & \all{hat begleitet} (aux. \all{haben}) \\
\all{der Abflug, -e} & le décollage, le départ (avion) & \all{Ab-} + \all{Flug} \\
\all{der Moment, -e} & le moment & \\
\all{die Grenze, -n} & la frontière & mot du lexique thématique \\
\all{die Passkontrolle, -n} & le contrôle des passeports & composé : \all{Pass} + \all{Kontrolle} \\
\all{passieren} & passer, franchir & ici transitif (contrôle) $\rightarrow$ aux. \all{haben} \\
\all{die Ankunft, -e} & l'arrivée & nominalisation de \all{ankommen} \\
\all{das Heimweh} (invar.) & le mal du pays, la nostalgie & \all{Heim} (foyer) + \all{Weh} (douleur) ; \all{Heimweh haben} \\
\all{die Hoffnung, -en} & l'espoir & mot du lexique thématique ; \all{Hoffnung fassen} (concevoir de l'espoir) \\
\all{wissen, wusste, gewusst} & savoir (un fait) & verbe irrégulier (Prät. \all{wusste}) \\
\all{dass} & que (conjonction) & verbe à la fin \\
\all{das Ausland} & l'étranger & mot du lexique thématique \\
\all{beenden, hat beendet} & terminer, achever & verbe à particule inséparable \all{be-} \\
\all{zurückkehren, ist zurückgekehrt} & rentrer, retourner & verbe à particule séparable $\rightarrow$ Part. \all{zurückgekehrt}, aux. \all{sein} \\
\all{sich einleben, hat sich eingelebt} & s'installer, s'habituer & verbe réfléchi, particule séparable \all{ein-} \\
\hline
\end{tabularx}
\end{center}

\begin{exercice}{Compréhension du modèle}
\begin{enumerate}
\item Quelles sont les quatre étapes du récit respectées dans ce texte ?
\item Relevez six verbes conjugués au \textbf{Perfekt} (phrase principale ou subordonnée) et donnez leur infinitif ainsi que leur auxiliaire (\all{haben} ou \all{sein}).
\item Identifiez la subordonnée en \all{dass} et celle en \all{wenn}. Quel est le temps du verbe dans la subordonnée en \all{wenn} ? Pourquoi ?
\item Combien de mots du lexique thématique (liste de la section 3) trouvez-vous ? Citez-les.
\end{enumerate}
\end{exercice}

\begin{corrige}{Exercice Compréhension du modèle}
\begin{enumerate}
\item 1) Le départ (l. 1: \all{letztes Jahr bin ich ... geflogen}), 2) Les préparatifs (l. 3-4: \all{habe ... beantragt, haben ... begleitet}), 3) Les sentiments (l. 5-8: \all{war ... schwierig, hat geweint, habe ... passiert, hatte ... Heimweh, habe ... Hoffnung gefasst}), 4) Les projets d'avenir (l. 9-10: \all{weiß, dass ... ; Wenn ... beendet habe, werde ... zurückkehren}).
\item \all{bin geflogen} (fliegen, sein) ; \all{habe beantragt} (beantragen, haben) ; \all{haben begleitet} (begleiten, haben) ; \all{hat geweint} (weinen, haben) ; \all{habe passiert} (passieren, haben) ; \all{habe ... Hoffnung gefasst} (Hoffnung fassen, haben) ; \all{beendet habe} (beenden, haben) — \textit{Note : \all{war, hatte, wollte, weiß, ist, werde ... zurückkehren} ne sont pas au Perfekt.}
\item \all{dass das Leben im Ausland am Anfang schwer ist} (verbe \all{ist} à la fin, Présent) ; \all{Wenn ich mein Studium beendet habe} (verbe \all{habe} à la fin, Perfekt : valeur de futur antérieur / action achevée avant l'action principale \all{werde ... zurückkehren}).
\item \all{das Ausland}, \all{der Pass}, \all{das Visum}, \all{die Grenze}, \all{das Heimweh}, \all{die Hoffnung}, \all{der Abflug}, \all{die Ankunft}, \all{die Passkontrolle}, \all{die Botschaft} (10 mots au total).
\end{enumerate}
\end{corrige}

\courssub{Analyse linguistique du modèle : structures cibles}

Le tableau ci-dessous récapitule les formes grammaticales obligatoires pour votre production. Utilisez-le comme aide-mémoire pendant la rédaction.

\begin{center}
\begin{tabularx}{\linewidth}{|X|X|X|}
\hline
\rowcolor{popblueL} \textbf{Contrainte} & \textbf{Forme attendue} & \textbf{Exemple du modèle} \\
\hline
Verbes au Perfekt (min. 5) & Auxiliaire (\all{haben}/\all{sein}) en 2e position + Participe II en fin de phrase & \all{Ich \textbf{habe} mein Visum \textbf{beantragt}.} \\
Subordonnée en \all{dass} & \all{dass} + sujet + ... + verbe conjugué \textbf{à la fin} & \all{Ich weiß, \textbf{dass das Leben schwer \textbf{ist}}.} \\
Subordonnée en \all{wenn} & \all{wenn} + sujet + ... + verbe conjugué \textbf{à la fin} & \all{\textbf{Wenn} ich mein Studium \textbf{beendet habe}, ...} \\
Vocabulaire thématique (min. 6) & Noms avec majuscule, article correct, pluriel connu & \all{der Pass, die Grenze, das Heimweh, die Hoffnung, das Ausland, das Visum} \\
Plan en 4 parties & Paragraphes ou phrases distinctes marquant la progression & 1. Abflug/Gründe 2. Vorbereitung 3. Gefühle 4. Zukunft \\
\hline
\end{tabularx}
\end{center}

\begin{propriete}[Rappel : Formation du Participe II (Partizip II)]
\begin{itemize}
\item \textbf{Verbes faibles (réguliers) :} \all{ge-} + radical + \all{-t} $\rightarrow$ \all{ge-macht, ge-lernt, be-antrag-t} (pas de \all{ge-} si préfixe inséparable : \all{be-, er-, ver-, ent-, zer-, ge-, miss-}).
\item \textbf{Verbes forts (irréguliers) :} \all{ge-} + radical modifié + \all{-en} $\rightarrow$ \all{ge-flog-en, ge-les-en, ge-seh-en}.
\item \textbf{Verbes en \all{-ieren} :} radical + \all{-t} (sans \all{ge-}) $\rightarrow$ \all{studier-t, telefonier-t, kopier-t}.
\item \textbf{Verbes à particule séparable :} \all{ge-} s'insère \textbf{entre} la particule et le radical $\rightarrow$ \all{an-ge-komm-en, zu-rück-ge-kehr-t, ein-ge-leb-t}.
\item \textbf{Auxiliaire \all{sein} :} verbes de déplacement (\all{gehen, fahren, fliegen, kommen, ankommen, zurückkehren, reisen, einreisen}) + changement d'état (\all{werden, sterben, aufwachen, einschlafen}) + \all{sein, bleiben, passieren} (au sens « se produire »).
\end{itemize}
\end{propriete}

\courssub{Méthodologie de la rédaction}

\begin{methode}[Rédiger un récit personnel cohérent et correct]
\begin{enumerate}
\item \textbf{Faire un plan détaillé} en français selon les 4 parties imposées : notez les mots-clés allemands (vocabulaire, verbes) pour chaque partie.
\item \textbf{Rédiger des phrases simples et correctes} plutôt que des phrases longues et fausses. Une phrase = une idée. Utilisez les connecteurs logiques (\all{zuletzt, zuerst, dann, aber, denn, weil}).
\item \textbf{Vérifier chaque verbe} systématiquement : 
\begin{itemize}
    \item Verbe principal en \textbf{2e position} (phrase principale) ?
    \item Au Perfekt : auxiliaire \textbf{conjugué} en 2e position + \textbf{participe II à la fin} ?
    \item Auxiliaire correct (\all{haben} ou \all{sein}) ?
    \item Participe II correct (\all{ge-}, radical, \all{-t/-en}, placement particule) ?
\end{itemize}
\item \textbf{Relire en traquant les pièges} : 
\begin{itemize}
    \item Virgule \textbf{obligatoire} avant \all{dass} et \all{wenn}.
    \item Verbe \textbf{à la fin} après \all{dass/wenn}.
    \item \textbf{Majuscule} à \textbf{tous les noms} (mots du lexique !).
    \item \textbf{Articles} corrects (genre, cas : ici surtout Nominatif/Akkusatif).
\end{itemize}
\end{enumerate}
\end{methode}

\courssub{Exemples de phrases clés pour votre rédaction}

Ces phrases « toutes faites » (chunks) vous serviront de briques pour construire votre texte. Adaptez les éléments entre crochets.

\begin{exemplebox}[Phrases-modèles pour « Mein Weg ins Ausland »]
\textbf{1. Le départ (Grund \& Zeit)}
\all{Letztes Jahr / Vor zwei Jahren bin ich nach [Stadt/Land] gereist.} « L'année dernière / Il y a deux ans, j'ai voyagé vers [Ville/Pays]. »
\all{Ich bin ausgewandert, weil ich [Grund: studieren / arbeiten / Familie] wollte.} « J'ai émigré car je voulais [étudier/travailler/famille]. »

\textbf{2. Les préparatifs (Pass, Visum, Abschied)}
\all{Zuerst habe ich meinen Pass und mein Visum beantragt.} « D'abord, j'ai demandé mon passeport et mon visa. »
\all{Meine Familie hat mich zum Flughafen [begleitet / gebracht].} « Ma famille m'a accompagné(e) à l'aéroport. »
\all{Der Abflug war [traurig / aufregend / schwer].} « Le départ était [triste / excitant / dur]. »

\textbf{3. Les sentiments (Heimweh, Hoffnung, Grenze, Ankunft)}
\all{An der Grenze habe ich die Passkontrolle passiert.} « À la frontière, j'ai passé le contrôle des passeports. »
\all{Nach der Ankunft in [Stadt] hatte ich großes Heimweh.} « Après mon arrivée à [Ville], j'avais un grand mal du pays. » (\textit{Präteritum usuel pour état})
\all{Aber ich habe große Hoffnung gefasst / habe ich mich gut eingelebt.} « Mais j'ai conçu de l'espoir / je me suis bien installé. » (\textit{Perfekt})

\textbf{4. Les projets (dass, wenn, Futur)}
\all{Ich weiß, \textbf{dass} das Leben im Ausland [nicht einfach / eine Chance] ist.} « Je sais que la vie à l'étranger [n'est pas simple / est une chance]. »
\all{\textbf{Wenn} ich mein [Studium / Ausbildung] beendet habe, werde ich nach Kamerun zurückkehren.} « Quand j'aurai terminé mes [études / formation], je retournerai au Cameroun. »
\all{Mein Ziel ist, [Beruf] zu werden und meiner Familie zu helfen.} « Mon but est de devenir [métier] et d'aider ma famille. »
\end{exemplebox}

\courssub{Grille d'auto-évaluation et points de vigilance}

Avant de remettre votre copie, cochez chaque case de cette liste de contrôle. C'est la garantie d'une note maximale sur la forme linguistique.

\begin{attention}[Check-list finale : « Habe ich alles richtig gemacht? »]
\begin{itemize}
\item[$\square$] \textbf{Position du verbe (phrase principale) :} Le verbe conjugué (auxiliaire au Perfekt) est-il bien en \textbf{2e position} ? (\all{Ich \textbf{habe} ...}, \all{Gestern \textbf{bin} ich ...})
\item[$\square$] \textbf{Participe II à la fin :} Le participe passé est-il bien à la \textbf{dernière place} de la phrase ? (\all{... Visum beantrag\textbf{t}.}, \all{... nach München ge\textbf{flogen}.})
\item[$\square$] \textbf{Formation du Participe II :} 
\begin{itemize}
    \item \all{ge-} présent pour les verbes faibles/forts séparables ? (\all{ge-mach-t, an-ge-komm-en, ein-ge-leb-t})
    \item \textbf{Pas de \all{ge-}} pour verbes en \all{-ieren} (\all{studier-t}) et verbes à préfixe inséparable (\all{be-antrag-t, ver-gess-en, be-endet}) ?
    \item Radical correct pour verbes forts ? (\all{ge-flog-en} pas \all{ge-flieg-t} ; \all{ge-seh-en} pas \all{ge-seh-t})
\end{itemize}
\item[$\square$] \textbf{Auxiliaire correct :} \all{sein} pour \all{fliegen, reisen, fahren, gehen, kommen, ankommen, zurückkehren, einreisen, aufstehen} ? \all{haben} pour les autres (\all{haben, machen, beantragen, lernen, passieren, weinen, begleiten}) ?
\item[$\square$] \textbf{Subordonnées (\all{dass} / \all{wenn}) :} 
\begin{itemize}
    \item Virgule \textbf{avant} la conjonction ? (\all{..., \textbf{dass} ...}, \all{..., \textbf{wenn} ...})
    \item Verbe conjugué \textbf{à la toute fin} de la subordonnée ? (\all{..., dass ich es \textbf{weiß}.}, \all{..., wenn ich ankomme \textbf{.}})
\end{itemize}
\item[$\square$] \textbf{Orthographe des noms :} \textbf{Majuscule} à \textbf{tous} les noms (\all{der Pass, die Grenze, das Heimweh, die Hoffnung, das Ausland, das Visum, der Abflug, die Ankunft, die Botschaft}) ?
\item[$\square$] \textbf{Articles et cas :} Article correct (genre) ? Cas correct (souvent Accusatif après \all{haben, beantragen, besuchen, passieren} : \all{mein\textbf{en} Pass, mein\textbf{e} Hoffnung, die Passkontrolle}) ?
\item[$\square$] \textbf{Contenu :} 4 parties du plan respectées ? 5 Perfekt ? 1 \all{dass} ? 1 \all{wenn} ? 6 mots lexique ? 8-10 lignes ?
\end{itemize}
\end{attention}

\courssub{Exercice résolu : correction d'un texte d'élève}

Voici un texte anonyme rédigé par un élève de Première. Il contient des erreurs types (francophonismes, oublis de règles). Appliquez la check-list ci-dessus pour les identifier, puis comparez avec la correction commentée.

\begin{exoresolu}{Correction d'un texte d'élève : « Mein Weg nach Deutschland »}
\textbf{Texte de l'élève (avec erreurs) :}
„Letztes Jahr ich bin nach Deutschland gereist. Ich wollte studieren Medizin. Ich habe mein Pass und Visum beantragt. Meine Mutter hat geweint beim Abflug. An der Grenze ich habe die Passkontrolle passiert. In München ich hatte Heimweh. Aber ich habe Hoffnung das ich schaffe es. Wenn ich fertig studiert habe, ich komme zurück nach Kamerun.“

\tcblower

\textbf{Analyse et correction détaillée :}

\begin{enumerate}
\item \all{Letztes Jahr \textbf{ich bin} ...} $\rightarrow$ \textbf{Erreur :} Sujet \all{ich} en 1re position, verbe en 3e. \textbf{Règle :} Verbe en 2e position. \textbf{Correction :} \all{Letztes Jahr \textbf{bin ich} nach Deutschland gereist.}
\item \all{studieren Medizin} $\rightarrow$ \textbf{Erreur :} Ordre des mots (objet avant infinitif). \textbf{Correction :} \all{Medizin studieren} (ou \all{Informatik studieren}).
\item \all{Ich habe \textbf{mein} Pass ...} $\rightarrow$ \textbf{Erreur :} \all{Pass} masculin, accusatif (COD de \all{beantragen}) $\rightarrow$ \all{mein\textbf{en} Pass}. \textbf{Correction :} \all{Ich habe \textbf{meinen} Pass und \textbf{mein} Visum beantragt.} (\all{Visum} neutre, accusatif = \all{mein Visum}).
\item \all{Meine Mutter hat geweint beim Abflug.} $\rightarrow$ \textbf{Correct} (mais style lourd). Mieux : \all{Beim Abflug hat meine Mutter geweint.} (Topicalisation du complément de temps).
\item \all{An der Grenze \textbf{ich habe} ...} $\rightarrow$ \textbf{Erreur :} Complément de lieu en 1re position $\rightarrow$ inversion sujet/verbe. \textbf{Correction :} \all{An der Grenze \textbf{habe ich} die Passkontrolle passiert.}
\item \all{In München \textbf{ich hatte} ...} $\rightarrow$ \textbf{Erreur :} Inversion manquante. \textbf{Correction :} \all{In München \textbf{hatte ich} großes Heimweh.} (Ajout adjectif \all{großes} pour le style).
\item \all{Hoffnung \textbf{das} ich schaffe es.} $\rightarrow$ \textbf{Erreurs multiples :} 
    \begin{itemize}
        \item Conjonction : \all{dass} (double s) pas \all{das} (article/démonstratif).
        \item Virgule obligatoire avant \all{dass}.
        \item Structure \all{schaffe es} : ordre standard en subordonnée = sujet (\all{ich}) + objet (\all{es}) + verbe (\all{schaffe}). Mais \all{schaffen} (réussir) + \all{es} est familier. Mieux : \all{dass ich es schaffe} (Présent) ou \all{dass ich es schaffen werde} (Futur).
    \end{itemize}
    \textbf{Correction :} \all{Aber ich habe große Hoffnung, \textbf{dass} ich \textbf{es schaffen werde}.}
\item \all{Wenn ich fertig \textbf{studiert habe}, \textbf{ich komme} ...} $\rightarrow$ \textbf{Erreurs :} 
    \begin{itemize}
        \item \all{fertig studiert habe} : structure possible mais \all{mein Studium beendet habe} est plus standard.
        \item \textbf{Verbe en 2e position dans la subordonnée !} \all{ich komme} $\rightarrow$ \textbf{Interdit.} Verbe à la fin.
        \item \all{zurückkommen} verbe séparable $\rightarrow$ particule \all{zurück} avant le verbe final : \all{... zurückkomme} ou futur \all{... werde ich zurückkehren}.
    \end{itemize}
    \textbf{Correction :} \all{\textbf{Wenn} ich mein Studium \textbf{beendet habe}, \textbf{werde ich} nach Kamerun \textbf{zurückkehren}.}
\end{enumerate}

\textbf{Version corrigée finale :}
\all{Letztes Jahr \textbf{bin ich} nach Deutschland \textbf{gereist}. Ich wollte \textbf{Medizin studieren}. Ich habe \textbf{meinen} Pass und mein Visum beantragt. \textbf{Beim Abflug} hat meine Mutter geweint. An der Grenze \textbf{habe ich} die Passkontrolle passiert. In München \textbf{hatte ich} großes Heimweh. Aber ich habe große Hoffnung, \textbf{dass} ich \textbf{es schaffen werde}. \textbf{Wenn} ich mein Studium \textbf{beendet habe}, \textbf{werde ich} nach Kamerun \textbf{zurückkehren}.}
\end{exoresolu}

\courssub{Production écrite personnelle : « Mein Weg ins Ausland »}

C'est à vous maintenant. Rédigez votre propre texte en respectant scrupuleusement les consignes.

\begin{exercice}{Votre tour : « Mein Weg ins Ausland »}
\textbf{Consignes :}
\begin{enumerate}
\item Rédigez un texte de \textbf{8 à 10 lignes} (environ 80-100 mots) intitulé \all{Mein Weg ins Ausland}.
\item Le récit peut être \textbf{réel} (votre expérience ou celle d'un proche) ou \textbf{imaginaire} (projet futur, scénario fictif).
\item \textbf{Plan imposé (4 parties distinctes) :}
\begin{enumerate}
\item \textbf{Le départ :} Quand ? D'où ? Vers où ? Pourquoi ? (Études, travail, famille, fuite ?)
\item \textbf{Les préparatifs :} Documents (\all{Pass, Visum}), formalités (\all{Botschaft}), au revoir (\all{Abflug, Flughafen}).
\item \textbf{Les sentiments :} Arrivée (\all{Ankunft}), frontière (\all{Grenze, Passkontrolle}), émotions (\all{Heimweh, Hoffnung, Angst, Freude}).
\item \textbf{Les projets d'avenir :} Études, métier, retour (\all{zurückkehren}), aide à la famille.
\end{enumerate}
\item \textbf{Contraintes linguistiques OBLIGATOIRES (points de notation) :}
\begin{itemize}
    \item Au moins \textbf{5 verbes au Perfekt} (auxiliaire + participe II bien formés).
    \item Exactement \textbf{1 subordonnée en \all{dass}} (avec virgule, verbe à la fin).
    \item Exactement \textbf{1 subordonnée en \all{wenn}} (avec virgule, verbe à la fin, temps logique).
    \item Au moins \textbf{6 mots du lexique thématique} (voir liste Section 3 : \all{auswandern, das Ausland, das Heimweh, die Hoffnung, der Pass, die Grenze, das Visum, die Botschaft, der Abflug, die Ankunft, die Passkontrolle, das Studium, die Ausbildung, der Beruf, die Familie}...).
    \item Respect de la \textbf{grille d'auto-évaluation} (voir boîte \textbf{Attention} ci-dessus).
\end{itemize}
\item Écrivez \textbf{lisiblement} sur votre copie, en sautant une ligne. Relisez-vous avec la check-list avant de rendre.
\end{enumerate}
\end{exercice}

\courssub{Production orale : présentation et interaction}

L'écrit prépare l'oral. Entraînez-vous à présenter votre récit sans le lire mot à mot.

\begin{experience}[Activité orale : « Mein Weg – Präsentation und Interview »]
\textbf{Déroulement (par binômes, 10 min total) :}
\begin{enumerate}
\item \textbf{Préparation (2 min) :} Relisez votre texte. Repérez les mots-clés (mots du lexique, verbes au Perfekt, conjonctions). Préparez 2 questions à poser à votre camarade (voir exemples ci-dessous).
\item \textbf{Présentation A $\rightarrow$ B (2 min) :} L'élève A raconte son histoire d'émigration \textbf{à l'oral}, en s'appuyant sur ses notes/mots-clés, \textbf{pas en lisant}. L'élève B écoute activement et prend des notes (mots-clés entendus : quand, où, pourquoi, sentiments, projets).
\item \textbf{Questions de compréhension B $\rightarrow$ A (2 min) :} L'élève B pose \textbf{2 questions} en allemand pour vérifier sa compréhension ou approfondir.
    \textit{Exemples de questions types :}
    \begin{itemize}
        \item \all{Wann bist du ausgewandert?}
        \item \all{Warum bist du nach [Land] gegangen?}
        \item \all{Was hast du vor der Reise vorbereitet? (Pass, Visum?)}
        \item \all{Wie war der Abflug? Warst du traurig?}
        \item \all{Hattest du Heimweh nach der Ankunft?}
        \item \all{Was sind deine Pläne für die Zukunft?}
        \item \all{Wann willst du nach Kamerun zurückkehren?}
    \end{itemize}
\item \textbf{Inversion des rôles (4 min) :} L'élève B présente, l'élève A écoute et pose 2 questions.
\item \textbf{Retour classe (optionnel) :} Quelques binômes présentent un résumé de l'histoire de leur partenaire à la classe (\all{Mein Partner ist vor zwei Jahren nach Berlin geflogen...}).
\end{enumerate}
\textbf{Critères d'évaluation orale :} Fluidité / Prononciation / Respect structures (Perfekt, \all{dass/wenn}) / Vocabulaire thématique / Interaction (questions/réponses pertinentes).
\end{experience}

\courssub{Frise chronologique du récit}

Cette frise visualise la structure narrative type d'une histoire d'émigration. Utilisez-la pour structurer votre plan oral ou écrit.

\begin{popfigure}
\begin{tikzpicture}[
    node distance=1.2cm and 1.5cm,
    etape/.style={rectangle, rounded corners, draw=popblue, thick, fill=popblueL, text width=2.8cm, align=center, minimum height=1.8cm, font=\small},
    fleche/.style={-Stealth, thick, popdark}
]
% Nœuds
\node[etape, align=center] (start){\textbf{1. Abflug in Douala}\\\all{letztes Jahr}\\Verabschiedung\\Familie \& Freunde};
\node[etape, right=of start, align=center] (grenz){\textbf{2. Grenze / Passkontrolle}\\\all{der Pass, das Visum}\\Kontrolle\\Dokumente prüfen};
\node[etape, right=of grenz, align=center] (ank){\textbf{3. Ankunft in München}\\\all{die Ankunft}\\Neue Umgebung\\Orientierung};
\node[etape, below=of grenz, align=center] (heim){\textbf{4. Heimweh}\\\all{das Heimweh}\\Sehnsucht nach\\Essen, Familie, Klima};
\node[etape, right=of heim, align=center] (zuk){\textbf{5. Hoffnung \& Zukunft}\\\all{die Hoffnung}\\Studium / Beruf\\\all{zurückkehren} (wenn...)};

% Flèches principales
\draw[fleche] (start) -- (grenz) node[midway, above, font=\tiny, popdark] {Reise};
\draw[fleche] (grenz) -- (ank) node[midway, above, font=\tiny, popdark] {Einreise};
\draw[fleche] (ank) -- (heim) node[midway, right, font=\tiny, popdark] {Gefühle};
\draw[fleche] (heim) -- (zuk) node[midway, above, font=\tiny, popdark] {Perspektive};
\draw[fleche, poporange] (ank) -- (zuk) node[midway, above, font=\tiny, poporange] {Motivation};

% Légende intégrée
\node[below=0.5cm of heim, font=\footnotesize\textit, popmuted, text width=10cm, align=center] {Structure narrative type : Départ $\rightarrow$ Formalités $\rightarrow$ Arrivée $\rightarrow$ Crise (Heimweh) $\rightarrow$ Projet (Hoffnung)};
\end{tikzpicture}
\legende{Frise narrative : les 5 étapes clés d'un récit d'émigration (vocabulaire cible en allemand).}
\end{popfigure}

\courssub{Le savais-tu ?}

\begin{savaistu}[Le Goethe-Institut à Yaoundé : votre passerelle vers l'Allemagne]
Pour étudier, travailler ou faire un stage en Allemagne, il faut généralement prouver sa maîtrise de l'allemand par un certificat officiel reconnu internationalement. Les deux examens de référence sont :
\begin{itemize}
\item \textbf{Le Goethe-Zertifikat} (niveaux A1 à C2) : délivré par le Goethe-Institut lui-même. Le \textbf{Goethe-Institut Kamerun} est situé à \textbf{Yaoundé} (quartier Bastos). Il organise des sessions d'examen plusieurs fois par an (souvent en mars, juin, septembre, décembre). Il propose aussi des cours de préparation.
\item \textbf{Le TestDaF} (Test Deutsch als Fremdsprache) : examen standardisé pour l'accès à l'université allemande (niveau B2/C1). Il se passe aussi dans des centres agréés, dont le Goethe-Institut de Yaoundé.
\end{itemize}
\textbf{Conseil pratique :} Si vous visez une bourse DAAD ou une inscription en université allemande (comme Kevin dans notre texte), visez le niveau \textbf{B2} (Goethe-Zertifikat B2 ou TestDaF Niveau 4) \textbf{avant} de déposer votre dossier. Les inscriptions aux examens se font en ligne sur le site du Goethe-Institut Kamerun (\texttt{www.goethe.de/kamerun}). Pensez à vous inscrire tôt, les places partent vite !
\end{savaistu}

\textbf{Devoir maison :} Finalisez votre texte « \all{Mein Weg ins Ausland} » à la maison en utilisant la grille d'auto-évaluation. Préparez-vous à le présenter oralement lors de la prochaine séance. Révisez les tableaux de conjugaison du Perfekt (verbes forts du chapitre) et la liste de vocabulaire thématique pour le contrôle continu.

\newpage

\coursec{Activité d'intégration}

\courssub{Objectifs de l'activité}

À l'issue de cette activité d'intégration, l'élève sera capable de :
\begin{itemize}
\item mobiliser le \cle{vocabulaire de l'émigration} (\all{auswandern}, \all{das Ausland}, \all{das Heimweh}, \all{die Hoffnung}, \all{der Pass}, \all{die Grenze}) dans une situation de communication réelle ;
\item utiliser correctement le \cle{parfait} (\all{ich bin ausgewandert}, \all{ich habe meinen Pass gezeigt}) pour raconter des événements passés ;
\item employer au moins une \cle{subordonnée en dass} et une \cle{subordonnée en wenn} à l'oral et à l'écrit ;
\item conduire un \cle{interview} en allemand : poser des questions, écouter, réagir ;
\item rédiger un court \cle{article de presse} rapportant les propos d'une personne (\all{Er sagt, dass ...}).
\end{itemize}

\courssub{Situation}

La radio locale de ta ville organise une émission spéciale intitulée \all{„Kameruner in der Welt“} (« Les Camerounais dans le monde »). Le journaliste reçoit un jeune auditeur qui a émigré en Allemagne il y a deux ans et qui est de retour au pays pour les vacances. Il raconte son départ, son arrivée, ses difficultés et ses espoirs. L'émission sera ensuite publiée sous forme d'article sur le site de la radio.

Voici le début de l'interview, tel qu'il a été diffusé :

\begin{textebox}[Radio Yaoundé — Ausschnitt aus der Sendung „Kameruner in der Welt“]
\textbf{Journalist:} Guten Tag, Herr Mbarga, und willkommen in unserem Studio. Sie sind vor zwei Jahren nach Deutschland ausgewandert. Warum haben Sie Kamerun verlassen?

\textbf{Herr Mbarga:} Guten Tag. Ich bin ausgewandert, weil ich in Deutschland Informatik studieren wollte. Ich hatte große Hoffnung, dass ich dort eine gute Ausbildung finde.

\textbf{Journalist:} Wie war die Reise? Was haben Sie an der Grenze gezeigt?

\textbf{Herr Mbarga:} Die Reise war lang. An der Grenze habe ich meinen Pass und mein Visum gezeigt. Ich war sehr nervös, aber alles war in Ordnung.

\textbf{Journalist:} Hatten Sie am Anfang Heimweh?

\textbf{Herr Mbarga:} Ja, natürlich! Wenn ich am Abend allein in meinem Zimmer war, habe ich oft an meine Familie in Yaoundé gedacht. Aber ich habe jeden Tag mit meiner Mutter telefoniert.
\end{textebox}

\begin{definition}[Glossaire]
\begin{itemize}
\item \all{die Sendung, -en} : l'émission (de radio)
\item \all{verlassen (hat verlassen)} : quitter
\item \all{die Ausbildung, -en} : la formation
\item \all{das Visum, die Visa} : le visa
\item \all{nervös} : nerveux
\item \all{telefonieren (hat telefoniert)} : téléphoner
\end{itemize}
\end{definition}

\courssub{Consignes}

Travail par \cle{groupes de deux} : un élève joue le rôle de l'émigrant, l'autre celui du journaliste.

\begin{enumerate}
\item \textbf{Lecture.} Lis le début de l'interview ci-dessus à voix haute avec ton camarade. Souligne tous les verbes au parfait et encadre les deux subordonnées (une en \all{dass}, une en \all{wenn}).

\item \textbf{Préparation du journaliste.} Rédige \cle{5 questions} pour ton interview. Utilise les modèles : \all{Warum bist du ausgewandert?} / \all{Was hast du an der Grenze gezeigt?} / \all{Hattest du Heimweh?} / \all{Was hast du in Deutschland gemacht?} / \all{Was willst du später machen?}

\item \textbf{Préparation de l'émigrant.} Prépare tes réponses : raconte ton histoire au \cle{parfait} (minimum 5 phrases). Ta fiche de personnage : tu t'appelles comme tu veux, tu viens de Douala ou de Yaoundé, tu es parti(e) en Allemagne (ou en Autriche, en Suisse) pour étudier ou travailler. Tu dois placer \cle{au moins une phrase en dass} et \cle{au moins une phrase en wenn}.

\item \textbf{Jeu de rôle.} Jouez l'interview devant la classe (3 à 4 minutes). Le journaliste commence par : \all{Guten Tag und willkommen in unserem Studio!} et termine par : \all{Vielen Dank für das Interview!}

\item \textbf{Rédaction.} Chaque journaliste rédige ensuite un court article de \cle{8 lignes} intitulé \all{„Ein Kameruner in Deutschland“}. Il rapporte les propos de l'émigrant avec les structures \all{Er/Sie sagt, dass ...} et \all{Wenn er/sie ..., will er/sie ...}.

\item \textbf{Relecture.} Échange ton article avec un autre groupe : vérifie le parfait (auxiliaire \all{haben} ou \all{sein}), la place du verbe dans les subordonnées et la majuscule aux noms.
\end{enumerate}

\courssub{Ressources à mobiliser}

\begin{propriete}[Rappel : das Perfekt]
Le parfait se forme avec l'auxiliaire \all{haben} ou \all{sein} (2\up{e} position) + le participe II (à la fin). On utilise \all{sein} avec les verbes de déplacement et de changement d'état : \all{ich \textbf{bin} ausgewandert}, \all{ich \textbf{bin} nach Deutschland geflogen}, mais \all{ich \textbf{habe} meinen Pass gezeigt}.
\end{propriete}

\begin{propriete}[Rappel : dass et wenn]
Dans une subordonnée introduite par \all{dass} ou \all{wenn}, le verbe conjugué va \textbf{à la fin} : \all{Ich hatte Hoffnung, dass ich eine Arbeit \textbf{finde}.} / \all{Wenn ich Heimweh \textbf{habe}, telefoniere ich mit meiner Familie.} Attention : quand la subordonnée en \all{wenn} commence la phrase, le verbe de la principale vient juste après la virgule.
\end{propriete}

\begin{aretenir}[Lexique utile pour l'interview]
\begin{itemize}
\item \all{auswandern (ist ausgewandert)} : émigrer ; \all{der Auswanderer, -} : l'émigrant
\item \all{das Ausland} : l'étranger ; \all{im Ausland leben} : vivre à l'étranger
\item \all{das Heimweh} : le mal du pays ; \all{Heimweh haben}
\item \all{die Hoffnung, -en} : l'espoir ; \all{Hoffnung haben, dass ...}
\item \all{der Pass, „e} : le passeport ; \all{die Grenze, -n} : la frontière
\item \all{zeigen (hat gezeigt)} : montrer ; \all{den Pass zeigen}
\end{itemize}
\end{aretenir}

\begin{attention}[Pièges fréquents]
Ne dis pas \all{*ich habe ausgewandert} : le verbe \all{auswandern} exprime un déplacement, il exige l'auxiliaire \all{sein} : \all{ich \textbf{bin} ausgewandert}. Et dans la subordonnée, le verbe ne reste pas en deuxième position comme en français : \all{Er sagt, dass er Heimweh \textbf{hat}} (et non \all{*dass er hat Heimweh}).
\end{attention}

\courssub{Proposition de production et corrigé commenté}

\begin{exoresolu}[Article modèle : „Ein Kameruner in Deutschland“]
Rédige un article de 8 lignes rapportant l'interview d'un jeune Camerounais émigré en Allemagne, avec le parfait, une subordonnée en \all{dass} et une en \all{wenn}.
\tcblower
\textbf{Production modèle :}

\all{Kevin Mbarga kommt aus Yaoundé. Er ist vor zwei Jahren nach Deutschland ausgewandert, weil er dort Informatik studieren wollte. An der Grenze hat er seinen Pass und sein Visum gezeigt. Am Anfang hatte er oft Heimweh. Er sagt, dass er jeden Tag mit seiner Familie telefoniert hat. In Deutschland hat er viele neue Freunde gefunden und gut Deutsch gelernt. Wenn er sein Studium beendet hat, will er nach Kamerun zurückkehren. Er hat große Hoffnung, dass er in seinem Land eine gute Arbeit findet. Wir wünschen ihm viel Erfolg!}

« Kevin Mbarga vient de Yaoundé. Il a émigré en Allemagne il y a deux ans parce qu'il voulait y étudier l'informatique. À la frontière, il a montré son passeport et son visa. Au début, il avait souvent le mal du pays. Il dit qu'il a téléphoné chaque jour à sa famille. En Allemagne, il a trouvé beaucoup de nouveaux amis et bien appris l'allemand. Quand il aura terminé ses études, il veut retourner au Cameroun. Il a le grand espoir de trouver un bon travail dans son pays. Nous lui souhaitons beaucoup de succès ! »

\textbf{Commentaire des choix de langue :}
\begin{itemize}
\item \all{ist ... ausgewandert} et \all{hat ... gezeigt / telefoniert / gefunden / gelernt} : parfait correct, avec \all{sein} pour le déplacement et \all{haben} pour les autres verbes ; le participe II est bien placé à la fin.
\item \all{Er sagt, dass er jeden Tag mit seiner Familie telefoniert hat.} : subordonnée en \all{dass}, verbe conjugué (\all{hat}) renvoyé à la toute fin, après le participe.
\item \all{Wenn er sein Studium beendet hat, will er nach Kamerun zurückkehren.} : la subordonnée en \all{wenn} ouvre la phrase ; le verbe de la principale (\all{will}) vient immédiatement après la virgule, avant le sujet.
\item Lexique du thème bien mobilisé : \all{auswandern}, \all{der Pass}, \all{die Grenze}, \all{das Heimweh}, \all{die Hoffnung}.
\item Tous les noms portent la majuscule : \all{Pass}, \all{Visum}, \all{Heimweh}, \all{Studium}, \all{Arbeit}, \all{Erfolg}.
\end{itemize}
\end{exoresolu}

\courssub{Bilan de l'activité}

Coche ce que tu sais maintenant faire :
\begin{itemize}
\item[$\square$] Je peux raconter une émigration au parfait avec le bon auxiliaire (\all{haben} / \all{sein}).
\item[$\square$] Je peux construire une subordonnée en \all{dass} et une en \all{wenn} avec le verbe à la fin.
\item[$\square$] Je connais et j'emploie le vocabulaire de l'émigration avec les bons articles.
\item[$\square$] Je peux poser des questions et rapporter les propos de quelqu'un dans un court article.
\end{itemize}

\begin{savaistu}[Le savais-tu ?]
L'Allemagne compte aujourd'hui une importante communauté camerounaise : des étudiants, des médecins, des ingénieurs et des sportifs. Beaucoup reviennent ensuite au pays avec de nouvelles compétences : on parle alors de \all{„Rückkehr“} (le retour) et de \all{„Brain Gain“} (le gain de cerveaux), par opposition au \all{„Brain Drain“} (la fuite des cerveaux).
\end{savaistu}

\newpage

\coursec{Méthodes et exercices résolus}

\courssub{Méthode 1 : Comprendre un texte sur l'émigration}

\begin{methode}[Lire et comprendre un texte : la démarche en 5 étapes]
\begin{enumerate}
\item \textbf{Lire le titre et la source.} Le titre \all{Kevins neues Leben in Deutschland} annonce le thème (l'émigration) et le point de vue (un jeune Camerounais en Allemagne). Pose-toi la question : de quoi va parler le texte ?
\item \textbf{Première lecture globale.} Lis le texte une fois sans dictionnaire. Repère : \emph{qui} parle, \emph{de quoi} il s'agit, \emph{où} et \emph{quand} se situe l'action. Souligne les mots que tu connais déjà (\all{Familie, Arbeit, Sprache, Heimat}).
\item \textbf{Deuxième lecture détaillée.} Souligne les mots-clés du thème : \all{auswandern, das Ausland, die Grenze, der Pass, das Heimweh, die Hoffnung}. Devine le sens des mots inconnus grâce au contexte et aux familles de mots (\all{wandern} $\rightarrow$ \all{auswandern} $\rightarrow$ \all{der Auswanderer}).
\item \textbf{Repérer les connecteurs.} Les subordonnées en \all{dass} (opinion, fait) et \all{wenn} (condition, temps) structurent les idées : repère-les pour comprendre les relations logiques.
\item \textbf{Répondre aux questions.} Réponds en phrases complètes, en allemand, en reprenant les mots de la question. Cite le texte pour justifier (\all{Im Text steht: „...“}).
\end{enumerate}
\textbf{Structures à retenir :} \all{Der Text handelt von...} (le texte parle de...) ; \all{Der Autor sagt, dass...} ; \all{Kevin ist ausgewandert, weil...}
\end{methode}

\begin{exoresolu}[Compréhension de texte — type Bac]
\textbf{Énoncé.} Lis le texte \all{Kevins neues Leben in Deutschland} et réponds en allemand par des phrases complètes :
\begin{enumerate}
\item \all{Warum ist Kevin nach Deutschland ausgewandert?}
\item \all{Was hat er vor der Abreise gemacht?}
\item \all{Was vermisst er in Deutschland?}
\end{enumerate}
\tcblower
\textbf{Raisonnement.} La question 1 porte sur la \emph{cause} : je cherche dans le texte un connecteur de cause (\all{weil}, \all{um...zu}) ou une subordonnée en \all{dass} exprimant l'espoir. La question 2 porte sur des actions \emph{avant le départ} : je cherche des verbes au Perfekt (\all{hat gemacht, hat beantragt}). La question 3 porte sur le manque : je cherche le mot \all{Heimweh} ou le verbe \all{vermissen}.

\textbf{Réponses rédigées.}
\begin{enumerate}
\item \all{Kevin ist nach Deutschland ausgewandert, weil er dort studieren und eine bessere Zukunft finden wollte.} « Kevin a émigré en Allemagne parce qu'il voulait y étudier et trouver un avenir meilleur. » — Subordonnée de cause avec \all{weil} : le verbe conjugué \all{wollte} va à la fin.
\item \all{Vor der Abreise hat er einen Pass beantragt und ein Visum bekommen. Er hat auch Deutsch gelernt.} « Avant le départ, il a demandé un passeport et obtenu un visa. Il a aussi appris l'allemand. » — Perfekt : auxiliaire \all{haben} en 2\textsuperscript{e} position + participe II à la fin (\all{beantragt, bekommen, gelernt}).
\item \all{In Deutschland vermisst er seine Familie, seine Freunde und das Essen aus Kamerun. Er hat oft Heimweh.} « En Allemagne, sa famille, ses amis et la cuisine du Cameroun lui manquent. Il a souvent le mal du pays. » — \all{vermissen} + accusatif ; \all{Heimweh} est neutre : \all{das Heimweh}.
\end{enumerate}
\textbf{Conclusion méthodologique.} Une bonne réponse reprend les mots de la question, cite une information précise du texte et respecte la place du verbe (2\textsuperscript{e} position en principale, dernière position en subordonnée).
\end{exoresolu}

\courssub{Méthode 2 : Employer le vocabulaire de l'émigration}

\begin{methode}[Mémoriser et utiliser le lexique thématique]
\begin{enumerate}
\item \textbf{Apprends chaque nom avec son article et son pluriel.} En allemand, le genre est imprévisible : \all{der Pass, die Pässe} ; \all{die Grenze, -n} ; \all{das Ausland} (sans pluriel) ; \all{das Heimweh} (sans pluriel) ; \all{die Hoffnung, -en}.
\item \textbf{Apprends les familles de mots.} \all{wandern} (marcher, migrer) $\rightarrow$ \all{auswandern} (émigrer) $\rightarrow$ \all{einwandern} (immigrer) $\rightarrow$ \all{der Auswanderer / die Auswanderin} $\rightarrow$ \all{die Auswanderung}.
\item \textbf{Apprends les verbes avec leur construction.} \all{auswandern nach + Dativ} (émigrer vers : \all{nach Deutschland}) ; \all{hoffen auf + Akkusativ} (espérer quelque chose) ; \all{vermissen + Akkusativ} (quelqu'un vous manque) ; \all{einen Pass beantragen} (demander un passeport).
\item \textbf{Distingue les points de vue.} \all{die Emigration} = vu du pays de départ ; \all{die Immigration} = vu du pays d'arrivée. Le même Kevin est \all{Auswanderer} pour le Cameroun et \all{Einwanderer} pour l'Allemagne.
\item \textbf{Réutilise le lexique dans tes propres phrases} pour l'expression écrite : au moins 4 mots du thème dans chaque production.
\end{enumerate}
\end{methode}

\begin{exoresolu}[Lexique — compléter un texte à trous]
\textbf{Énoncé.} Complète avec le mot qui convient : \all{das Ausland — die Grenze — der Pass — das Heimweh — die Hoffnung}.
\begin{enumerate}
\item \all{Kevin hat große \dots : er will in Deutschland studieren.}
\item \all{An der \dots\ zwischen Nigeria und Kamerun muss man seinen \dots\ zeigen.}
\item \all{Viele junge Leute träumen vom Leben im \dots .}
\item \all{Nach zwei Monaten in Berlin hat Kevin \dots : er vermisst seine Mutter.}
\end{enumerate}
\tcblower
\textbf{Raisonnement.} Je regarde l'article demandé et le sens de la phrase.
\begin{enumerate}
\item \all{große Hoffnung} : article féminin à l'accusatif (\all{große} trahit le féminin) ; sens : « il a un grand espoir ».
\item \all{An der Grenze} : \all{an} + datif féminin (\all{der}) ; sens : « à la frontière ». \all{seinen Pass} : accusatif masculin (\all{seinen}) ; « montrer son passeport ».
\item \all{im Ausland} : \all{in dem} $\rightarrow$ \all{im} ; « à l'étranger ». Attention : \all{das Ausland} n'a pas de pluriel.
\item \all{Heimweh} : sens « mal du pays » ; nom neutre, employé ici sans article (comme \all{Hunger, Durst}).
\end{enumerate}
\textbf{Texte complet :} \all{Kevin hat große Hoffnung: er will in Deutschland studieren. An der Grenze zwischen Nigeria und Kamerun muss man seinen Pass zeigen. Viele junge Leute träumen vom Leben im Ausland. Nach zwei Monaten in Berlin hat Kevin Heimweh: er vermisst seine Mutter.}

\textbf{Conclusion.} Toujours vérifier trois choses : le sens, l'article (le genre !) et le cas exigé par la préposition ou le verbe.
\end{exoresolu}

\courssub{Méthode 3 : Construire le Perfekt}

\begin{methode}[Former et employer le Perfekt (le parfait)]
Le Perfekt est le temps du passé le plus utilisé à l'oral et dans les récits personnels (lettres, dialogues, réponses aux questions).
\begin{enumerate}
\item \textbf{Structure :} sujet + \all{haben} ou \all{sein} conjugué (2\textsuperscript{e} position) + participe II (à la fin de la phrase). \all{Kevin \textbf{hat} Deutsch \textbf{gelernt}.}
\item \textbf{Choix de l'auxiliaire \all{sein} :} verbes de déplacement ou de changement d'état (\all{gehen, kommen, reisen, auswandern, fliegen, sterben, aufstehen}) + \all{sein} et \all{bleiben}. \all{Er \textbf{ist} nach Deutschland \textbf{geflogen}.}
\item \textbf{Participe II des verbes faibles :} \textbf{ge + radical + t} : \all{lernen $\rightarrow$ gelernt} ; \all{machen $\rightarrow$ gemacht} ; \all{arbeiten $\rightarrow$ gearbeitet} (e de liaison après -t/-d).
\item \textbf{Participe II des verbes forts :} \textbf{ge + (radical modifié) + en} : \all{gehen $\rightarrow$ gegangen} ; \all{fliegen $\rightarrow$ geflogen} ; \all{sprechen $\rightarrow$ gesprochen}. Apprends la liste par cœur.
\item \textbf{Verbes à particule :} séparables : \textbf{ge} entre la particule et le verbe : \all{auswandern $\rightarrow$ ausgewandert} ; \all{ankommen $\rightarrow$ angekommen}. Inséparables (\all{be-, ver-, er-}) : pas de \textbf{ge} : \all{bekommen $\rightarrow$ bekommen} ; \all{verstehen $\rightarrow$ verstanden}.
\end{enumerate}
\textbf{Réflexe :} dans une phrase au Perfekt, vérifie toujours le couple auxiliaire (2\textsuperscript{e} position) + participe (dernière position).
\end{methode}

\begin{exoresolu}[Conjugaison — mettre au Perfekt]
\textbf{Énoncé.} Mets les phrases suivantes au Perfekt.
\begin{enumerate}
\item \all{Kevin lernt zwei Jahre Deutsch.}
\item \all{Er fliegt im September nach Frankfurt.}
\item \all{Seine Familie versteht seine Entscheidung.}
\item \all{Er steht jeden Tag um sechs Uhr auf.}
\item \all{Seine Schwester kommt ihn besuchen.}
\end{enumerate}
\tcblower
\textbf{Raisonnement et solutions.}
\begin{enumerate}
\item \all{lernen} : verbe faible, auxiliaire \all{haben}. $\rightarrow$ \all{Kevin \textbf{hat} zwei Jahre Deutsch \textbf{gelernt}.}
\item \all{fliegen} : verbe fort de déplacement $\rightarrow$ auxiliaire \all{sein} ; participe \all{geflogen}. $\rightarrow$ \all{Er \textbf{ist} im September nach Frankfurt \textbf{geflogen}.}
\item \all{verstehen} : particule inséparable \all{ver-} $\rightarrow$ pas de \textbf{ge} ; auxiliaire \all{haben}. $\rightarrow$ \all{Seine Familie \textbf{hat} seine Entscheidung \textbf{verstanden}.}
\item \all{aufstehen} : particule séparable + changement d'état $\rightarrow$ \all{sein} ; participe \all{aufgestanden}. $\rightarrow$ \all{Er \textbf{ist} jeden Tag um sechs Uhr \textbf{aufgestanden}.}
\item \all{kommen} : déplacement $\rightarrow$ \all{sein} ; participe \all{gekommen}. $\rightarrow$ \all{Seine Schwester \textbf{ist} ihn besuchen \textbf{gekommen}.} (double infinitif : \all{besuchen gekommen} à la fin)
\end{enumerate}
\textbf{Conclusion.} Trois questions à se poser : (1) déplacement ou changement d'état ? $\rightarrow$ \all{sein}, sinon \all{haben} ; (2) verbe faible (\all{ge...t}) ou fort (\all{ge...en}) ? ; (3) particule séparable (\all{aufge...}) ou inséparable (pas de \all{ge}) ?
\end{exoresolu}

\courssub{Méthode 4 : Construire les subordonnées en dass et wenn}

\begin{methode}[Subordonnées en dass et wenn : la règle d'or]
\begin{enumerate}
\item \textbf{\all{dass} = « que »} : introduit un fait, une opinion, un espoir. \all{Kevin hofft, \textbf{dass} er ein Stipendium \textbf{bekommt}.} « Kevin espère qu'il obtiendra une bourse. »
\item \textbf{\all{wenn} = « quand / si »} : temps (répétition ou futur) ou condition. \all{\textbf{Wenn} er Heimweh \textbf{hat}, ruft er seine Mutter an.} « Quand il a le mal du pays, il appelle sa mère. »
\item \textbf{Règle d'or :} dans la subordonnée, le verbe conjugué va \textbf{à la toute dernière place}. C'est la différence majeure avec le français.
\item \textbf{Subordonnée en tête :} si la subordonnée commence la phrase, le verbe de la principale vient juste après la virgule (il reste en 2\textsuperscript{e} position) : \all{\textbf{Wenn} Kevin an seine Familie \textbf{denkt}, \textbf{ruft} er sie \textbf{an}.}
\item \textbf{Ne pas confondre :} \all{wenn} (quand, à chaque fois, si) $\neq$ \all{als} (quand, une seule fois au passé) ; \all{wenn} $\neq$ \all{ob} (si interrogatif indirect : \all{Ich weiß nicht, ob er kommt.}).
\end{enumerate}
\begin{center}
\begin{tabularx}{\linewidth}{|X|X|X|}
\hline
\rowcolor{popblueL} Français & Deutsch & Remarque \\
\hline
Il dit qu'il est content. & \all{Er sagt, dass er zufrieden ist.} & verbe à la fin \\
\hline
Quand il a le mal du pays, il téléphone. & \all{Wenn er Heimweh hat, telefoniert er.} & verbe après la virgule \\
\hline
S'il réussit, il restera. & \all{Wenn er Erfolg hat, bleibt er.} & condition = \all{wenn} \\
\hline
\end{tabularx}
\end{center}
\end{methode}

\begin{exoresolu}[Grammaire — transformation et construction de phrases]
\textbf{Énoncé A.} Relie les deux phrases avec \all{dass} ou \all{wenn}.
\begin{enumerate}
\item \all{Kevin hofft. Er findet eine Arbeit in Berlin.}
\item \all{Er hat Heimweh. Er hört kamerunische Musik.}
\item \all{Seine Mutter sagt. Sie ist stolz auf ihn.}
\end{enumerate}
\textbf{Énoncé B.} Commence la phrase par la subordonnée : \all{Kevin ruft seine Familie an, wenn er Zeit hat.}
\tcblower
\textbf{Raisonnement A.} J'identifie la relation logique : opinion/espoir $\rightarrow$ \all{dass} ; circonstance (quand, chaque fois que) $\rightarrow$ \all{wenn}. Puis je mets le verbe conjugué à la fin.
\begin{enumerate}
\item Espoir $\rightarrow$ \all{dass} : \all{Kevin hofft, dass er eine Arbeit in Berlin \textbf{findet}.} « Kevin espère qu'il trouvera un travail à Berlin. »
\item Circonstance $\rightarrow$ \all{wenn} : \all{Wenn Kevin Heimweh \textbf{hat}, hört er kamerunische Musik.} « Quand il a le mal du pays, il écoute de la musique camerounaise. »
\item Parole rapportée $\rightarrow$ \all{dass} : \all{Seine Mutter sagt, dass sie stolz auf ihn \textbf{ist}.} « Sa mère dit qu'elle est fière de lui. »
\end{enumerate}
\textbf{Raisonnement B.} La subordonnée passe en tête : elle occupe la 1\textsuperscript{re} position, donc le verbe de la principale (\all{ruft}) vient immédiatement après la virgule, \emph{avant} le sujet.

\textbf{Solution :} \all{Wenn Kevin Zeit hat, ruft er seine Familie an.} « Quand Kevin a le temps, il appelle sa famille. » — Attention à \all{anrufen}, verbe à particule séparable : la particule \all{an} retourne à la fin de la principale.

\textbf{Conclusion.} La virgule est obligatoire entre les deux propositions, et le verbe conjugué de la subordonnée se place toujours en dernière position.
\end{exoresolu}

\courssub{Méthode 5 : Rédiger un court paragraphe sur l'émigration}

\begin{methode}[Produire un texte guidé : « mon histoire d'émigration »]
\begin{enumerate}
\item \textbf{Prépare un plan en 3 parties :} (a) la situation de départ (qui, d'où, pourquoi) ; (b) le voyage et l'arrivée (au Perfekt) ; (c) la vie actuelle et les sentiments (\all{Heimweh, Hoffnung}).
\item \textbf{Utilise au moins 5 mots du lexique du thème :} \all{auswandern, das Ausland, der Pass, die Grenze, das Heimweh, die Hoffnung}.
\item \textbf{Alterne les temps :} Perfekt pour le passé (\all{ich bin ausgewandert, ich habe gelernt}), Präsens pour le présent (\all{ich wohne, ich arbeite}).
\item \textbf{Insère au moins une subordonnée en \all{dass} et une en \all{wenn}} pour montrer ta maîtrise de la grammaire.
\item \textbf{Relis-toi :} majuscule à tous les noms, verbe en 2\textsuperscript{e} position, participe II à la fin, Umlaute (\all{ä, ö, ü}) correctement écrits.
\end{enumerate}
\textbf{Amorce possible :} \all{Ich heiße... und ich komme aus... Vor zwei Jahren bin ich nach... ausgewandert, weil...}
\end{methode}

\begin{exoresolu}[Expression écrite — rédaction guidée (60-80 mots)]
\textbf{Énoncé.} Raconte l'histoire (réelle ou imaginaire) d'un jeune Camerounais qui émigre en Allemagne. Utilise : le Perfekt, une subordonnée en \all{dass}, une subordonnée en \all{wenn}, et au moins 5 mots du lexique de l'émigration.
\tcblower
\textbf{Raisonnement.} Je suis le plan en 3 parties : départ $\rightarrow$ voyage (Perfekt) $\rightarrow$ vie actuelle. Je vérifie chaque phrase : place du verbe, auxiliaire du Perfekt, verbe final dans les subordonnées.

\textbf{Production modèle :}

\all{Mein Freund Arno kommt aus Yaoundé. Er hat große Hoffnungen für seine Zukunft, aber in Kamerun hat er keine Arbeit gefunden. Deshalb ist er vor einem Jahr nach Deutschland ausgewandert. Vor der Abreise hat er einen Pass beantragt und ein Visum bekommen. An der Grenze hat er alle Dokumente gezeigt. Jetzt wohnt er in München und lernt jeden Tag Deutsch. Er sagt, dass das Leben im Ausland nicht einfach ist. Wenn er Heimweh hat, kocht er Eru und ruft seine Familie an. Aber er hofft, dass er bald ein Stipendium bekommt.}

« Mon ami Arno vient de Yaoundé. Il a de grands espoirs pour son avenir, mais au Cameroun il n'a pas trouvé de travail. C'est pourquoi il a émigré en Allemagne il y a un an. Avant le départ, il a demandé un passeport et obtenu un visa. À la frontière, il a montré tous ses documents. Maintenant il habite à Munich et apprend l'allemand tous les jours. Il dit que la vie à l'étranger n'est pas facile. Quand il a le mal du pays, il cuisine de l'eru et appelle sa famille. Mais il espère qu'il obtiendra bientôt une bourse. »

\textbf{Vérification finale :} 6 mots du thème (\all{Hoffnungen, ausgewandert, Pass, Grenze, Ausland, Heimweh}) ; Perfekt avec \all{sein} pour \all{ausgewandert} (déplacement) ; \all{dass...ist} et \all{wenn...hat} avec verbe final ; majuscules aux noms respectées.
\end{exoresolu}

\begin{attention}[Les 5 erreurs qui coûtent des points au Bac]
\begin{enumerate}
\item \textbf{Oublier la place du verbe dans la subordonnée.} Faux : \all{..., dass er findet eine Arbeit.} Juste : \all{..., dass er eine Arbeit \textbf{findet}.} Le verbe conjugué va \textbf{à la fin} après \all{dass}, \all{wenn}, \all{weil}.
\item \textbf{Confondre les auxiliaires du Perfekt.} Faux : \all{Er \textbf{hat} nach Deutschland \textbf{geflogen}.} Juste : \all{Er \textbf{ist} nach Deutschland geflogen.} Tout verbe de déplacement exige \all{sein}.
\item \textbf{Oublier la majuscule des noms.} Faux : \all{das heimweh, der pass, die grenze.} Juste : \all{das \textbf{H}eimweh, der \textbf{P}ass, die \textbf{G}renze.} En allemand, \emph{tous} les noms prennent la majuscule, partout dans la phrase.
\item \textbf{Calquer le français « émigrer en Allemagne ».} Faux : \all{in Deutschland auswandern} pour dire « vers ». Juste : \all{\textbf{nach} Deutschland auswandern}. La direction vers un pays sans article s'exprime avec \all{nach}.
\item \textbf{Confondre \all{wenn} et \all{als}, et les faux amis.} « Quand j'étais enfant » (une seule période passée) = \all{\textbf{Als} ich ein Kind war}, pas \all{wenn}. Attention aussi : \all{das Ausland} = l'étranger (pas « le pays ») ; \all{bekommen} = recevoir (pas « devenir » : \all{werden}).
\end{enumerate}
\end{attention}

\newpage

\coursec{Exercices}

\courssub{Je vérifie mes connaissances}

\begin{exercice}[QCM : le parfait]
Entourez la bonne réponse.
\begin{enumerate}
\item Kevin \textbf{...} nach Deutschland geflogen. \quad a) hat \quad b) ist \quad c) wird
\item Er \textbf{...} einen Pass beantragt. \quad a) ist \quad b) hat \quad c) war
\item Seine Familie \textbf{...} in Yaoundé geblieben. \quad a) hat \quad b) ist \quad c) hatte
\item Sie \textbf{...} viel über die Emigration gelernt. \quad a) ist \quad b) haben \quad c) sein
\end{enumerate}
\end{exercice}

\begin{exercice}[Vrai ou faux ?]
Répondez par \all{richtig} ou \all{falsch} et justifiez avec une phrase du cours.
\begin{enumerate}
\item \all{Das Heimweh} bedeutet : die Freude am neuen Land.
\item Au parfait, les verbes de mouvement se conjuguent avec \all{sein}.
\item Dans une subordonnée en \all{dass}, le verbe conjugué est à la deuxième place.
\item \all{Wenn} peut exprimer une condition ou une répétition dans le temps.
\end{enumerate}
\end{exercice}

\begin{exercice}[Vocabulaire : articles et pluriels]
Complétez le tableau avec l'article et le pluriel de chaque nom.
\begin{center}
\begin{tabularx}{\linewidth}{|X|X|X|}
\hline
\rowcolor{popblueL} Nom & Article & Pluriel \\
\hline
Pass & \ldots & \ldots \\
\hline
Grenze & \ldots & \ldots \\
\hline
Hoffnung & \ldots & \ldots \\
\hline
Ausland & \ldots & \ldots \\
\hline
Heimweh & \ldots & \ldots \\
\hline
\end{tabularx}
\end{center}
\end{exercice}

\begin{exercice}[Conjugaison à compléter]
Mettez les verbes entre parenthèses au parfait.
\begin{enumerate}
\item Kevin \ldots\ einen Koffer \ldots\ (packen).
\item Er \ldots\ nach Frankfurt \ldots\ (fliegen).
\item Seine Mutter \ldots\ am Flughafen \ldots\ (weinen).
\item Er \ldots\ eine Wohnung \ldots\ (finden).
\item Er \ldots\ oft an Kamerun \ldots\ (denken).
\end{enumerate}
\end{exercice}

\courssub{Je m'entraîne}

\begin{exercice}[Texte à trous]
Complétez le texte avec les mots suivants : \all{das Ausland, die Grenze, der Pass, das Heimweh, die Hoffnung, ausgewandert}.
\begin{textebox}[Ein neues Leben]
Viele junge Kameruner träumen vom Leben im \ldots\ldots\ldots . Bevor sie reisen, brauchen sie einen \ldots\ldots\ldots . An \ldots\ldots\ldots zeigen sie ihre Dokumente. Viele sind schon nach Europa \ldots\ldots\ldots . Sie leben mit \ldots\ldots\ldots auf eine bessere Zukunft, aber oft haben sie auch \ldots\ldots\ldots und vermissen ihre Familie.
\end{textebox}
\end{exercice}

\begin{exercice}[Appariements]
Reliez chaque début de phrase à la bonne fin.
\begin{enumerate}
\item Kevin sagt, \quad a) wenn er an seine Familie denkt.
\item Er hat Heimweh, \quad b) dass er oft an Kamerun denkt.
\item Er ruft seine Mutter an, \quad c) wenn er Zeit hat.
\item Er bleibt in Deutschland, \quad d) dass er dort Arbeit findet.
\item Er hofft, \quad e) wenn er genug Geld gespart hat.
\end{enumerate}
\end{exercice}

\begin{exercice}[Transformation : dass et wenn]
Transformez les phrases simples en subordonnées avec \all{dass} ou \all{wenn}, comme dans le modèle.\par
Modèle : \all{Kevin denkt oft an Kamerun. Er sagt:} $\rightarrow$ \all{Er sagt, dass er oft an Kamerun denkt.}
\begin{enumerate}
\item Er vermisst seine Familie. Er sagt:
\item Er kommt bald zurück. Er hofft:
\item Er hat viel gelernt. Er meint:
\item Er besucht Kamerun. Er freut sich, \ldots
\item Seine Mutter ist krank. Er kommt sofort, \ldots
\end{enumerate}
\end{exercice}

\begin{exercice}[wenn, als ou wann ?]
Complétez avec la bonne conjonction (\all{wenn}, \all{als} ou \all{wann}), puis traduisez les phrases 1 à 4 en français.
\begin{enumerate}
\item \ldots\ Kevin nach Deutschland gekommen ist, hat er zuerst Deutsch gelernt.
\item \ldots\ er Heimweh hat, ruft er seine Mutter an.
\item \ldots\ bist du ausgewandert? — Im Jahr 2022.
\item \ldots\ ich ein Kind war, lebten wir in Douala.
\item \ldots\ er genug Geld hat, kauft er ein Ticket.
\item Ich weiß nicht, \ldots\ der Zug abfährt.
\item \ldots\ es regnete, blieben wir zu Hause.
\item \ldots\ du Zeit hast, besuche mich!
\end{enumerate}
\end{exercice}

\courssub{Je me prépare au Bac}

\begin{exercice}[Compréhension écrite — type Bac]
Lisez le texte suivant, puis répondez aux questions \textbf{en allemand} et par des phrases complètes.
\begin{textebox}[Lena geht nach Kanada]
Lena Müller ist 24 Jahre alt und kommt aus München. Vor einem Jahr ist sie nach Kanada ausgewandert. „Ich habe in Deutschland Informatik studiert, aber ich habe keine Arbeit gefunden“, erzählt sie. „Dann habe ich ein Angebot aus Toronto bekommen.“ Zuerst war das Leben dort schwer: Der Winter war sehr kalt, und sie hatte oft Heimweh. Sie hat ihre Familie und die bayerischen Berge vermisst. Aber Lena hat schnell neue Freunde gefunden und viel Englisch gelernt. „Wenn ich an Deutschland denke, rufe ich meine Eltern an“, sagt sie. „Ich glaube, dass ich hier bleibe, aber ich hoffe, dass meine Familie mich bald besucht.“ Heute arbeitet Lena in einer großen Firma und ist zufrieden. Sie sagt: „Das Ausland hat mir eine Chance gegeben.“
\end{textebox}
\textbf{Fragen :}
\begin{enumerate}
\item Woher kommt Lena?
\item Warum ist sie nach Kanada ausgewandert?
\item Was war am Anfang schwer für sie? (2 Gründe)
\item Was macht sie, wenn sie an Deutschland denkt?
\item Richtig oder falsch? Begründen Sie mit dem Text : \all{Lena will sofort nach Deutschland zurückkehren.}
\item Trouvez dans le texte : un verbe au parfait avec \all{sein} et une subordonnée en \all{dass}.
\end{enumerate}
\end{exercice}

\begin{exercice}[Thème et version — type Bac]
\textbf{A. Version.} Traduisez en français :\par
\all{Kevin ist vor zwei Jahren ausgewandert. Er sagt, dass er oft Heimweh hat. Wenn er an Kamerun denkt, wird er traurig, aber er verliert die Hoffnung nicht.}\par\medskip
\textbf{B. Thème.} Traduisez en allemand :
\begin{enumerate}
\item Il dit qu'il a le mal du pays.
\item Quand il est arrivé à la frontière, il a montré son passeport.
\item Beaucoup de jeunes sont partis à l'étranger.
\item Elle espère que sa famille viendra bientôt.
\item Quand j'ai le mal du pays, j'écoute de la musique camerounaise.
\end{enumerate}
\end{exercice}

\begin{exercice}[Production écrite — type Bac]
\textbf{Sujet :} \all{Erzählen Sie die Geschichte einer Person, die ins Ausland gegangen ist.}\par
Rédigez un texte de \textbf{10 lignes environ (80 à 100 mots)}. Consignes obligatoires :
\begin{itemize}
\item utilisez au moins \textbf{3 verbes au parfait} (attention au choix de l'auxiliaire) ;
\item intégrez au moins \textbf{une subordonnée en \all{dass}} et \textbf{une en \all{wenn}} ;
\item employez au moins \textbf{4 mots du vocabulaire de la leçon} (\all{das Ausland, der Pass, die Grenze, das Heimweh, die Hoffnung, auswandern}).
\end{itemize}
\end{exercice}

\newpage

\coursec{Corrigés des exercices}

\begin{corrige}[Exercice 1 — QCM : le parfait]
\textbf{1. b) ist} — \all{Kevin ist nach Deutschland geflogen.} Le verbe \all{fliegen} exprime un déplacement : auxiliaire \all{sein}.\par
\textbf{2. b) hat} — \all{Er hat einen Pass beantragt.} \all{beantragen} est un verbe transitif (avec accusatif : \all{einen Pass}) : auxiliaire \all{haben}. Attention : \all{beantragen} est un verbe à préfixe inséparable \all{be-}, son participe passé ne prend donc pas de \all{ge-} : \all{beantragt} (et non \all{gebeantragt}).\par
\textbf{3. b) ist} — \all{Seine Familie ist in Yaoundé geblieben.} \all{bleiben} se conjugue toujours avec \all{sein} au parfait (exception à mémoriser, comme \all{sein} lui-même : \all{ist gewesen}).\par
\textbf{4. b) haben} — \all{Sie haben viel über die Emigration gelernt.} \all{lernen} prend \all{haben} ; attention à l'accord du pluriel : \all{sie haben} (et non \all{hat}).
\end{corrige}

\begin{corrige}[Exercice 2 — Vrai ou faux ?]
\textbf{1. falsch} — \all{Das Heimweh} bedeutet nicht die Freude am neuen Land, sondern die Sehnsucht nach der Heimat (le mal du pays, la nostalgie de la patrie).\par
\textbf{2. richtig} — Au parfait, les verbes de mouvement (\all{gehen, fliegen, kommen, reisen, auswandern}) se conjuguent avec \all{sein} : \all{Kevin ist nach Deutschland geflogen.}\par
\textbf{3. falsch} — Dans une subordonnée en \all{dass}, le verbe conjugué est à la \textbf{dernière place} : \all{Er sagt, dass er oft an Kamerun denkt.}\par
\textbf{4. richtig} — \all{Wenn} exprime une condition (\all{Wenn er Geld hat, kauft er ein Ticket.}) ou une répétition dans le temps (\all{Wenn er Heimweh hat, ruft er seine Mutter an.}).
\end{corrige}

\begin{corrige}[Exercice 3 — Vocabulaire : articles et pluriels]
\begin{center}
\begin{tabularx}{\linewidth}{|X|X|X|}
\hline
\rowcolor{popblueL} Nom & Article & Pluriel \\
\hline
Pass & der & die Pässe \\
\hline
Grenze & die & die Grenzen \\
\hline
Hoffnung & die & die Hoffnungen \\
\hline
Ausland & das & (kein Plural) \\
\hline
Heimweh & das & (kein Plural) \\
\hline
\end{tabularx}
\end{center}
\textbf{Points de vigilance :} \all{der Pass} prend un tréma au pluriel : \all{die Pässe}. \all{Das Ausland} et \all{das Heimweh} sont des noms abstraits sans pluriel usuel. On retiendra aussi la famille de mots : \all{auswandern} (émigrer, verbe) $\rightarrow$ \all{die Auswanderung} (l'émigration) $\rightarrow$ \all{der Auswanderer} (l'émigrant).
\end{corrige}

\begin{corrige}[Exercice 4 — Conjugaison à compléter]
\begin{enumerate}
\item \all{Kevin hat einen Koffer gepackt.} (verbe faible, auxiliaire \all{haben})
\item \all{Er ist nach Frankfurt geflogen.} (verbe de mouvement, auxiliaire \all{sein} ; participe irrégulier : \all{fliegen — flog — ist geflogen})
\item \all{Seine Mutter hat am Flughafen geweint.} (verbe faible, \all{haben})
\item \all{Er hat eine Wohnung gefunden.} (verbe fort : \all{finden — fand — hat gefunden})
\item \all{Er hat oft an Kamerun gedacht.} (verbe irrégulier : \all{denken — dachte — hat gedacht})
\end{enumerate}
\textbf{Rappel :} le participe passé se place à la fin de la phrase ; l'auxiliaire conjugué à la deuxième place.
\end{corrige}

\begin{corrige}[Exercice 5 — Texte à trous]
Viele junge Kameruner träumen vom Leben im \textbf{Ausland}. Bevor sie reisen, brauchen sie einen \textbf{Pass}. An \textbf{der Grenze} zeigen sie ihre Dokumente. Viele sind schon nach Europa \textbf{ausgewandert}. Sie leben mit \textbf{der Hoffnung} auf eine bessere Zukunft, aber oft haben sie auch \textbf{Heimweh} und vermissen ihre Familie.\par
\textbf{Traduction :} Beaucoup de jeunes Camerounais rêvent de la vie à l'étranger. Avant de voyager, ils ont besoin d'un passeport. À la frontière, ils montrent leurs documents. Beaucoup sont déjà partis (ont émigré) en Europe. Ils vivent avec l'espoir d'un avenir meilleur, mais souvent ils ont aussi le mal du pays et leur famille leur manque.\par
\textbf{Points de vigilance :} \all{im Ausland} (datif contracté de \all{in dem Ausland}) ; \all{an der Grenze} (datif, lieu) ; \all{sind ausgewandert} (verbe de mouvement avec \all{sein}, particule \all{aus-} recollée dans le participe : \all{ausgewandert}, placé à la fin).
\end{corrige}

\begin{corrige}[Exercice 6 — Appariements]
\textbf{1 — b :} \all{Kevin sagt, dass er oft an Kamerun denkt.} (Kevin dit qu'il pense souvent au Cameroun.)\par
\textbf{2 — a :} \all{Er hat Heimweh, wenn er an seine Familie denkt.} (Il a le mal du pays quand il pense à sa famille.)\par
\textbf{3 — c :} \all{Er ruft seine Mutter an, wenn er Zeit hat.} (Il appelle sa mère quand il a le temps.)\par
\textbf{4 — e :} \all{Er bleibt in Deutschland, wenn er genug Geld gespart hat.} (Il reste en Allemagne quand il a assez économisé.)\par
\textbf{5 — d :} \all{Er hofft, dass er dort Arbeit findet.} (Il espère qu'il trouvera du travail là-bas.)\par
\textbf{Justification :} après \all{sagen} et \all{hoffen} (verbes de parole et d'opinion), on emploie \all{dass} ; \all{wenn} introduit une circonstance temporelle ou une condition. Dans les deux cas, le verbe conjugué va à la fin de la subordonnée.
\end{corrige}

\begin{corrige}[Exercice 7 — Transformation : dass et wenn]
\begin{enumerate}
\item \all{Er sagt, dass er seine Familie vermisst.} (Il dit que sa famille lui manque.)
\item \all{Er hofft, dass er bald zurückkommt.} (Il espère qu'il reviendra bientôt.)
\item \all{Er meint, dass er viel gelernt hat.} (Il pense qu'il a beaucoup appris.)
\item \all{Er freut sich, wenn er Kamerun besucht.} (Il se réjouit quand il visite le Cameroun.)
\item \all{Er kommt sofort, wenn seine Mutter krank ist.} (Il vient tout de suite quand sa mère est malade.)
\end{enumerate}
\textbf{Points de vigilance :} dans la subordonnée, le verbe conjugué passe à la dernière place (\all{vermisst}, \all{zurückkommt}, \all{hat}, \all{besucht}, \all{ist}). Pour le verbe à particule \all{zurückkommen}, la particule se recolle au verbe en fin de subordonnée : \all{zurückkommt} (et non \all{kommt zurück}). Au parfait dans une subordonnée, l'auxiliaire conjugué se place après le participe : \all{dass er viel gelernt hat}.
\end{corrige}

\begin{corrige}[Exercice 8 — wenn, als ou wann ?]
\begin{enumerate}
\item \all{\textbf{Als} Kevin nach Deutschland gekommen ist, hat er zuerst Deutsch gelernt.} — Quand Kevin est arrivé en Allemagne, il a d'abord appris l'allemand. (passé, événement unique $\rightarrow$ \all{als})
\item \all{\textbf{Wenn} er Heimweh hat, ruft er seine Mutter an.} — Quand il a le mal du pays, il appelle sa mère. (présent, répétition $\rightarrow$ \all{wenn})
\item \all{\textbf{Wann} bist du ausgewandert? — Im Jahr 2022.} — Quand es-tu parti à l'étranger ? — En 2022. (question directe $\rightarrow$ \all{wann})
\item \all{\textbf{Als} ich ein Kind war, lebten wir in Douala.} — Quand j'étais enfant, nous vivions à Douala. (passé, période unique $\rightarrow$ \all{als})
\item \all{\textbf{Wenn} er genug Geld hat, kauft er ein Ticket.} — Quand il a assez d'argent, il achète un billet. (condition au présent $\rightarrow$ \all{wenn})
\item \all{Ich weiß nicht, \textbf{wann} der Zug abfährt.} — Je ne sais pas quand le train part. (question indirecte $\rightarrow$ \all{wann})
\item \all{\textbf{Als} es regnete, blieben wir zu Hause.} — Quand il a plu, nous sommes restés à la maison. (passé, événement unique $\rightarrow$ \all{als})
\item \all{\textbf{Wenn} du Zeit hast, besuche mich!} — Quand tu as le temps, viens me voir ! (présent/futur, condition $\rightarrow$ \all{wenn})
\end{enumerate}
\textbf{Règle récapitulative :} \all{wann} = question (directe ou indirecte) ; \all{als} = passé, une seule fois ; \all{wenn} = présent, futur, répétition ou condition.
\end{corrige}

\begin{corrige}[Exercice 9 — Compréhension écrite — type Bac]
\textbf{Barème indicatif : 12 points} (2 points par question ; 1 point pour le contenu, 1 point pour la correction de la langue).\par\medskip
\textbf{1.} \all{Lena kommt aus München.} (Lena vient de Munich.)\par
\textbf{2.} \all{Sie ist nach Kanada ausgewandert, weil sie in Deutschland keine Arbeit gefunden hat und ein Angebot aus Toronto bekommen hat.} (Elle a émigré au Canada parce qu'elle n'a pas trouvé de travail en Allemagne et qu'elle a reçu une offre de Toronto.)\par
\textbf{3.} \all{Am Anfang war der Winter sehr kalt, und sie hatte oft Heimweh.} (Au début, l'hiver était très froid et elle avait souvent le mal du pays.)\par
\textbf{4.} \all{Wenn sie an Deutschland denkt, ruft sie ihre Eltern an.} (Quand elle pense à l'Allemagne, elle appelle ses parents.)\par
\textbf{5.} \all{Falsch. Sie sagt: „Ich glaube, dass ich hier bleibe.“} (Faux. Elle dit : « Je crois que je reste ici. »)\par
\textbf{6.} Verbe au parfait avec \all{sein} : \all{ist sie nach Kanada ausgewandert} (ou \all{ist ... gekommen}). Subordonnée en \all{dass} : \all{Ich glaube, dass ich hier bleibe.} (ou \all{ich hoffe, dass meine Familie mich bald besucht}).\par\medskip
\textbf{Points de vigilance :} répondre par des phrases complètes avec le verbe en deuxième position ; reprendre le vocabulaire du texte sans le copier mot pour mot sur toute la longueur ; attention au parfait avec \all{sein} pour \all{auswandern}.
\end{corrige}

\begin{corrige}[Exercice 10 — Thème et version — type Bac]
\textbf{Barème indicatif : 10 points} (Version : 4 points ; Thème : 6 points, environ 1,2 point par phrase).\par\medskip
\textbf{A. Version.} « Kevin a émigré il y a deux ans. Il dit qu'il a souvent le mal du pays. Quand il pense au Cameroun, il devient triste, mais il ne perd pas espoir. »\par
\textbf{Points de vigilance :} \all{vor zwei Jahren} = il y a deux ans ; \all{die Hoffnung nicht verlieren} = ne pas perdre espoir ; \all{wird er traurig} = il devient triste (inversion après la subordonnée en tête).\par\medskip
\textbf{B. Thème.}
\begin{enumerate}
\item \all{Er sagt, dass er Heimweh hat.}
\item \all{Als er an der Grenze angekommen ist, hat er seinen Pass gezeigt.} (passé, événement unique $\rightarrow$ \all{als} ; \all{ankommen} avec \all{sein})
\item \all{Viele junge Leute sind ins Ausland gegangen.} (mouvement $\rightarrow$ \all{sein})
\item \all{Sie hofft, dass ihre Familie bald kommt.} (ou : \all{bald kommen wird})
\item \all{Wenn ich Heimweh habe, höre ich kamerunische Musik.} (subordonnée en tête $\rightarrow$ inversion dans la principale)
\end{enumerate}
\textbf{Points de vigilance :} ne pas traduire « quand » au passé par \all{wenn} (phrase 2) ; majuscule aux noms (\all{Pass, Grenze, Ausland, Musik}) ; verbe conjugué en fin de subordonnée.
\end{corrige}

\begin{corrige}[Exercice 11 — Production écrite — type Bac]
\textbf{Barème indicatif : 10 points.} Respect des consignes et contenu : 4 points ; grammaire (parfait, subordonnées) : 3 points ; vocabulaire thématique : 2 points ; orthographe et majuscules : 1 point.\par\medskip
\textbf{Proposition de rédaction modèle :}\par\medskip
\emph{Aminas Weg nach Deutschland}\par
\all{Meine Cousine Amina ist vor drei Jahren ins Ausland gegangen. Sie hat zuerst einen Pass beantragt und dann ein Visum bekommen. An der Grenze hat sie ihre Dokumente gezeigt. Dann ist sie nach Berlin geflogen. Am Anfang war alles schwer: Sie hatte Heimweh und hat ihre Familie sehr vermisst. Sie sagt, dass sie oft an Kamerun denkt. Wenn sie traurig ist, ruft sie ihre Mutter an. Aber Amina hat die Hoffnung nicht verloren: Sie hat Deutsch gelernt und eine Arbeit gefunden. Heute ist sie zufrieden, und sie hofft, dass ihre Familie sie bald besucht.}\par\medskip
\textbf{Traduction :} Ma cousine Amina est partie à l'étranger il y a trois ans. Elle a d'abord demandé un passeport, puis obtenu un visa. À la frontière, elle a montré ses documents. Ensuite, elle a pris l'avion pour Berlin. Au début, tout était difficile : elle avait le mal du pays et sa famille lui manquait beaucoup. Elle dit qu'elle pense souvent au Cameroun. Quand elle est triste, elle appelle sa mère. Mais Amina n'a pas perdu espoir : elle a appris l'allemand et trouvé un travail. Aujourd'hui, elle est satisfaite et elle espère que sa famille lui rendra bientôt visite.\par\medskip
\textbf{Points de vigilance :}
\begin{itemize}
\item Auxiliaire \all{sein} pour les verbes de mouvement : \all{ist gegangen, ist geflogen} ; \all{haben} pour les autres : \all{hat beantragt, hat gelernt, hat gefunden}.
\item Verbe conjugué à la fin des subordonnées : \all{dass sie oft an Kamerun denkt} ; \all{wenn sie traurig ist}.
\item Inversion après une subordonnée en tête : \all{Wenn sie traurig ist, ruft sie ... an.}
\item Majuscule à tous les noms : \all{Pass, Grenze, Heimweh, Hoffnung, Familie, Arbeit}.
\end{itemize}
\end{corrige}

\newpage

\coursec{Fiche bilan}

\begin{popfigure}
\begin{tikzpicture}[mindmap, grow cyclic, every node/.style={concept, align=center, font=\small},
  root concept/.append style={concept color=popblue, font=\bfseries},
  level 1/.append style={level distance=3.4cm, sibling angle=60},
  level 1 concept/.append style={font=\small}]
\node[root concept, align=center]{Die Emigration\\ (ALA24)}
  child[concept color=poporange]{ node[align=center]{Texte\\ \all{Kevins neues}\\ \all{Leben in}\\ \all{Deutschland}} }
  child[concept color=popgreen]{ node[align=center]{Lexique\\ \all{auswandern}\\ \all{das Heimweh}\\ \all{die Grenze}} }
  child[concept color=poppurple]{ node[align=center]{Perfekt\\ \all{haben/sein}\\ + Partizip II} }
  child[concept color=poppink]{ node[align=center]{\all{dass}\\ verbe\\ à la fin} }
  child[concept color=popgold]{ node[align=center]{\all{wenn}\\ condition\\ + répétition} }
  child[concept color=popblue!60]{ node[align=center]{Méthode\\ lire -- comprendre\\ -- résumer} }
  child[concept color=popgreen!60]{ node[align=center]{Production\\ mon histoire\\ d'émigration} };
\end{tikzpicture}
\legende{Carte mentale de la leçon : l'émigration}
\end{popfigure}

\begin{bilan}
\textbf{1. Lexique clé (à connaître avec article et pluriel)}

\begin{tabularx}{\linewidth}{|X|X|}\hline
\rowcolor{popblueL} \textbf{Deutsch} & \textbf{Français} \\ \hline
\all{die Emigration} (nur Singular) & l'émigration \\ \hline
\all{auswandern} (wanderte aus, ist ausgewandert) & émigrer \\ \hline
\all{der Auswanderer, -} / \all{die Auswanderin, -nen} & l'émigrant / l'émigrante \\ \hline
\all{das Ausland} (nur Singular) & l'étranger \\ \hline
\all{das Heimweh} (nur Singular) & le mal du pays \\ \hline
\all{die Hoffnung, -en} & l'espoir \\ \hline
\all{der Pass, ⸚e} & le passeport \\ \hline
\all{die Grenze, -n} & la frontière \\ \hline
\all{das Visum, die Visa} & le visa \\ \hline
\all{die Heimat, -en} & la patrie, le pays natal \\ \hline
\all{verlassen} (verließ, hat verlassen) & quitter \\ \hline
\all{ankommen} (kam an, ist angekommen) & arriver \\ \hline
\all{sich integrieren} & s'intégrer \\ \hline
\all{vermissen} (hat vermisst) & manquer (regretter) \\ \hline
\end{tabularx}

\textbf{2. Grammaire à connaître par cœur}

\begin{tabularx}{\linewidth}{|X|X|}\hline
\rowcolor{popgreenL} \textbf{Règle} & \textbf{Exemple} \\ \hline
Perfekt = \all{haben} ou \all{sein} au Présent + Partizip II \textbf{en fin de phrase} & \all{Kevin \textbf{ist} nach Deutschland \textbf{geflogen}.} \\ \hline
\all{sein} avec les verbes de déplacement et de changement d'état (\all{gehen, kommen, fliegen, auswandern, sterben}) & \all{Er \textbf{ist} in Yaoundé \textbf{angekommen}.} \\ \hline
\all{haben} avec tous les autres verbes & \all{Sie \textbf{hat} ihre Familie \textbf{vermisst}.} \\ \hline
Partizip II : faibles \all{ge...t} (\all{gemacht}), forts \all{ge...en} (\all{gegangen}) & \all{gearbeitet / gefahren} \\ \hline
Verbes à particule : \all{ge} entre particule et radical & \all{\textbf{ausge}wandert, \textbf{ange}kommen} \\ \hline
Subordonnée en \all{dass} : verbe conjugué \textbf{à la fin} & \all{Ich hoffe, \textbf{dass} er bald zurück\textbf{kommt}.} \\ \hline
\all{wenn} = condition (si) et répétition (quand, chaque fois que) ; verbe à la fin & \all{\textbf{Wenn} er Heimweh \textbf{hat}, ruft er seine Mutter \textbf{an}.} \\ \hline
Subordonnée devant : le verbe de la principale vient juste après la virgule & \all{\textbf{Wenn} er Zeit \textbf{hat}, \textbf{besucht} er seine Freunde.} \\ \hline
\end{tabularx}

\textbf{3. Méthodes express}
\begin{itemize}
\item \textbf{Comprendre un texte} : lire le titre, lire une première fois pour l'idée générale, souligner les mots du thème, relire pour les détails, répondre par des phrases complètes.
\item \textbf{Résumer} : repérer qui ? où ? quand ? pourquoi ? ; reformuler en 3 à 5 phrases au Perfekt ou au Präsens.
\item \textbf{Produire un court texte} : introduction (qui, d'où), déroulement (voyage, arrivée), bilan (sentiments, espoirs) ; vérifier la place du verbe et les majuscules des noms.
\end{itemize}
\end{bilan}

\begin{aretenir}[Checklist avant le Bac]
$\square$ Je connais le lexique de l'émigration avec l'article et le pluriel de chaque nom.

$\square$ Je sais conjuguer un verbe au Perfekt avec le bon auxiliaire (\all{haben} ou \all{sein}).

$\square$ Je place toujours le Partizip II en fin de phrase au Perfekt.

$\square$ Je forme correctement le Partizip II des verbes faibles, forts et à particule séparable.

$\square$ Je mets le verbe conjugué à la fin dans une subordonnée en \all{dass}.

$\square$ Je distingue \all{wenn} (si, chaque fois que) et \all{als} (une seule fois au passé).

$\square$ Quand la subordonnée ouvre la phrase, je place le verbe de la principale juste après la virgule.

$\square$ Je mets une majuscule à tous les noms allemands et j'utilise les guillemets „ … “.

$\square$ Je sais répondre à une question de compréhension par une phrase complète.

$\square$ Je sais rédiger un court paragraphe (5 à 8 phrases) sur une histoire d'émigration.

$\square$ Je relis ma production : place du verbe, auxiliaire, orthographe des Umlaute (ä, ö, ü, ß).
\end{aretenir}

\findecours

\end{document}
